%% GEB LaTeX v2 — Part 06 / 7
%% 自动生成  split_v2.py
%% 独立编译（从 GEB_LaTeX_v2/ 目录）：
%%   xelatex -output-directory=split part06.tex

%% GEB 精排模板 ── XeLaTeX
\PassOptionsToPackage{unicode}{hyperref}
\PassOptionsToPackage{hyphens}{url}

%% ─────────────────────────────────────────────
%%  文档类：memoir（精排书籍首选）
%% ─────────────────────────────────────────────
\documentclass[
  12pt,
  twoside,
  openright,
  oldfontcommands,
]{memoir}

%% memoir 纸张尺寸（必须在字体等 usepackage 之前）
\setstocksize{260mm}{185mm}   % 大16开
\settrimmedsize{\stockheight}{\stockwidth}{*}
\settrims{0pt}{0pt}

%% ─────────────────────────────────────────────
%%  字体
%% ─────────────────────────────────────────────
\usepackage{fontspec}
\usepackage{xeCJK}

% ※（U+203B）Georgia 无此字符，强制用 CJK 字体渲染
\xeCJKDeclareCharClass{CJK}{"203B}

% geb.ttf：EPUB 自定义字体，含 Tai Viet 占位字形（U+AA9E–U+AAAF）
% 以及 U+31CF、U+39DF 等特殊字形
\newfontfamily\gebfont{geb.ttf}[Path=./fonts/]

% 中文正文：思源宋体（Noto Serif CJK SC）
\setCJKmainfont[
  BoldFont    = {Noto Serif CJK SC SemiBold},
  ItalicFont  = {Kaiti SC},          % 楷体作斜体
]{Noto Serif CJK SC}

% 中文无衬线（标题等）
\setCJKsansfont[
  BoldFont = {Heiti SC},
]{Heiti SC}

% 中文等宽
\setCJKmonofont{PingFang SC}

% 西文正文：Palatino（等高数字，基线与思源宋体对齐更佳）
\setmainfont[
  BoldFont       = {Palatino Bold},
  ItalicFont     = {Palatino Italic},
  BoldItalicFont = {Palatino Bold Italic},
]{Palatino}

% 西文无衬线
\setsansfont{Helvetica Neue}

% 数学字体
\usepackage{amsmath,amssymb}
\usepackage{unicode-math}

%% ─────────────────────────────────────────────
%%  版面：大16开（185×260mm）页边距（memoir 原生布局）
%% ─────────────────────────────────────────────
\setlrmarginsandblock{25mm}{20mm}{*}
\setulmarginsandblock{22mm}{25mm}{*}
\setheaderspaces{*}{7mm}{*}
\setheadfoot{\baselineskip}{12mm}
\checkandfixthelayout

%% ─────────────────────────────────────────────
%%  颜色
%% ─────────────────────────────────────────────
\usepackage{xcolor}
\definecolor{accent}{HTML}{1A3A5C}     % 深蓝（章号、页眉线）
\definecolor{accentlight}{HTML}{4A7DAA}% 中蓝（节标题）
\definecolor{rulegray}{HTML}{BBBBBB}   % 分隔线

%% ─────────────────────────────────────────────
%%  章节样式
%% ─────────────────────────────────────────────
\usepackage{titlesec}

% 章：大号宋体居中，上方装饰线，下方细实线
\titleformat{\chapter}[display]
  {\normalfont\centering}
  {\color{accent}\small\CJKfamily{}\sffamily 第\ \thechapter\ 章}
  {0pt}
  {\color{accent}\LARGE\bfseries\CJKfamily{}\rmfamily}
  [\vspace{4pt}\textcolor{rulegray}{\titlerule[0.4pt]}]

\titlespacing*{\chapter}{0pt}{30pt}{40pt}

% 节：宋体加粗左对齐，左侧蓝色竖线装饰
\titleformat{\section}
  {\normalfont\large\bfseries\color{accentlight}}
  {}
  {0pt}
  {\colorbox{white}{\raisebox{0pt}{\rule[-1pt]{3pt}{14pt}}\hspace{6pt}}}
  []

\titlespacing*{\section}{0pt}{18pt}{8pt}

% 小节
\titleformat{\subsection}
  {\normalfont\normalsize\bfseries}
  {}
  {0pt}
  {}

%% ─────────────────────────────────────────────
%%  页眉页脚（memoir 内建）
%% ─────────────────────────────────────────────
\makepagestyle{geb}
\makeevenhead{geb}{\small\thepage}{}{\small\itshape\leftmark}
\makeoddhead {geb}{\small\itshape\rightmark}{}{\small\thepage}
\makeheadrule {geb}{\textwidth}{0.4pt}
\pagestyle{geb}
\renewcommand{\chaptermark}[1]{\markboth{第\ \thechapter\ 章\enspace #1}{第\ \thechapter\ 章\enspace #1}}
\renewcommand{\sectionmark}[1]{\markright{\thesection\enspace #1}}

%% ─────────────────────────────────────────────
%%  中文标准名称
%% ─────────────────────────────────────────────
\renewcommand{\contentsname}{目录}
\renewcommand{\figurename}{图}
\renewcommand{\tablename}{表}
\renewcommand{\appendixname}{附录}
\renewcommand{\bibname}{参考文献}
\renewcommand{\indexname}{索引}

%% ─────────────────────────────────────────────
%%  段落间距与行距（memoir 内建命令）
%% ─────────────────────────────────────────────
\setSpacing{1.3}
\setlength{\parindent}{2em}
\setlength{\parskip}{3pt plus 1pt minus 1pt}
\raggedbottom

%% ─────────────────────────────────────────────
%%  图片与浮动体
%% ─────────────────────────────────────────────
\usepackage{graphicx}
\usepackage{float}
\usepackage{caption}
\captionsetup{
  font=small,
  labelfont={bf,color=accent},
  labelsep=quad,
  width=0.88\linewidth,
  justification=centering,
  singlelinecheck=false,
}

%% pandocbounded：memoir 兼容版，限制图片最大宽度和高度
\newcommand*\pandocbounded[1]{%
  \begingroup
  \setkeys{Gin}{width=\linewidth,height=0.8\textheight,keepaspectratio}%
  #1%
  \endgroup
}
\def\fps@figure{htbp}

%% ─────────────────────────────────────────────
%%  引用块样式
%% ─────────────────────────────────────────────
\usepackage{mdframed}
\mdfdefinestyle{gebquote}{
  leftline=true,
  rightline=false,
  topline=false,
  bottomline=false,
  linewidth=2pt,
  linecolor=accentlight,
  backgroundcolor=black!3,
  innerleftmargin=12pt,
  innerrightmargin=8pt,
  innertopmargin=6pt,
  innerbottommargin=6pt,
}
\renewenvironment{quote}
  {\begin{mdframed}[style=gebquote]\small}
  {\end{mdframed}}

%% ─────────────────────────────────────────────
%%  表格
%% ─────────────────────────────────────────────
\usepackage{longtable,booktabs,array,multirow,calc}
\usepackage{etoolbox}
\makeatletter
\patchcmd\longtable{\par}{\if@noskipsec\mbox{}\fi\par}{}{}
\makeatother
\IfFileExists{footnotehyper.sty}{\usepackage{footnotehyper}}{\usepackage{footnote}}
\makesavenoteenv{longtable}

% 美化 booktabs 规则
\setlength{\aboverulesep}{0pt}
\setlength{\belowrulesep}{0pt}
\setlength{\heavyrulewidth}{0.8pt}
\setlength{\lightrulewidth}{0.4pt}

%% ─────────────────────────────────────────────
%%  脚注
%% ─────────────────────────────────────────────
\usepackage[
  bottom,
  hang,
  flushmargin,
]{footmisc}
\setlength{\footnotemargin}{1.5em}
\renewcommand{\footnoterule}{%
  \vspace{4pt}\noindent\textcolor{rulegray}{\rule{0.3\linewidth}{0.4pt}}\vspace{4pt}%
}

%% ─────────────────────────────────────────────
%%  目录样式
%% ─────────────────────────────────────────────
\usepackage{titletoc}

%% memoir 的 chapternumberline/partnumberline 需映射到 \numberline
%% 才能被 titletoc 识别为「有编号」条目
\makeatletter
\renewcommand{\chapternumberline}[1]{\numberline{#1}}
\renewcommand{\partnumberline}[1]{\numberline{#1}}
\makeatother

\titlecontents{chapter}[0em]
  {\addvspace{6pt}\bfseries\color{accent}}
  {\makebox[3em][l]{\thecontentslabel}}
  {}
  {\hfill\contentspage}

\titlecontents{section}[3em]
  {\addvspace{2pt}\small}
  {}
  {}
  {\titlerule*[0.5em]{.}\contentspage}

%% ─────────────────────────────────────────────
%%  超链接
%% ─────────────────────────────────────────────
\usepackage{bookmark}
\IfFileExists{xurl.sty}{\usepackage{xurl}}{}
\urlstyle{same}
\usepackage{hyperref}
\hypersetup{
  pdftitle={哥德尔、艾舍尔、巴赫——集异璧之大成},
  pdfauthor={〔美〕侯世达},
  pdflang={zh-CN},
  colorlinks=true,
  linkcolor={accent},
  urlcolor={accentlight},
  citecolor={accent},
  pdfcreator={XeLaTeX + GEB Template},
}

%% ─────────────────────────────────────────────
%%  杂项
%% ─────────────────────────────────────────────
\usepackage{microtype}
\setlength{\emergencystretch}{3em}
\providecommand{\tightlist}{%
  \setlength{\itemsep}{0pt}\setlength{\parskip}{0pt}}

% 对话角色加粗
\newcommand{\speaker}[1]{\textbf{#1}}

%% ─────────────────────────────────────────────
%%  GEB v2 字体命令
%% ─────────────────────────────────────────────
\newCJKfontfamily\kaitifamily[BoldFont={Kaiti SC Bold}]{Kaiti SC}
\newcommand{\kaititext}[1]{{\kaitifamily #1}}
\newcommand{\kaitismall}[1]{{\kaitifamily\small #1}}
\newcommand{\kaitibold}[1]{{\kaitifamily\bfseries #1}}
\newcommand{\heititext}[1]{{\sffamily\bfseries #1}}
\newcommand{\heitismall}[1]{{\sffamily\bfseries\small #1}}
\newcommand{\dialogname}[1]{{\sffamily\bfseries #1}}
\newcommand{\westerntext}[1]{{\rmfamily #1}}
\newcommand{\westernital}[1]{{\rmfamily\itshape #1}}

%% ─────────────────────────────────────────────
%%  GEB v2 排版环境（memoir 内建 hangparas，不需要 hanging 宏包）
%% ─────────────────────────────────────────────

% 对话悬挂缩进行
\newenvironment{hangdialog}{%
  \begin{hangparas}{2em}{1}%
}{%
  \par\end{hangparas}%
}

% 楷体舞台指示（居中小楷）
\newenvironment{stagedir}{%
  \begin{center}\kaitifamily\small%
}{%
  \end{center}%
}

% 对话引导语（楷体斜体缩进）
\newenvironment{dialogguide}{%
  \begin{adjustwidth}{1em}{0em}\kaitifamily\itshape\small%
}{%
  \end{adjustwidth}%
}

% 仿宋/楷体引用块（带左竖线）
\newenvironment{fsquote}{%
  \begin{mdframed}[style=gebquote]\kaitifamily\small%
}{%
  \end{mdframed}%
}

% 楷体加粗引用框
\newenvironment{ktquote}{%
  \begin{mdframed}[style=gebquote]\kaitifamily\small%
}{%
  \end{mdframed}%
}

% 禅宗公案块
\newenvironment{zenkoan}{%
  \begin{mdframed}[style=gebquote]\kaitifamily\small%
}{%
  \end{mdframed}%
}

% 禅宗公案诗
\newenvironment{zenkoanpoem}{%
  \begin{center}\kaitifamily\small%
}{%
  \end{center}%
}

% 图片副说明（小楷居中，灰色）
\newenvironment{imgsub}{%
  \begin{center}\kaitifamily\small\color{rulegray}%
}{%
  \end{center}%
}

% 对话诗歌
\newenvironment{dialogpoem}{%
  \begin{adjustwidth}{3em}{0em}\kaitifamily%
}{%
  \end{adjustwidth}%
}

% 章节前置数字大标题（如"第一章"）
\newenvironment{chappreheader}{%
  \begin{center}\large\sffamily\bfseries\color{accent}%
}{%
  \end{center}\vspace{-10pt}%
}

% 部分前置大标题
\newenvironment{partpreheader}{%
  \begin{center}\Large\sffamily\bfseries\color{accent}%
}{%
  \end{center}\vspace{-10pt}%
}

%% ─────────────────────────────────────────────
%%  扉页（自定义，替换 maketitle）
%% ─────────────────────────────────────────────
\newcommand{\gebmaketitle}{%
  \begin{titlingpage}
    \vspace*{6\baselineskip}
    \begin{center}
      {\fontsize{28}{34}\selectfont\bfseries\color{accent}
        哥德尔、艾舍尔、巴赫}\\[12pt]
      {\Large\color{rulegray}——集异璧之大成}\\[40pt]
      \textcolor{rulegray}{\rule{0.4\linewidth}{0.4pt}}\\[20pt]
      {\large \textit{Gödel, Escher, Bach: An Eternal Golden Braid}}\\[6pt]
      {\normalsize \color{accentlight} Douglas R.\ Hofstadter}\\[40pt]
      {\normalsize 〔美〕侯世达　著\\[4pt]
        本书翻译组　译}\\[60pt]
      {\small\color{rulegray} 商务印书馆}
    \end{center}
    \vfill
  \end{titlingpage}
}

%% pandoc 在 body 中会注入 \babel@toc，提供空定义避免报错
\makeatletter
\providecommand{\babel@toc}[2]{}
\makeatother

%% memoir + XeLaTeX 下需手动告知 xdvipdfmx 纸张尺寸
\AtBeginDocument{%
  \special{papersize=\the\paperwidth,\the\paperheight}%
}

%% ═══════════════════════════════════════════
%%  正文开始
%% ═══════════════════════════════════════════

\begin{document}
\mainmatter
\setcounter{chapter}{40}
\markboth{}{}

\chapter{的确该赞美螃蟹}\label{Dialog17_split_001.xhtml_calibre_pb_0}

明媚的春光里，乌龟和阿基里斯一起在林中散步。他们打算登上一座小山。据说，山顶上有间极好的茶馆，里面有各种精美的糕点。

\dialogname{阿基里斯：}\kaitismall{（一边爬山一边唱）}我心尊主为大，我灵以上帝我的救主为乐。

\dialogname{乌龟：}升D？我对音乐可是个内行，阿基，我敢说你唱的这支曲子是D调的，而不是升D。

\dialogname{阿基里斯：}的确，是D调的。我刚才的唱词是``上帝''，不是``升D''。我唱的这支曲子是巴赫作的《D调的赞美歌》，你难道从没听过？

\dialogname{乌龟：}恐怕没有。

\dialogname{阿基里斯：}真可惜，这可是一部了不起的作品。这表明你的音乐知识严重不足，龟兄。我敢肯定我们的朋友螃蟹一定非常熟悉这部作品，无论是把它正着演奏，还是倒着演奏。这个螃蟹几乎在所有方面都很聪明。呃，他至少要比任何活着的螃蟹聪明两倍，或许三倍，说不定------

\dialogname{乌龟：}你心尊蟹为大，你灵以螃蟹你的朋友为乐。

\dialogname{阿基里斯：}的确，我要赞美螃蟹，这是因为我恰好是他的一个崇拜者\ldots\ldots{}

\dialogname{乌龟：}别解释了。我也很钦佩他。说到螃蟹的崇拜者，我跟你说起过螃蟹不久前收到了一封崇拜者写来的莫名其妙的信吗？

\dialogname{阿基里斯：}我不信会有这事。谁发的信？

\dialogname{乌龟：}盖着印度邮戳，是一个我们过去谁也没听说过的人来的------我记得是叫衍奴玛拉。

\dialogname{阿基里斯：}怪了，从来不认识螃蟹的人怎么会给他来信呢------还有，他怎么会知道他的地址呢？

\dialogname{乌龟：}看来那准是个误把螃蟹当成个数学家的人。那封信里有一大堆研究结果，全是------哎，嘿！说到曹操，曹操就到！那不是螃蟹来了，正下山呢。

\dialogname{螃蟹：}再见！很愉快又跟你们聊了一次。哦，我看我最好还是走吧。我撑坏了------无论如何也不能再吃了。我自己刚才就在上面------很值得去。你们到过山顶上的那家茶馆吗？哦，阿基也在，你好啊。哈，哈，噢，噢，那不是龟兄吗。

\dialogname{乌龟：}你好，蟹兄。你爬上山顶去那家茶馆了吗？

\dialogname{螃蟹：}对呀，是的，我去了。你怎么猜着的？我很想吃一点他们那特有的宫廷点心------一种极美味的小吃。我饿极啦，简直能吞掉一只青蛙。哦，阿基也在，你好吗，阿基？

\dialogname{阿基里斯：}不错。

\dialogname{螃蟹：}妙极了！那么我就陪你们走走，但不插嘴了，你们继续讨论吧。

\dialogname{乌龟：}真是有意思透了，我刚要叙说叙说几周前你收到的那封由那个印度家伙寄给你的神秘的信------可是你现在就在这儿，我想还是让阿基听听你亲口讲这个故事吧。

\dialogname{螃蟹：}那好吧，事情是这样的。衍奴玛拉这个家伙显然从未受过任何正规的数学训练，可他却创立了一套他自己的新方法，推出一些新的数学真理。他的某些发现完全征服了我，不管怎么说，我以前从未见过类似的东西。比如说，他搞出一张印度地图，用不少于1729种颜色给它上色。

\dialogname{阿基里斯：}1729！你是说1729吗？

\dialogname{螃蟹：}是啊------怎么啦？

\dialogname{阿基里斯：}哦，你要知道，1729是个很令人感兴趣的数。

\dialogname{螃蟹：}真的吗，我没看出来。

\dialogname{阿基里斯：}具体地说，是这么回事，1729碰巧是我今天早晨去找龟兄时所乘的出租汽车的号码！

\dialogname{螃蟹：}真神！你知不知道你明天早晨去龟兄那儿时所要乘坐的汽车号码？

\dialogname{阿基里斯}难说，不过我想那个数字一定很大。

\dialogname{乌龟：}阿基对这些事有一种出色的直觉。

\dialogname{螃蟹：}是的。好了，我接着往下说吧。衍奴玛拉在信里还证明了每个偶素数都是两个奇数之和、方程

a\textsuperscript{n}+b\textsuperscript{n}=c\textsuperscript{n}　　对n=0

没有正整数解。

\dialogname{阿基里斯：}什么？这些数学老古董就这么一下子全解决了？他准是第一流的天才！

\dialogname{乌龟：}不过，阿基------你甚至一点都不怀疑吗？

\dialogname{阿基里斯：}什么？噢，对------怀疑。我当然怀疑。你不会认为我有那么傻，竟然相信老蟹收到过这样一封信吧？我可是什么当也不会上。龟兄，收到那封信的人一定是你！

\dialogname{乌龟：}噢，不，阿基，关于老蟹收到信这一节完全是真的。我的意思是说，你不怀疑这封信的内容------那些言过其实的论断吗？

\dialogname{阿基里斯：}凭什么我该怀疑？嗯嗯\ldots\ldots 对，我当然怀疑。我是个多疑的人，这一点你们俩现在应该非常了解了。让我相信什么事情很难，不管它多真多假。

\dialogname{乌龟：}说得好，阿基。你当然对自己的智力活动有极好的自我意识。

\dialogname{阿基里斯：}我的朋友，你们就没想过衍奴玛拉的这些论断可能不对吗？

\dialogname{螃蟹：}说老实话，阿基，我自己很保守也很正统，刚收到信时，我也想到了这个问题。其实，我当初就怀疑这是一个彻头彻尾的骗局。不过转念一想，又觉得没有什么人能仅凭想象就得出如此前所未闻的复杂结果。事实上，把它归结起来就是这样一个问题：``哪一个更可能呢：到底是一个独出心裁的骗子，还是一个天资过人的数学家？''不久我就明白了，其概率显然偏向于前者。

\dialogname{阿基里斯：}不过，你没有把他这些叫人吃惊的论断拿一个来直接核对一下吗？

\dialogname{螃蟹：}我干嘛要核对？那个基于概率的论证是我所想到的最令人信服的东西：没有一个数学证明能同它媲美。可是这位龟兄要坚持严格性，我最后也只好迁就，把衍奴玛拉的所有结果都核对了一遍。令人吃惊的是，他的每个结果都对。可我永远也不会知道他是怎么发现这些结果的。他肯定有一种惊人的、不可思议的东方式的洞察力，而我们这些西方人对此却一窍不通。目前在我看来，只有这种理论还多少能说得过去。

\dialogname{乌龟：}老蟹总是比我更容易接受那些神秘的或怪诞的解释。我完全有把握认为，不管怎么说，衍奴玛拉用他的方法能做到的事，在正统数学里都有与它完全相对应的东西。依我看，在数学研究上不会有什么方法能与我们已知的方法全然不同。

\dialogname{阿基里斯：}这倒是个有趣的观点。我看这与丘奇-图灵论题和其它一些有关的论题有些关系。

\dialogname{螃蟹：}好了，好了，在这样一个好天里，还是把这些学术问题放到一边去吧，让我们好好享受一下森林的宁静、鸟儿的啭鸣和洒在嫩叶新蕾上的阳光吧。嘿！

\dialogname{乌龟：}我赞成这一动议。说到底，世世代代的乌龟都非常迷恋大自然。

\dialogname{螃蟹：}就像世世代代的螃蟹一样。

\dialogname{阿基里斯：}你没带长笛来吧，老蟹？

\dialogname{螃蟹：}哪儿的话，我当然带了，走到哪儿带到哪儿。你想听上一两首吗？

\dialogname{阿基里斯：}在这种田园风光里，一定会令人心旷神怡的。你能凭记忆演奏吗？

\dialogname{螃蟹：}说来惭愧，我没这个本事，我只能看着乐谱演奏音乐。不过这没关系。我这只盒子里就有几段十分欢快的乐曲。

（他打开一只细长的盒子，拿出几张纸来，最上面一张上有下面的符号：

\gebfont{∀}a:\textasciitilde Sa=0

他把这张纸放到固定在长笛上的一个小谱架上，然后吹了起来，曲子很短。）

\dialogname{阿基里斯：}挺迷人的。\kaitismall{（盯着长笛上的那张纸，脸上露出一种疑惑的表情。）}挂在你长笛上的那个数论语句是干什么用的？\kaitismall{（螃蟹瞥了一眼长笛，又看了看乐谱，把头摇来摇去的，有点茫然。）}

\dialogname{螃蟹：}我不明白，什么数论语句？

\dialogname{阿基里斯：}``零不是任何自然数的后继。''就在那儿，长笛的谱架上。

\dialogname{螃蟹：}这是皮亚诺写的第三首长笛前奏曲。一共有五首来着。它们很好懂，同时也很吸引人。

\dialogname{阿基里斯：}我所不懂的是，数论语句怎么能当音乐演奏？

\dialogname{螃蟹：}我再说一遍，这不是数论语句------它是一支长笛前奏曲！你不想听点别的吗？

\dialogname{阿基里斯：}我会中魔的。

（螃蟹又把一张纸放在长笛上，这次阿基里斯瞧得更仔细了。）

哦，我看到你的眼光了，你刚才看着纸上的公式。你确实认为那是音乐符号吧？我敢发誓，它太像人们在形式化数论中所使用的那种符号了。

\dialogname{螃蟹：}真逗！要让我说，毫无疑问这是乐谱，而不是什么数论语句！当然，无论按数学这个词的哪种意义来说，我都不是什么数学家。你还想听别的曲子吗？

\dialogname{阿基里斯：}那当然。你还有什么？

\dialogname{螃蟹：}多得很。

\kaitismall{（他又拿出一张纸，挂在长笛上。纸上有下面这些符号}：

\textasciitilde{}\gebfont{∃}a:\gebfont{∃}b:(SSa·SSb)=SSSSSSSSSSSSS0

螃蟹演奏时，阿基里斯盯着它。）

好听吗？

\dialogname{阿基里斯：}嗯，的确是支悦耳的小曲。可我不得不说，我越看它越觉得它像数论。

\dialogname{螃蟹：}天哪！这只是我通常使用的音乐符号，仅此而已。我简直不知道你怎么会把那些非音乐的涵义硬塞进这些直接代表各种声音的符号里去。

\dialogname{阿基里斯：}也许你不会反对演奏一段我自己作的曲子吧？

\dialogname{螃蟹：}非常荣幸。你带来了吗？

\dialogname{阿基里斯：}没有，不过我想我完全可以自己写些曲子。

\dialogname{乌龟：}我必须告诉你，阿基，老蟹对别人的音乐作品可是个苛刻的批评家，所以，要是他对你的努力并不热心，你可别觉得失望。

\dialogname{阿基里斯：}多承关照。可我还是想试试\ldots\ldots{}

\kaitismall{(他写道：}

((SSS0·SSS0)+(SSSS0·SSSS0))=(SSSSS0·SSSSS0)

螃蟹拿过来看了一会儿，然后放在谱架上，吹了起来。）

\dialogname{螃蟹：}噢，满好。阿基，我喜欢稀奇古怪的节奏。

\dialogname{阿基里斯：}这段曲子的节奏有什么奇特之处吗？

\dialogname{螃蟹：}噢，当然啦，以你作曲家的眼光看它，似乎很平常。可对我的耳朵来说，这从3/3到4/4最后到5/5的节奏变化带了点异国情调。如果你还有别的曲子，我也很乐意演奏。

\dialogname{阿基里斯：}多谢。我过去从来没写过什么曲子，可我必须说，作曲这种事与我过去的想象完全不是一码事。让我再试一首。\kaitismall{（草草地写下了一行。）}

\textasciitilde{}\gebfont{∃}a:\gebfont{∃}b:(SSa·SSb)=SSSSSSSSSSSSSS0

\dialogname{螃蟹：}嗯\ldots\ldots 这不就是把我先前的那段曲子抄了一遍吗？

\dialogname{阿基里斯：}噢，不！我多加了一个S。你那行里有十三个，我这有十四个。

\dialogname{螃蟹：}噢，是这么回事。好吧（\kaitismall{他演奏起来，看上去挺严峻}）。

\dialogname{阿基里斯：}我希望你不至于不喜欢我的乐曲吧！

\dialogname{螃蟹：}阿基，恐怕你完全没有抓住我的那首曲子的精妙之处，你是照猫画虎。可我又怎么能指望你听上一遍就能理解呢？人们并不总能了解美的本质所在，很容易把一首曲子表面的东西错当成它的美，并且去模仿它。而美本身似乎是不可分析的，它藏在音乐的深处。

\dialogname{阿基里斯：}恐怕你这篇深奥的评论把我弄得有点晕头转向了。我知道我的乐曲够不上你的高标准，可我还没弄清我错在哪里。你也许能告诉我，我这个作品中的哪些方面有毛病吧？

\dialogname{螃蟹：}阿基，对你的作品可能会有的一种补救方法是：在末尾那一串S中再插进三个S（五个也行）。这将会造成一种精巧的不同寻常的效果。

\dialogname{阿基里斯：}我明白了。

\dialogname{螃蟹：}不过修改你的作品还有一些可以选用的方法。依我个人的看法，要是在前面再加一个弯号，会极富于感染力的。那样一来首和尾之间就会产生很好的平衡。你知道，一行中有两个弯会给乐曲一个轻快的小转折。

\dialogname{阿基里斯：}要是我同时采纳你的两项建议，会怎么样呢？那样曲子就成了：

\textasciitilde\textasciitilde{}\gebfont{∃}a:\gebfont{∃}b:(SSa·SSb)=SSSSSSSSSSSSSSSSS0

\dialogname{螃蟹}好了，阿基，听听下面这一课对你至关重要：决不要企图把太多的东西加到一支曲子里。总是有一个界限的，超过了它反而会弄巧成拙。这个例子就是如此。你这种把两条建议兼施并行的想法，并没有获得所期望的更多的美，恰好相反，却造成了一种不平衡，这下子把原来所有的妙处都弄没了。

\dialogname{阿基里斯：}你那支有十三个S，我这支有十四个S，这样两支十分相似的乐曲按你的音乐价值标准衡量，怎么会如此不同呢？除了这么一个细节，二者完全一样啊。

\dialogname{螃蟹：}天哪！你我两人的乐曲有着天壤之别。也许在这里无法用语言来表达那种心理感受。确实，我敢说不存在什么现成的规则，可以用来描述究竟是什么决定了一支曲子是不是美，而且永远也不会有这样的规则。美感只能存在于有意识的心灵之中，这种心灵靠生活经验的积累所达到的深度，超越了一切仅由一组规则所能做到的解释。

\dialogname{阿基里斯：}我将永远牢记你对美的本质的这种生动的阐述。这么说，是不是类似的评判也可以用于真理概念，是吧？

\dialogname{螃蟹：}毫无疑问。真和美密切相关，就像，就像\ldots\ldots{}

\dialogname{阿基里斯：}比方说，就像数学和音乐？

\dialogname{螃蟹：}嘿！真是英雄所见略同！你怎么就知道我在想什么呢？

\dialogname{乌龟：}阿基十分聪明，老蟹，你千万可别低估他的洞察力。

\dialogname{阿基里斯：}你是想说，在一个具体的数学语句是真还是假，与一支乐曲是美还是不美之间可能有一种可以想见的联系？还是说，这只是我的一种缺乏现实根据的猜想？

\dialogname{螃蟹：}你要是这么说我觉得就太走极端了。在我说到音乐和数学的相互关联时，你要明白，我用了一种非常比喻化的说法。至于说到具体的乐曲和具体的数学语句之间的直接关系，我可是极怀疑其可能性的。我冒昧地奉劝你不要花过多的时间去做这种呆想。

\dialogname{阿基里斯：}你当然是对的。那样做确实徒劳无益。也许我该集中力量再作一些新曲子，以增进我的音乐感受力，你愿意做我的辅导老师吗，老蟹？

\dialogname{螃蟹：}我非常高兴能帮助你向理解音乐的方向迈进。

(于是，阿基里斯拿起笔，显得非常聚精会神，写道：

∧OOa\gebfont{∀}\textquotesingle∨～∧∧:b+cS(\gebfont{∃}\gebfont{∃}=O∧→((～d)＜∨(\gebfont{∀}S·+(＞∨

螃蟹显得大吃一惊。）

你真的想叫我演奏这个------这个------这个莫名其妙的东西吗？

\dialogname{阿基里斯：}啊，请吧！

（螃蟹吹奏起来，显然十分吃力。）

\dialogname{乌龟：}真棒！真棒！阿基，你最喜欢的作曲家是约翰·卡奇吗？

\dialogname{阿基里斯：}实际上，他是我最喜欢的``不作曲家''，因为他写了许多部由静默构成的作品，例如《四分三十三秒》等等。不管怎么说，我很希望你喜欢我的音乐。

\dialogname{螃蟹：}听听这种毫无意义的不合谐音调，你们俩可能觉得挺有趣的，不过我向你们保证，对于一个敏感的作曲家来说，受这种令人难以忍受的、空洞无物的不和谐音以及毫无内容的节奏的折磨，一点也不会感到愉快。阿基，我认为你对音乐有某种良好的感觉。你前面的几支曲子都有可取之处，它们难道只是歪打正着的吗？

\dialogname{阿基里斯：}哦，请原谅，老蟹。我是在寻找你的这种音乐符号表达能力的极限。我想直接搞清楚我写下的某种类型的记号序列会得到哪一种声音效果，以及你如何评价不同风格的曲子。

\dialogname{螃蟹：}见鬼！我可不是一台自动音乐机，更不是一只专装音乐垃圾的垃圾桶。

\dialogname{阿基里斯：}对不起。不过我觉得，写过这几支曲子之后，我已经学到了很多东西，而且我敢说，要不是我有意想试试你，我现在肯定写的比我以前写的好。如果你肯再演奏一两首我的曲子，我想你会觉得我的乐感已经比以前好些了。

\dialogname{螃蟹：}那好，你写下来，我再给你一次机会。

（阿基里斯写道：

\gebfont{∀}a:\gebfont{∀}b:＜(a·a)=(SS0·(b·b))→a=0＞

随后螃蟹吹了起来。）

你行了，阿基。你好像已经完全恢复了你的音乐敏感力。这可是金不换啊！你是怎么创作出来的？我从来没听过这样的东西，它符合所有的和声律，而且还有某种------怎么说呢------不可思议的感染力。我把握不准。可正是由于这个原因，我非常喜欢它。

\dialogname{阿基里斯：}我猜你会喜欢的。

\dialogname{乌龟：}你给它起名字了吗，阿基？也许你可以叫它``毕达哥拉斯之歌''。你该记得，毕达哥拉斯及其弟子是最先研究乐音的人。

\dialogname{阿基里斯：}是的，确实如此。这是个极好的标题。

\dialogname{螃蟹：}毕达哥拉斯不也是最先发现了两个平方数的比值不可能是2吗？

\dialogname{乌龟：}我相信你正确。当时这被看作一项十分不吉利的发现。因为在此之前，从未有人认识到会有一些数------诸如2的平方根------不是整数的比值。这项发现深深地困扰着毕达哥拉斯学派的人们，他们认为，它揭示出在数的抽象世界中，有一个不容置疑的、怪诞的缺陷。不过，我到不认为这会造成天下大乱，以至于连中国的茶叶价钱都受影响。

\dialogname{阿基里斯：}哎，说起茶，前面不正是我们要去的那家茶馆吗？

\dialogname{乌龟：}是这样，我们三五分钟之内就可以到那儿了。

\dialogname{阿基里斯：}嗯\ldots\ldots 这段时间正好够我用口哨给你们吹一遍今天早上出租汽车司机用他的收音机放出的一只曲子。它是这样的：

\dialogname{螃蟹：}等一会儿，让我从盒子里拿几张纸把你的曲子记下来\kaitismall{（在盒子里翻了一阵，找出一张白纸）}，开始吧，我准备好了。

（阿基里斯用口哨吹了一段相当长的曲子。螃蟹手忙脚乱地跟着他的口哨记录）

你能不能把最后几个小节重吹一遍？

\dialogname{阿基里斯：}噢，当然可以。

（这么重复了两遍之后，这一阵子忙乱也就结束了。螃蟹得意地炫耀着它的抄件：

＜((SSSSS0·SSSSS0+(SSSSS0·SSSSS0))=((SSSSSSS0·SSSSSSS0)+(S0·S0))∧～\gebfont{∃}b:＜\gebfont{∃}c:(Sc+b)=((SSSSSSS0·SSSSSSS0)+(S0·S0))∧～\gebfont{∃}d:＜\gebfont{∃}c:(Sc+b)=((SSSSSSS0·SSSSSSS0+(S0·S0))∧\gebfont{∃}d:\gebfont{∃}d\textquotesingle:\gebfont{∃}e:\gebfont{∃}e\textquotesingle:＜～∨＜d=e∨d=e\textquotesingle＞∧＜b=((Sd·Sd)+(Sd\textquotesingle·Sd\textquotesingle))∧b=((Se·Se)+(Se\textquotesingle·Se\textquotesingle))＞＞＞＞

随后螃蟹自己吹奏起来）

\dialogname{乌龟：}这音乐别具风味，不是吗？我听着有点像印度音乐。

\dialogname{螃蟹：}我想它太简单了，不像是来自印度。不过当然啦，我对这类事情知道得极少。

\dialogname{乌龟：}好了，到茶馆了，我们坐在外边的阳台上好不好？

\dialogname{螃蟹：}要是你不反对，我宁愿进去，今天我可晒够太阳了。

（他们走进茶馆，找了一张舒适的木桌子，在旁边坐下，要了蛋糕和茶。一会儿，一辆小车装着看上去美味诱人的点心推了上来。他们各自拣了自己爱吃的）

\dialogname{阿基里斯：}老蟹，我很愿意知道你对我刚构思好的又一支曲子的看法。

\dialogname{螃蟹：}你能让我看看吗？写这儿吧，写在这块儿餐巾上。

阿基里斯写道：

\gebfont{∀}a:\gebfont{∃}b:\gebfont{∃}C:＜～\gebfont{∃}d:\gebfont{∃}e:＜(SSd·SSe)=b∨(SSd·SSe)=c＞∧(a+a)=(b+c)＞

螃蟹和乌龟颇感兴趣地琢磨着它。）

\dialogname{乌龟：}蟹兄，依你的看法，这一首也很美吗？

\dialogname{螃蟹：}这个，唔\ldots\ldots{}\kaitismall{（在椅子里挪了一下身子，看上去有点儿不自在）}

\dialogname{阿基里斯：}怎么样？是不是判断这支曲子是否美要比判断别的难些？

\dialogname{螃蟹：}啊，嗯\ldots\ldots 不，不是的\ldots\ldots 一点儿也不。是这么回事，哦\ldots\ldots 在我能说清对一支曲子有多喜欢之前，我确实得先听听它。

\dialogname{阿基里斯：}那好啊，吹它一遍！再不让我知道你是否觉得它美，我就活不下去了。

\dialogname{螃蟹：}当然，我极愿意为你演奏，唯一的问题是\ldots\ldots{}

\dialogname{阿基里斯：}你不能为我演奏？出了什么事？你怎么不太情愿？

\dialogname{乌龟：}难道你还不明白吗？阿基，老蟹觉得满足你的要求将会显得极不礼貌，这会打扰这家店铺的顾客和店员。

\dialogname{螃蟹}是这么回事，我们没有权力把我们的乐曲强加给别人。

\dialogname{阿基里斯}噢，咳！可我是多么想知道你对这支曲子的看法啊？

\dialogname{螃蟹：}咳！可算躲过去了！

\dialogname{阿基里斯：}你在说什么呢？

\dialogname{螃蟹：}哦！没什么，只不过是那儿的那个侍者，他让另一名侍者撞了一下，差一点把一整壶茶都倒在一位女士的衣襟上，所以我说，差点儿躲不开了，龟兄，你说呢？

\dialogname{乌龟：}我想说，茶不多了，你不同意吗，阿基？

\dialogname{阿基里斯：}噢，是喝得差不多了。

\dialogname{螃蟹：}一点儿不错，不知你们二位意下如何，我觉得我好像是该走了，哎，我们刚才的情景，让我想起我看过的一个片子。

\dialogname{阿基里斯：}什么？你的意思是说我们都是骗子？

\dialogname{螃蟹：}这可是你说的，阿基。

\dialogname{阿基里斯：}我明白了，噢，我得把这记住。

\dialogname{螃蟹：}今天下午真快活，阿基，我衷心地希望我们改天再交流更多的音乐作品。

\dialogname{阿基里斯：}我热切期待着那一天，老蟹，好，再见。

\dialogname{乌龟：}再见，老蟹。

（螃蟹沿着山的另一边下去了。）

\dialogname{阿基里斯：}走了一个才华横溢的伙伴\ldots\ldots 依我看他至少比所有活着的螃蟹聪明四倍，也许甚至有五\ldots\ldots{}

\dialogname{乌龟}我心尊蟹为大，我灵以螃蟹我的朋友为乐\ldots\ldots{}

\phantomsection\label{Chapter17.xhtml}{}

第十七章

\chapter{丘奇、图灵、塔斯基及别的人}

形式的和非形式的系统

现在可以展开本书的一个主要论题了，那就是：思维的每一个方面，都可以看成是从较高的层次上描述的一个位于较低层、受某些简单的乃至形式的规则支配的系统。当然，这个``系统''就是大脑------除非谈论的是在另一种媒质（比如一台计算机的电路）中流动的思维过程。形象地说，这是一个支撑着``非形式系统''的形式系统，那个非形式系统能一语双关、能发现数字中的模式、会忘记人的姓名、会走臭棋、等等。我们从外部看到的是它的非形式化的、公开的、软件的层次。作为对照，它还有一个形式化的、隐蔽的、硬件的层次（或叫``基质''），这是一部令人望而生畏的机器，它根据物理地嵌入其中的某些确定的规则以及与之密切相联的信号输入，在各个状态之间进行转换。

像这样一种对大脑的看法，肯定会有很多哲学的和其它方面的推论。我打算在这一章里详细地说明这样一些推论。其中之一是说，这种看法似乎蕴涵着：从本质上讲，大脑是某种``数学''对象。其实，这充其量也不过是看待大脑的一种十分笨拙的办法。因为即便在技术意义或抽象意义上说大脑是某种类型的形式系统，仍然无法摆脱下面的事实：数学家只与那些简单的、考究的形式系统打交道，在那些系统中，一切东西都有着极为清晰的定义。而大脑则与此大相径庭，它的上百亿个不完全独立的神经元彼此是近乎随机地相联接的。所以，数学家决不会去研究实际的大脑网络。如果你把``数学''定义为数学家喜欢作的事情，那么大脑的性质就不是数学性质了。

要想了解像大脑这样的复杂系统，唯一的方法是在越来越高的层次上对之组块，因而每一步都要损失一些严格性。出现在最顶层上的，是一个``非形式系统''，它要服从许多复杂到我们找不到合适的词汇去思考的规律。而这正是人工智能的研究所想找到的东西。它具有与数学研究大不相同的味道。尽管如此，还是与数学有个松散的联系：从事人工智能研究的人，通常都有很强的数学背景，并且，数学家有时也会对自己大脑的工作感兴趣。下面这段引自斯坦尼斯拉夫·乌兰姆的自传《一位数学家的奇遇》的文字，正好说明了这一点。

依我看，在揭示\ldots\ldots 联想的本质时，有了计算机提供的实验手段，就可以做更多的事情了。这种研究必定要把概念、符号、符号组成的类、类组成的类等等进行分级，就像在研究数学和物理结构的复杂性时所做的那样。

人脑的思维过程一定有某种窍门------某种递归公式。一组神经元能自动开始工作，有时用不着外来剌激，而是靠一个具有递增模式的迭代过程。它在大脑中游来荡去，而且，它的发生方式必定依赖于对类似模式的记忆。\footnote{}

直觉和值得赞美的螃蟹

人工智能{[}Artificial Intelligence{]}常被简称为``AI''。如果要我来解释这两个字母的涵义，我会说它也可以被理解为``人工直觉''{[}Artificial Intuition{]}或``人工意象''{[}Artificial Imagery{]}。人工智能的目标是要弄清，当人的大脑在一种十分复杂的环境中不动声色地从大量可能性中选择哪一种最为合理时，会发生什么事。在现实生活的许多场合，演绎推理之所以不适用，并不是因为它会得出错误答案，而是因为它会得到过多的正确然而无关的断言。为了满足推理本身的充分性，需要同时考虑的东西也就太多了。请看这样一段很短的对话：

``那天，我从报上读到------''

``噢------你曾读报来着？那就可知你有眼睛。或者说至少有一只眼睛。或者更确切地说你那时至少有一只眼睛。''

这就需要一种辨别的能力------``在这里，什么东西重要，什么不重要？''与此相关联的，是一种对简洁的感受力和对美的感受力。这些直觉是从哪来的呢？它们是怎样从作为基础的形式系统中显现出来的？

《的确该赞美螃蟹》里显示了螃蟹智慧的一些非凡力量。按他自己对这种力量的说法，他只不过是在听音乐，并区分优美的音乐和不优美的音乐（显然，他有一条明确的分界线）。然后，阿基里斯找到了另一种描述螃蟹能力的方法：螃蟹把数论陈述分为真的和假的两个范畴。可是螃蟹坚持说，如果他真的做到了这一点，那也纯属偶然，因为按照他自己所承认的，他在数学方面是无知的。然而更加使阿基里斯不解的是，螃蟹的行为看上去直接违背了阿基里斯所熟悉的一项备受称颂的元数学结果：

\heititext{丘奇定理}：没有一个切实可靠的方法总能区分开TNT的定理和非定理。

这是美国逻辑学家丘奇于1936年证明的。与此密切相关的是下面这个结果，我称之为

\heititext{塔斯基-丘奇-图灵定理}：没有一种切实可靠的方法总能区分开真的数论语句和假的数论语句。

丘奇-图灵论题

为了更好地理解\heititext{丘奇定理}和\heititext{塔斯基-丘奇-图灵定理}，我们先得描述一下作为它们的基础的一种思想，那就是丘奇-图灵论题（通常也叫``丘奇论题''），因为丘奇-图灵论题确实是数学、大脑及思维的哲学中最重要的概念之一。

实际上，像泡茶一样，丘奇-图灵论题也能有各种不同的浓度。所以，我要以好几种形式来表述它，并且看看这些不同的形式各自蕴涵着什么。

第一种形式听起来十分幼稚------事实上几乎是空洞无物：

\heititext{丘奇-图灵论题，同义反复形式}：数学问题只能通过数学推演来解决。

自然，整体的意义寓于作为组成成分的各术语的意义之中。所谓``数学问题''，我指的是判定某个数是否具有一个给定的算术性质。看得出，借助于哥德尔配数法和相应的编码手段，几乎每一个数学分支的每一个问题都能归入这种形式。所以，``数学问题''一词仍保持其通常含义。``数学推演''又怎么讲呢？当你试图确定一个数是否具有某种性质时，似乎只有一遍又一遍地结合使用很少量的几个运算------加法、乘法、验证相等与不等。这就是说，循环地使用这些运算，好像就是使我们得以深究数的世界的唯一工具。要注意``好像''这个词。这是有关丘奇-图灵论题的关键字眼。我们可以给出一个``修订本''：

\heititext{丘奇-图灵论题，标准形式}：假设有一种方法，一个有感知能力的生物可以根据这种方法逐个把数分成两类。又假定这种方法总能在有穷时间内得出答案，而且对于给定的数，这种方法总给出相同的答案。那么：存在一个有终止的FlooP程序（即一般递归函数），它给出的答案恰好与这个有感知能力的生物的方法所得到的答案一样。

中心假设说穿了就是：把数分成两类的任何一个心智过程都可以用FlooP程序来描述。直觉上的一个信念是，这种心智过程中没有使用什么超出了FlooP能力的工具，而且使用这些工具的方式也没有超出不限次数的迭代（FlooP允许这么做）。丘奇-图灵论题不是数学定理意义上的一个可证的事实------它是有关人脑中所使用的过程的一个假设。

大众过程形式

此刻，有些人可能会感到，丘奇-图灵论题的这种形式断定的东西太多。这些人可以提出如下的反对意见：``也许有像螃蟹这样的人------他对数学有着几乎是神奇的洞察力，但像其他人一样，他自己对这种特殊能力也一无所知------那么这种人的心智机制大概就能执行在FlooP中找不到对应物的运算。''这种想法就是说：或许我们会有一些下意识的潜力，可以做出超越意识过程的事情------这些事情大概无法用基本FlooP运算表示。对这些反对派，我们将给出该论题的一个弱形式，它在大众的和个人的心智过程之间作了区分。

\heititext{丘奇-图灵论题，大众过程形式}：假设有一种方法，一个有感知能力的生物可以根据它逐个把数分成两类。又假设这种方法总能在有穷的时间内得出结果，并且对给定的数它总能给出相同的结果。同时，还要假定这种方法可以通过语言由一个有感知能力的生物不走样地传达给另一个有感知能力的生物。那么，存在一个有终止的FlooP程序（即一般递归函数），它给出的答案恰好和这个有感知能力的生物的方法所得到的答案一样。

这就是说，大众接受的方法受``FlooP化过程''的管辖，但对个人特有的方法，它却什么也没断定。当然它并没有说个人特有的方法都是FlooP力所不及的，不过至少还是留有这样的余地。

湿利尼吠萨·拉玛奴衍

作为反驳任何一种强形式丘奇-图灵论题的证据，我们来看看二十世纪初叶著名的印度数学家湿利尼吠萨·拉玛奴衍（1887-1920）（图105）这个实例。他出生于印度最南部的泰米尔邦，高中时学过一点数学。一天，一个认为拉玛奴衍有数学天分的人，介绍他看一本略为过时的分析教科书，拉玛奴衍立即狼吞虎咽起来。于是他开始突入分析世界，二十三岁前后，他就已作出了一些他认为值得花费时间的发现。他不知道该向谁求教。可不知怎么回事，有人就谈到了远在英国的一位数学教授------哈代。于是拉玛奴衍把他最好的一些结果全都装进一个纸袋，连同朋友们帮他用英文写的一封说明信一起，寄给事先毫不知情的哈代。下面这段话录自哈代对自己接到这包东西时的描述：

\ldots\ldots 一下子就能看出，拉玛奴衍一定还有很多更普遍的定理，肯定还藏着很多东西\ldots\ldots{}（有些公式）完全使我折服了。以前我从未见过任何一种哪怕是稍许类似的东西。仅仅看一眼就足以说明只有最高水平的数学家才写得出来。它们必定是真的，因为假如它们不真，不可能会有人具备杜撰它们的想象力。最后\ldots\ldots 作者肯定百分之百地诚实，因为大数学家比具有如此难以置信的技巧的窃贼或骗子更常见。\footnote{}

\pandocbounded{\includegraphics[keepaspectratio]{./media/OEBPS/Images/img_105.jpg}}

图105(a)．湿利尼吠萨·拉玛奴衍

这次通信的结果是，拉玛奴衍根据哈代的建议，于1913年来到英国，随后就开始了一场以拉玛奴衍三十三岁早逝于肺结核为结局的紧张合作。

\pandocbounded{\includegraphics[keepaspectratio]{./media/OEBPS/Images/img_105-1.png}}

图105(b)．拉玛奴衍发现的神奇的印度旋律之一

拉玛奴衍具有一些使他有别于大多数数学家的极不寻常的特点。其中之一是他缺乏严格性。他常常会简单地陈述一个结果，然后坚称这个结果完全来自自己模糊的直觉，这与那个有意识的进行研究探索的领域是格格不入的。事实上，他常说女神娜摩吉利在梦中给他灵感。这种情形一再发生。使之更加神秘------也许甚至是使之充满了神奇特性------的是，拉玛奴衍的这些``直觉定理''中有很多是错的。世上有一种奇怪的悖论效果：一件你认为不会促使人们轻信，而是会使易轻信的人变得开始怀疑的事情，实际上却具有相反的效果。因为它击中了轻信者心理上的弱点，用暗示人类本性中一些令人困惑的非理性侧面的办法吊起了他们的胃口。这就是拉玛奴衍出错时所发生的情形：很多受过教育的人，以一种强烈的倾慕之心，相信某些被看成拉玛奴衍的直觉能力的东西，认为它证明了有一种神奇的、洞察真理的力量存在，而他难免出错这一事实甚至可能是加强了------而不是削弱了------这种信念。

当然，没人否认拉玛奴衍来自印度最落后的地区。在那里，行乞化缘和其它一些古里古怪的印度习俗已经沿袭了数千年，而且还将继续下去，其频繁程度大概要超过高等数学课。拉玛奴衍的那些偶然的失误并没有提醒人们认清他只是个人，反倒自相矛盾地激起了一种想法：他的错误总有某种``更深刻的正确性''------某种西方人的心智或许不可企及的``东方的''正确性。这是多么诱人而且近乎无法抵御的想法啊！即便是哈代------他该是第一个否认拉玛奴衍有特异功能的人------在谈到拉玛奴衍的某次失败时也说：``从某种意义上说，我并不认为他的这次失败就不比任何一次成功更精彩。''

拉玛奴衍的数学个性中另一个突出的特点（照他的同事李特伍德的说法），就是``与整数的友谊''。相当多的数学家都在不同程度上具有这样的特点，而拉玛奴衍则登峰造极。有两则轶事可以说明他的这种特殊能力。第一则为哈代所讲：

记得有一次他生病住在普特涅，我去看他，乘的是一部号码为1729的出租马车，我注意到对这个数我翻不出什么花样，只希望它不是什么凶兆。``不''，他听后答道，``这是个很有意思的数：它是能用两种不同方法表示成两个数的立方和的最小的数。''自然，我就问他是否知道对应于四次方的这样一个问题的答案。他想了一会，回答说他一下子找不到例子，第一个这样的数一定很大。\footnote{}

实际上，对于四次方，答案是：

635318657=134\textsuperscript{4}+133\textsuperscript{4}=158\textsuperscript{4}+59\textsuperscript{4}

而关于二次方的类似问题则要容易的多，读者可以试着动手解一解，说不定会觉得很好玩的。

考虑一下哈代为什么会直接跳到四次方去，其实是很有意思的。因为毕竟是存在好几种合理而又自然的方式，能在不同的维度上推广方程

u\textsuperscript{3}+v\textsuperscript{3}=x\textsuperscript{3}+y\textsuperscript{3}

例如，把一个数用三种不同方法表示成两个立方之和的问题：

r\textsuperscript{3}+s\textsuperscript{3}=u\textsuperscript{3}+v\textsuperscript{3}=x\textsuperscript{3}+y\textsuperscript{3}

或者，可以用三个不同的立方数：

u\textsuperscript{3}+v\textsuperscript{3}+w\textsuperscript{3}=x\textsuperscript{3}+y\textsuperscript{3}+z\textsuperscript{3}

或者，甚至可以同时对所有的维度都推广：

r\textsuperscript{4}+s\textsuperscript{4}+t\textsuperscript{4}=u\textsuperscript{4}+v\textsuperscript{4}+w\textsuperscript{4}=x\textsuperscript{4}+y\textsuperscript{4}+z\textsuperscript{4}

即：``能用三种不同方法表示成三个数的四次方之和的最小数是什么？''不过，从某种意义上来说，哈代的推广是``大多数数学家所倾向的''。这样一种数学美学的意义最终能不能得以程序化呢？

另一则轶事出自拉玛奴衍的同胞拉戈那罕为他所写的传记，传记中称这件事为``拉玛奴衍的闪念''，那是拉玛奴衍在剑桥时代的印度朋友玛哈拉诺比斯博士讲述的。

又有一次，我去他家吃午饭，当时第一次世界大战已经打了一些时候了。我手里拿着一份《海滨杂志》月刊，当时这家杂志常发表一些谜题供读者解答。拉玛奴衍正把锅坐在火上，在里面搅拌着什么东西准备午饭，我则坐在桌子旁边翻阅杂志。有一个涉及两个数之间的关系的问题吸引了我。问题的细节已经忘了，不过我还记得它的类型。是说两名英国军官住在巴黎一条长街上的两所房子里，其门牌号码有某种特殊的联系，问题就是：求出这两个数。题目并不难，我通过尝试、订正，花了几分钟就解出来了。

\dialogname{玛哈拉诺比斯}有个问题你。

\dialogname{拉玛奴衍：}什么问题？说吧。\kaitismall{（他继续在锅里搅着。）}

我把《海滨杂志》上的问题读了一遍。

\dialogname{拉玛奴衍：}请把答案记下来。\kaitismall{（他念出一个连分数。）}

连分数的第一项是我已经得到的解，而如果街上的门牌号数可以无限制地增长下去的话，后续各项就将依次成为这个问题的后续的解。

我感到十分惊奇。

\dialogname{玛拉哈诺比斯：}你一眨眼之间就得到这个解了吗？

\dialogname{拉玛奴衍：}我一听到这个问题就明显感到，解显然是一个连分数，于是我就想：``是个什么样的连分数呢？''然后答案就在我脑子里出现了。就是这么简单。\footnote{}

拉玛奴衍逝世之后，常有人向他的最密切的合作者哈代问起，拉玛奴衍的思维方式是否有什么超自然的或异国情调的奇异因素。哈代作了如下的评述：

我曾不时地被人问起：拉玛奴衍是否有什么特殊的奥秘？他的方法是否与其他数学家有实质上的不同？他的思维方式是否确实有什么不正规的东西？我无法以任何自信的或有说服力的方式来回答这些问题，不过我并不相信那些假设。我的信念是，从本质上讲，所有的数学家都在用同一性质的方法思维，拉玛奴衍也不例外。\footnote{}

哈代在这里其实是在陈述他自己对丘奇-图灵论题的一种说法。我来发挥一下，就是：

\heititext{丘奇-图灵论题，哈代形式}：从本质上讲，所有的数学家都同构。

这还没有把每个数学家的数学潜力都等同于一般递归函数的潜力。要想做到这一点，你所需要做的事情也就是设法说明某个数学家的心智能力与递归函数的能力是一样的。这样，如果你相信丘奇-图灵论题的哈代形式，那就可以知道全体数学家都是如此了。

此后，哈代将拉玛奴衍和一些计算奇才作了比较：

他的记忆力和计算能力很不寻常，不过也不能说它们``反常''。当他需要做两个大数的乘法时，他是用普通的方法，只是他可以用不同寻常的速度和准确度来做好它，但也并不比任何一名天生迅速而又有计算习惯的数学家更快。\footnote{}

哈代描述了他所发现的拉玛奴衍突出的智力特征：

由于他的记忆、耐心和计算才能，他把综合能力、对形式的感觉能力和对假设的快速修订能力结合了起来，这经常出人意料，并使他在那个时代中自己的那个领域内无可匹敌。\footnote{}

这段文字中提到的那几种能力，依我看是对一般智能的某些最微妙的特征的精彩刻划。最后哈代颇怀眷恋地总结说：

（他的工作）没有最伟大的工作所必须具备的那种简洁性和必然性。如果不那么古怪，它就会更伟大一些。他的工作具有一种不可否认的天赋------意义深刻和无往不胜的独创性。如果他早在青年时代就被发现并受到一些训练，他就可能成为一名伟大的数学家，他就会完成更多新的、而且无疑是极为重要的发现。另一方面，那也将使他变得不那么像拉玛奴衍，而更像一名欧洲教授。这也许是得不偿失的。\footnote{}

哈代在说起拉玛奴衍时所用的浪漫方式，充分表达了他对拉玛奴衍的珍重。

``心算家''

还有一类人，他们的数学能力似乎得不到合理的解释------那就是所谓``心算家''，这些人能在心中（或别的什么地方）一眨眼之间完成复杂的计算。生于1824年，卒于1861年，并为欧洲的好几个政府雇来作计算工作的戴斯就是一个突出的例子。他不仅能心算两个一百位数的乘法，而且有一种不可思议的数量感。也就是说，他可以不用数数就``分辨出''一块地里有多少只羊或者一个句子里有多少个词等等，可以多到30个左右------这和我们大多数人显著不同，我们的这种感觉只在6个以内时才是可信的。

我不打算描述很多这类让人眼花缭乱的、但却有案可稽的``瞬息计算者''，因为这不是我此处的目的。不过我觉得，消除那种认为他们是用某种神秘的、不可分析的方法计算的想法，是重要的。尽管实际情况常常是：这种奇才的计算能力远远超过他们解释自己的结果的能力，不过偶尔也会出现一个具有其它智力天赋，同时也具有这种惊人的数字能力的人。根据这些人的自省以及广泛的心理学研究已经弄清：在瞬息计算者进行运算的过程中，并没有发生什么超自然的事情，而只是他们的心智非常自信地加速越过了各个中间步骤，就像天生的运动员迅速完美地从事一项复杂运动时那样。他们不是依赖某种顿悟般的瞬息闪念来得到答案（尽管在主观上他们当中有些人似乎感到是这么一回事），而是------像我们这些人一样------由一系列计算得到的，也就是说用FlooP（或BlooP）之类的过程得到的。

此外，能够说明并不存在什么``通往上帝的热线''的一条最明显的线索，是一个简单的事实：当用到的数字变大时，答案就出来得慢。可以设想，如果是上帝或``神谕''在提供答案，那么当数字变大时就不该变慢。人们也许可以找到一个更好的解释，说明瞬息计算所花的时间是如何随用到的数字的大小及用到的运算种类而变化的，并由此来推定所用的算法的某些特征。

丘奇-图灵论题的同构形式

这最终把我们带到加强了的丘奇-图灵论题标准形式：

\heititext{丘奇-图灵论题，同构形式}：假定有一种方法，一个有感知能力的生物可以利用它逐个把数分成两类。又假定这种方法总能在有穷的时间内得到答案，而且对于一个给定的数总是给出相同的答案。那么，存在某个有终止的FlooP程序（即一般递归函数），它给出的答案恰好与这个有感知能力的生物的方法所得到的答案一样。而且，这个心智过程与这个FlooP程序在下述意义上同构：在某个层次上，计算机和大脑各自执行的那些步骤之间存在一个对应。

应当注意，在这里，不仅结论加强了，而且也抹去了吞吞吐吐的大众过程形式中的那个有关通迅的附加条件。这个大胆的形式正是我们现在所要讨论的。

简而言之，丘奇-图灵论题的这种形式断言：当一个人计算某个问题时，他的心智活动可以同构地在某个FlooP程序中得到反映。不过读者应该清楚，这并不意味着大脑实际上是按照一个用FlooP语言写下的、尽是些BEGIN、END、ABORT之类的FlooP程序而运行的------完全不是。这只不过是说各个步骤之间的次序与FlooP程序中相应的次序完全相同，而且计算的逻辑结构也可以在FlooP程序中得到反映。

现在，为了弄清这一思想的含义，我们得在计算机以及大脑中区分出不同的层次，否则这就会被误认为是十足的胡说。很可能，在人的头脑中进行的这些计算步骤都处于最高层次，并从较低的层次上得到支持，因而最终被``硬件''所支持。所以如果我们说到同构，那就意味着我们心照不宣地做了一个假设，即能把最高层隔离开来，从而使我们得以脱离其它层次来讨论这里发生的事情，并能把这个顶层对应到FlooP。说得更严格一些，这个假定就是说：存在一些软件实体充当着种种数学思维的角色，它们以某些能够严格反映为FlooP的方法起着作用（见图106）。使这些软件实体得以实现的东西，是第十一章、第十二章以及《前奏曲，蚂蚁赋格》中讨论过的整个基础结构。对大脑和计算机的较低层次（例如神经元和位）上的同构作用，则没下任何断言。

\pandocbounded{\includegraphics[keepaspectratio]{./media/OEBPS/Images/img_106.png}}

图106．联系着自然数、计算机和人脑的同构关系

自然数的性质能在人类大脑和计算机程序中得到反映，这两个不同的描述因而能在一个适当的抽象层次上彼此对应

如果不拘泥于字面，丘奇-图灵论题的同构形式之所以能为人理解，是由于它说一个心算家在计算（比如说π的对数）时所做的事情，同构于一部袖珍计算器在计算它时所做的事情。这里的同构是在算术步骤的层次上成立，而不是在更低的层次------神经元和集成电路之间------成立的。（当然，计算任何东西都可以有不同的方法，但不妨假定那部袖珍计算器------如果不提那个人------可以按任意特定的方式完成所要求的计算。）

对关于现实世界的知识的表示

当我们所考虑的范围是数论的时候，这件事看上去的确是合情合理的，因为此时发生各种事件的那个论域既小又清晰。而论域本身的边界、居民、法规也都是良定义的，就像用标准几何图形画出的迷宫。这样的一个天地远不如我们生活于其中的那个无边无际且定义模糊的天地来得复杂。一个数论问题一经提出就完全是自足的了，然而现实问题却不然，根本不能绝对有把握地把它与现实世界的任何一个部分隔离开来。比方说换一只烧蹩了的电灯泡的任务，大概免不了要拿一个板凳来；而这又可能意外地弄撒一盒药，于是又不得不去扫地以免孩子误食撒了的药丸；等等、等等。药丸、凳子、小孩及灯泡本是现实世界中风马牛不相及的几个部分------然而某些日常事件却会把它们紧密地联系起来。这里还没有说到其它微小变化会使预期的事情变成什么样。总之，如果给你一个数论问题，那么要想解决它，就决不致于手忙脚乱地考虑什么药丸、小孩、凳子乃至扫帚之类的无关事物（当然，你对这些东西的直观知识可能有助于你下意识地生成一种内心的图像，从而有助于用几何术语想象问题的解答------不过这是另外一回事了）。

鉴于世界的这种复杂性，很难设想能有一部小型的袖珍计算器，当你按下一些印有``小孩''、``凳子''、``灯泡''之类字样的按钮时，就能回答输入它的种种问题。实际上，现在已经证实，要想有一台大型的高速计算机，能回答一些在我们看来是现实世界中极为简单的一个小领域中的问题，都是极为困难的。为做到``理解''，看来就得用一种高度综合的方式把大量知识都考虑进去。我们可以把现实世界中的思维过程比做一棵树，其可见部分坚实地长在地面之上，但却性命攸关地依赖于在地下四处延伸、看不见的根系。根使它稳定并提供给它营养。这里，根象征着那些复杂过程，这是一些发生在大脑的意识层次之下的过程------其影响遍布我们自己觉察不到的思维方式中。这就是第十一章和第十二章讨论过的``符号的触发模式''。

关于现实世界的思维与我们作两个数的乘法时所发生的事情很不相同。作乘法时一切东西都在``明面上''，随时可以检验。在算术中，可以``撇出''顶层，并使用很多不同种类的硬件等效地完成它，例如用机械加法器、袖珍计算器、大型计算机、人脑等等。而这正是丘奇-图灵论题所说到的全部东西。但是，谈到现实世界中的理解问题，好像并没有什么简单的方法``撇出''这个顶层而单用程序来实现：那个``符号的触发模式''太复杂了。必须得有若干个层次使思维得以``渗透''和``冒出''。

特别地------这就回到了第十一章和第十二章的主题------大脑中对现实世界的描述尽管在某种程度上根植于同构，但却包含有某些在外部世界根本没有对应物的成分。也就是说，对于这种同构，存在有比表示``小孩''、``扫帚''等东西的简单意识结构多得多的东西。这些符号全都存在，这毫无疑问------不过它们的内部结构极为复杂，而且在很大程度上不便于有意识地检查核对。不仅如此，若还想把一个符号的内部结构的各个方面，映射到现实世界的什么具体性状上，那将是徒劳无功的。

难以撇出的过程

由于这个理由，大脑看起来就像一个很有特色的形式系统了。因为在它的最底层------神经元的层次，也就是由一些``规则''控制并改变其状态的层次------原始的元素（神经元的发射或者更低层的事件）可以没有解释。而在顶层，却出现了有意义的解释------即从我们曾称为``符号''的那个神经元活动的大``云雾''到现实世界上的一个映射。这里与哥德尔的构造有几分相似。按照这种相似性，一个高层的同构可以使高层次的意义注入符号串之中。不过，在哥德尔的构造中，较高层次的意义是``骑''在较低层次的意义上的------也就是说，只要引入了哥德尔配数法的概念，它就可以从较低的层次上导出。而在大脑中，神经元层次上的活动并不附属于现实世界的解释，它们根本不是在模仿什么东西。它们纯粹只是实现较高层次的一种基质，就像袖珍计算器中的晶体管，其作用仅仅是支持计算器的反映数字世界的活动。言下之意就是说，没有什么方法能把最高的层次撇出，而用程序造出一个同构拷贝。如果要反映大脑中对现实世界的理解过程，那就必须反映正在发生的某些较低层次的事情------``大脑的语言''。而这并不一定意味着我们非得一直下到硬件层次，尽管可能到头来的确如此。

如果为了达到对``在外面''的东西有一个``智能的''（即类似于人的）内部表示，而开发一个程序，那么在开发过程中，有时可能要被迫使用一些结构和过程，它们不容许有任何直接了当的解释------也就是说，它们不能直接映到现实世界的成分之上。要想理解该程序的这些较低的层次，只能靠它们对高于自己的那些层次的催化关系，而不是靠它们与外部世界之间所具有的某种直接联系（这种思想的一个具体形象曾由食蚁兽在《蚂蚁赋格》中提示过：试图只凭笔划就理解一本书的``难以言状的噩梦''）。

依我看，一旦一些包括想象和类比的过程变成程序的重要成分时，这种从概念上把握系统的多层建筑就变成必不可少的了------这与那些被看成是严格地执行演绎推理的过程恰成对照。执行演绎推理的过程，从本质上说都能在一个单一的层次上程序化，因而根据定义就是可以撇出的。所以，按照我的假定，想象和类比的思维过程本质上都需要有若干层次的基质，因而本质上都不可撇出。此外我还相信：就是在这些地方，创造性开始浮现出来------这将蕴涵着：创造性本质上依赖于某种``不可解释的''低层事物。作为类比思维的基础的那些层次无疑是极为吸引人的，下面两章会给出有关它们性质的一些思考。

简化论的几个信念

考察大脑中较高层次和较低层次之间关系的一种方法是：装配一个神经元网络，在一个局部（神经元到神经元）的层次上，该网络表现得与大脑中的一个神经元网络没有什么区别，但它完全没有较高层次的意义。较低的层次由相互作用的神经元组成，这一事实并不见得一定能导出有什么较高层次的意义------就像一碗``笔划汤''里有笔划这一事实，并不能导出在碗里游泳时就一定能碰到有意义的句子一样。高层意义是神经元网络的一种带有随意性的特征------一种作为（进化中）环境压力的后效而出现的东西。

\pandocbounded{
\includegraphics[keepaspectratio]{./media/OEBPS/Images/img_107.png}}

图107．大脑的神经和符号运动

大脑的符号层漂浮于神经元活动之上，从而反映了世界。不过能在计算机上模拟的那种神经元活动本身并不能产生思维，那得靠组织中的一些更高的层次。

图107说明了意义的较高层次的出现是有选择余地的。向上的箭头表示一种基质在出现时可以没有高层意义，但反之不然：较高的层次必须从一个较低层次的性质中导出。

这个图表明，可以用计算机模拟一个神经元网络。这在原则上是行得通的。不论网络多么复杂，只要单个的神经元的行为能用计算机可执行的计算来描述就行了。这是一个微妙的假设，没有多少人甚至会想起来怀疑它。然而，这却是``简化论的信念''之一，可以把它看成丘奇-图灵论题的``微观形式''。下面就把它明确地叙述一遍：

\heititext{丘奇-图灵论题，微观形式}：一个生物体的各组成部分的行为能用计算机来模拟。也就是说，任何元素（为典型起见，就假定是一个细胞）的行为，都能用一个FlooP程序（即一般递归函数）------在给定该元素的内部状态和外部环境的一个足够精确的描述之后------计算到任意精确的程度。

丘奇-图灵论题的这种形式是说：大脑过程------尽管拥有更多的组织层次------并不比（例如）胃的过程具有更多的神秘性。今天，要是提出人消化食物不是通过普通的化学过程，而是通过某种神秘的、魔术式的``同化''过程，那就会使人感到不可思议了。丘奇-图灵论题的这种形式，不过就是把这种常识性的道理推广到大脑过程。简言之，这等于说，相信大脑在原则上是用一种可以了解的方式活动的。这是简化论的一条信念。

微观的丘奇-图灵论题的一项推论，是下面这样一个十分简洁的新微观形式：

\heititext{丘奇-图灵论题，简化论形式}：全部的大脑过程都可以从一个可计算的基质中导出。

这句话大概是能支持``最终可能实现人工智能''这种观点的最强有力的理论基础。

当然，人工智能的研究并不是以模拟神经元网络为目的的，因为它建立在另一种信念之上，即：可能有一些意义重大的智能特征是漂浮在一些与生物大脑的基质完全不同种类的基质之上的。图108表明人工智能、自然智能和现实世界之间的大致关系。

\pandocbounded{\includegraphics[keepaspectratio]{./media/OEBPS/Images/img_108.png}}

图108．``撇出''大脑的最高层次

使人热衷于从事人工智能研究的关键，在于这样一种观念：心智的符号层次能从它们的神经元基质上被``撇出''，并用诸如计算机的电子基质之类的方法实现。至于对大脑的复制需要做到什么深度，现在还完全不淸楚

人工智能研究能否与对大脑的模拟平行发展？

有人认为，要想实现人工智能，就总有一天得要模拟或复制大脑的实际硬件。这种想法至少迄今为止是使很多人工智能工作者极为反感的。人们仍在纳闷：``我们得把大脑模拟到什么精度才算实现了人工智能？''实际答案可能是：这完全依赖于你想要模拟人类意识的多少特征。

玩好跳棋的能力是否足够成为智能的指标？如果是，那么人工智能就已经存在了，因为下跳棋的程序是世界第一流水平。或者智能是不是一种像大学一年级微积分课程中用纸笔求函数积分的能力？如果是，那么人工智能就已经存在了，因为符号积分运算的程序在大多数情形已胜过最熟练的人。或者，智能是不是弈棋之类的能力？如果是，那么人工智能就已经上了正路，因为下国际象棋的程序已能战胜最好的业余棋手了，而且机器棋手的水平仍然可以继续改善。

历史上，人们对于什么性质在机械化了之后能算是无可争议地构成了智能，一直是很幼稚的。有些时候，当我们朝着人工智能方向前进了一步之后，却仿佛不是造出了某种大家都承认的确是智能的东西，而只是弄清了实际智能不是哪一种东西。如果智能包括学习、创造、情感响应、美的感受力、自我意识，那前面的路就还长，而且可能一直要到我们完全复制了一个活的大脑，才算是实现了这些。

美感、螃蟹和灵魂

上面所说的这一切，与螃蟹在阿基里斯面前的那番高超表演有什么联系呢？这里有两个缠在一起的问题，它们是：

⑴．在任何环境下，任何大脑过程是否能够不违背丘奇-图灵论题，而完全可靠地区分真的和假的TNT陈述------这是不是一件原则上做不到的事？

⑵．美感是不是大脑过程？

先来答复（1）。如果允许背离丘奇-图灵论题，那么，对话中的古怪事情就似乎不会遇到什么本质性的障碍了。所以我们感兴趣的是，一个相信丘奇-图灵论题的人是不是必然不相信螃蟹的能力。这完全依赖于究竟相信丘奇-图灵论题的哪种形式。比方说，如果你只在大众过程形式下签名，你只需说螃蟹的能力是不可传播的，就很容易使两者和睦共处了。反之，如果你相信简化论形式，你就很难相信螃蟹所诡称的能力（根据\heititext{丘奇定理}------我们马上就要证明它）。相信各种中间形式将使你在争论中模棱两可。不过，如果为了方便而常常改变自己的立场，那就纯属见风使舵了。

好像我们还应当再给出丘奇-图灵论题的一种新的形式，一种为相当大量的人所心领神会的形式，一种为很多著作者以各种方式公开使用的形式。这批作者中，较为著名的有：哲学家休伯特·德雷福斯、贾基、摩蒂莫·陶布以及卢卡斯，生物学家兼哲学家迈克尔·波拉尼（最出色的整体论者），杰出的澳大利亚神经生理学家约翰·埃克勒斯。我确信还有很多其他人的著作也表达过类似的思想，也有不计其数的读者是赞同的。下面我来尽量概括出他们的立场中共同的部分。我做得可能并不完全正确，但我将力图尽可能准确地保持其原有的韵味：

\heititext{丘奇-图灵论题，唯灵论形式}：大脑所能做的某些种类的事情可以大致地由一台计算机来模拟，不过不是大多数事情，而是些不那么吸引人的事情。不管怎么说，即使都能模拟，灵魂仍将留待解释，而且没有什么方法能让计算机来承担这个任务。

丘奇-图灵论题的这种形式与《赞美螃蟹》中的故事有两种关联。第一种情形是，信奉它的人可能会把那个故事看成不可相信的傻话，不过原则上并不禁止它。第二种情形，他们可能会断言，对于诸如美感这类品德的评价，是与难以捉摸的心灵相联的一种特性，因而生来就只能为人所具有，单单机器是不行的。

我们过些时候再回到这第二点，但首先，既然我们在讨论``唯灵论''，我们得先把上面的丘奇-图灵论题表述成一种更为极端的形式，因为近年来大量受过良好教育的人都认可这种形式。

\heititext{丘奇-图灵论题，反科学形式}：计算机是荒唐的。一般说来科学也都如此。

在那些认为带有数量和严格性味道的东西是对人类价值的一种威胁的人们中间，这种观点颇为流行。真是令人遗憾，他们意识不到深究诸如人类心智等抽象结构时所包含的深刻性、复杂性和美，正是在那里，需要密切联系``人是什么''这样一个根本问题。

再回到美感问题上来。我们要考虑一下美感是否一种大脑过程。如果是，它是否可以用计算机来模拟。那些相信不能用大脑解释美感的人自然很难相信计算机会有美感。相信美感是大脑过程的人，仍是按照他们相信丘奇-图灵论题的哪种形式而取不同的态度。整个简化论学派都相信：任何一个大脑过程都能在原则上转换成计算机程序，然而其他人可能会感到美是一个定义过于不清楚的概念，无法使计算机程序吸收。或许他们觉得对美的鉴赏要求一种非理性的因素，因而与计算机的条理性不相容。

非理性的东西与理性的东西可以共存于不同的层次

不过，``非理性与计算机不相容''这个观念，是由严重的层次混淆引起的。这个错误观念出自一种想法：由于计算机是运行得完美无缺的机器，所以它们在所有的层次上都必然得``逻辑地''行动。而十分明显的是，一台计算机可以按指令打印出一系列不合逻辑的语句------或者，换个方式说，打印出一批具有随机真值的语句。然而，按照这样一类指令工作，计算机并不会出任何错误！相反，仅当打印出了指令规定的语句之外的东西时，它才算是出错了。这说明，在一个层次上完美无缺的运行如何会支持较高层次上的符号操作------而高层上的目标却可能与真理的传播毫无关系。

看透这一点的另一种方法，是记住大脑也是由一些运行得完美无缺的元素------神经元------所组成的集合体。一旦全部新来的信号跨入某个神经元的门槛，乓！------它就发射了。决不会有一个神经元忘记它的算术知识------把它的输入马马虎虎地加起来得到一个错误答案。即使一个神经元死了，它仍能在下述意义上正确行动：它的组成元素继续遵从数学和物理的规律。而且，众所周知，神经元完全有可能在它自己的层次上以极其出人意料的方式支持错误的高层行为。图109就是要说明层次之间的这种冲突：在心智的软件中有一个错误的信念，它却得到具有完美无缺的行为的大脑硬件的支持。

\pandocbounded{\includegraphics[keepaspectratio]{./media/OEBPS/Images/img_109.png}}

图109．大脑的低层与高层之间的冲突

大脑是理性的，而心智则可能不是{[}作者绘{]}

这种观点------在先前讨论各种各样的问题时曾多次提到------其实就是说：意义可以存在于符号操作系统的两个或多个不同的层次上，而且与意义一道，正确性与错误性能在所有这些层次上存在。意义在一个给定层次上的出现，取决于现实世界是否是以一种同构（还是较为松散的）的形式在这个层次上得到反映。所以，神经元总是正确地执行加法（其实，是比加法复杂得多的计算），这一事实完全不能担保它们的机制所支持的最高层结论是正确的。即便大脑的最高层忙于证明布尔佛教的公案，或者忙于冥想禅宗代数的定理，它的神经元却仍在合理地行动着。同理，大脑中感受美的经验的高层符号过程，其底层是完全理性的，那里发生着完美无缺的行动。任何非理性的东西，如果存在，就是在较高的层次上，而且是低层事件的旁效现象一个后果而已。

为了从另一个角度来看这个问题，我们来假定你正难以做出决定，不知是该要一盘茄汁鳜鱼还是黄连汁鳜鱼。这是不是就蕴涵着你的神经元细胞也受到阻碍，而难以决定它们是否发射呢？当然不是。你的鳜鱼混乱是一个高层王国的事件，以很有组织的方式依赖于数以千计的神经元的总体的有效发射。这多少有点讽刺，然而只要想一想，就会觉得是完全明白的。尽管如此，断定有关心智和计算机的几乎全部的混乱都来源于这种基本的层次混乱，大概是十分公平合理的。

没有什么理由使人相信，一台计算机中运转得完美无缺的硬件不可能支持那些体现诸如混乱、遗忘以及美感这类复杂事件的高层符号行为。这将需要大量的子系统按照一种复杂的``逻辑''相互作用。表面的行为可能表现为理性或非理性，然而暗地里则应该是可靠的逻辑硬件的工作。

再驳卢卡斯

这种层次区别倒顺便给我们提供了与卢卡斯辩论的某种新武器。卢卡斯的论证建立在``哥德尔定理按其定义可用于机器''这样一种思想之上。事实上，卢卡斯极力宣称：

哥德尔定理必定适用于控制论机器，因为作为机器的本质特征，机器该是形式系统的一个具体特例。\footnote{}

如我们已经看到的，在硬件层次上这是真的------然而由于或许还存在更高的层次，所以这不是关于这个论题的最后断语。卢卡斯强调说，在他所讨论的模拟意识的机器里，只有一个层次上发生了符号操作。例如分离规则（在他的文章中叫``假言推理''）就要做成硬件，并将成为这样一台机器的一个不可改变的特性。他更进一步地宣称，如果假言推理不是该机器系统的永远不变的栋梁，而是可以偶然撇开的，那么：

该系统将不再是一个形式的逻辑系统，而且该机器将胜任心智的模型这个角色。\footnote{}

在人工智能研究中正在开发的很多程序与那些生成数论真理的程序------一些具有不变的推导规则和固定的公理集合的程序------没有多少共同之处，而且它们当然是有意做成``心智模型''的。在它们的顶层------``非形式''层次------可能有一些意象的操作、类比的形式、念头的遗忘、概念的混淆以及差别的抹杀等等。不过，这并不与下述事实相矛盾：它们就像大脑依赖其神经元的正确行动一样，依赖于支持它们的硬件的正确行动。所以，人工智能程序依然是``形式系统的具体实例''------不过他们不是能应用卢卡斯所曲解了的哥德尔证明的那种机器。卢卡斯的论证只适用于它们的底层，可是它们的智能------不论高低------并不依赖那个底层。

我们还可以从另一个角度考察卢卡斯对``心智过程必须怎样在计算机程序内部表示''这个问题的那个过于简单的看法。在讨论一致性问题时，他写到：

如果我们实际是些不一致的机器，我们就应满足于这种不一致性，从而痛痛快快地肯定一对矛盾的两个方面。更进一步，我们该是可以随便讲任何事情------但我们没有这样。容易说明，在一个不一致的形式系统内，一切都是可证的。\footnote{}

最后这句话表明，卢卡斯假定了命题演算一定得包含在任何一个能做推理的形式系统之中。具体地说，他是在考虑命题演算的定理＜＜P∧\textasciitilde P＞→Q＞。显然，他错误地认为这是机械化推理的一条必不可少的要素。无论如何，要是说像命题推理这样的逻辑思维过程是一个人工智能程序的一般智力的结果，而不是预先程序化了的东西，是完全有可能的。这恰是在人类身上发生的事情！没有什么特别的理由该假定狭义命题演算------连同它们所必需的刻板规则，以及有关一致性的呆笨的定义------都得要出现在这样一个程序之中。

人工智能的一个基点

我们可以把这一通闲扯都归结到层次的区分。在离开这个话题之前，我们给出丘奇-图灵论题的一个最强的------也是最后的------形式：

\heititext{丘奇-图灵论题，人工智能形式}：任何种类的心智过程都可以用一个计算机程序来模拟，而该程序的基础语言与FlooP一样强，也就是说全体部分递归函数都能用这种语言程序化。

还需要指出的是，实际上很多人工智能研究者信奉的是与丘奇-图灵论题密切相关的另一个信条，我把它叫做人工智能论题，叙述如下：

人工智能论题：随着智能机的发展，它的基础机制会逐渐收敛于人类智能的基础机制。

换句话说，一切智能都只是同一主题的各种变奏。为了创造真正的智能，人工智能工作者如果想要使他们的机器达到我们所具有的能力，他们就得坚持深入那些较低的层次，使之越来越接近大脑的机制。

丘奇定理

现在，让我们回过头来看螃蟹，再看看他的定理资格判定过程（伪装成一种判定音乐美的过滤器）是否与现实世界相容这一问题。实际上，根据对话中提到的事实，我们没有办法推定螃蟹的才智是不是一种分清定理和非定理的能力，或者换个说法，分清真陈述和假陈述的能力。当然，在很多情形下这是一回事，不过\heititext{哥德尔定理}表明这并不总是一回事。但是，这没有关系：如果我们相信丘奇-图灵论题的人工智能形式，那么这两种说法都不可能是真的。在任何强如TNT的形式系统中都不可能有判定定理的过程------这个命题叫\heititext{丘奇定理}。没有一个判定数论真理的过程------如果这种真理存在的话。读者在看过TNT的全部分叉现象之后，很可能会怀疑这种存在性------这样一个命题可以很快地从\heititext{塔斯基定理}（发表于1933年。实际上其思想在此之前就早已为塔斯基所知）推出。

这两个极为重要的元数学结果的证明是很相似的。它们都可以很快地从一个自指结构中得到。我们先来考虑TNT定理的判定过程。如果存在一个统一的方法使人们可以判定一个给定的公式X究竟应归入``定理''类还是``非定理''类，那么，根据丘奇-图灵论题（标准形式），就该有一个有终止的FlooP程序（一般递归函数），当给出公式X的哥德尔数作为输入时，该程序能给出与上述方法相同的判定。其关键步骤在于（读者应该还记得）：能用一个有终止的FlooP程序检测的任何一种性质在TNT中都有一个体现。这意味着``是TNT定理''这条性质能在TNT内部体现（与仅仅表示是不同的）。然而，如我们马上就能看到的那样，这会使我们陷入困境，因为如果``是定理''是一个可体现的属性，那么哥德尔公式G就会是与说谎者悖论同样可怕了。

这全是因为G在说：``G不是TNT定理''。假如G是定理，那么由于``是定理''按假定应可体现，所以断言``G是定理''的那个TNT公式应当是TNT定理。而这个公式就是\textasciitilde G，即G的否定。所以TNT就不一致了。另一方面，假如G不是定理，那么根据假设的定理资格的可体现性，断言``G不是定理''的那个公式就应是TNT定理。然而这个公式就是G，所以我们仍得到悖论。与前面的情形不同，这里根本无法消除这个悖论。问题之所以出现，就是由假定``是定理''这一性质可由某个TNT公式体现。因而，我们就必须追溯回去，摒弃这个假设。这就迫使我们得出：没有一个FlooP程序能分清定理的哥德尔数和非定理的哥德尔数。最后，如果我们接受丘奇-图灵论题的人工智能形式，那我们就得进一步地回溯，得出：无论如何都不会有什么方法能使人类可靠地分清定理和非定理------这也包括基于美感的判断。那些只同意丘奇-图灵论题的大众过程形式的人，可能仍然认为螃蟹的成就是可能的。不过，在丘奇-图灵论题的所有各种形式当中，这一个也许是最难得到支持的。

塔斯基定理

现在我们接着看塔斯基的结果。塔斯基问，是否可以有一种在TNT中体现数论真理概念的方法。我们已然看到``是定理''这个性质是可表示的（尽管不是可体现的）。塔斯基则对于有关真理概念的类似问题产生了兴趣。尤其是，他希望确定是否会有一个只含一个自由变元a的TNT公式，它可以翻译成：

``哥德尔数为a的公式表示一条真理''。

像塔斯基那样，我们先来假定有这样一个公式------把它缩写成TRUE\{a\}。然后，所要做的事情就是使用对角线方法制造一个句子，它对自己下断语，说自己不是真的。我们完全照搬哥德尔的方法，从一个``服''号串入手：

\gebfont{∃}a:＜\textasciitilde TRUE\{a\}∧ARITHMOQUINE\{a", a\}＞

把这个``服''号串的哥德尔数记作t，然后就算术\gebfont{㧟}摁它，从而得到塔斯基公式T：

\begin{center}

\pandocbounded{\includegraphics[keepaspectratio]{./media/OEBPS/Images/Formula13.png}}

\end{center}

如果解释一下，它是在说：

``t的算术\gebfont{㧟}摁化是个假陈述的哥德尔数。''

可是由于t的算术\gebfont{㧟}摁化是T自身的哥德尔数，所以塔斯基公式T恰好在TNT中重现了说谎者悖论，它在说自己：``我是假话''。当然，这就导致了结论：它既真又假（或既不真又不假）。于是又产生了一个令人感兴趣的问题：说谎者悖论的出现有什么不好？这有什么要紧的吗？毕竟汉语中已经有了这个悖论，而中国语言并没有因此就灰飞烟灭。

那个被赞美的螃蟹是不可能存在的

其答案就在于：别忘了这里牵涉到了意义的两个层次。一个层次是我们刚才用过的，而另一个层次则是作为一个数论陈述。如果塔斯基公式T真的存在，那它就是一个同时既真又假的论及自然数的陈述！难就难在这里，尽管对于汉语来说，说谎者悖论总能因其话题（它自己的真理性）过于抽象而清除掉，可这一个却不然，因为它变成了一个严格的数论陈述！如果我们相信这是一种荒唐可笑的事态，那就只好丢掉TRUE\{a\}是存在的这一假设。就是说没有一种方法能在TNT内部表示真理的概念。应当看到，这个结果就使得``是真理''成为比``是定理''更难以捉摸的性质，因为``是定理''还是可表示的。按照前面（涉及丘奇-图灵论题的人工智能形式）谈过的回溯推理，我们得到结论：

``螃蟹的心智辨认起真理来，一点也不比辨认TNT定理更强。''

因为前一种辨认违背\heititext{塔斯基-丘奇-图灵定理}（``不存在确认算术真理的判定过程''），而后者则违背了\heititext{丘奇定理}。

形式的两种类型

当把``形式''这个词用于各式各样形状复杂的结构时，想想它的意义是很有意思的。例如，当我们观赏一幅画并感到美时，引起我们反应的是什么？是视网膜上的点和线的``形式''吗？显然一定是的，因为它就是如此进入我们头脑的分析机制的------不过这一处理过程的复杂性使我们感到并不只是看到了一个二维的表面，我们在对这幅画的某种内在意义------某种莫名其妙地陷入了二维的那个多维对象------产生共鸣。这里重要的是``意义''这个词。我们的心智有一些接受二维模型的解释机制，然后再从中``抽''出高维概念，这些高维概念复杂得我们无法有意识地对之进行描述。其实，我们对音乐的反应也是一样。

主观感觉是，抽出内在意义的机制完全不同于检验有还是没有某种特定性质（比如一个符号串是否良构）的判定过程。也许这是因为内在意义过一段时间以后会泄露出更多的东西。我们永远不可能像对良构性那样有把握地确定什么时候才算大功告成了。

这提示了一种区分，在我们所分析的模式中，对``形式''的两种含义可以作个区分。首先，有一些性质，如良构性，它们可以用预知有终止的检验来检查，比如用BlooP程序。我把这些性质叫做形式的句法性质。直观上可以察觉到，形式的句法方面比较接近于表面，因而并不能激发多维认知结构的创造。

作为对照，形式的语义方面是不能在预定时间期限内检验的：它仍需要无尽头的检验。如我们已经看到的，检验TNT符号串是不是定理这种情形就是如此。你无法把某个标准的检验法用于一个符号串，以辨认出它是否是一条定理。不知什么原因，``涉及到了意义''这一事实，与弄清一个符号串是否是TNT定理的难度，是密切相关的。从本质上说，抽出一个符号串的意义的行为，涉及到建立它与其它所有符号串之间相关联时的全部意味，而这当然就导致一个无尽头的征途。所以，语义性质就与无终止的搜索联系在一起了。因为，在一个很重要的意义上，一个客体的意义并不局限于该客体自身之内。这自然并不是要说在时间终结之前就不可能理解任何客体的意义。因为，随着时间的流逝，越来越多的意义变明朗了。但是，不论过多久，也不可能发现其意义的所有方面。

意义来自认知结构间的联系

让我们从符号串转向乐曲，变个花样看看。如果你愿意的话，仍可把每一处提到乐曲的地方都换成``符号串''。这种讨论必然是一般性的，只不过我觉得还是谈音乐更容易使人理解其韵味。乐曲的意义有一种少见的两重性：一方面，凭借着它与世界上其它很多事物的关系，似乎要向四周蔓延------而另一方面，一段乐曲的意义显然源自音乐本身，所以它必定能定位于音乐内部的某个地方。

考虑解释机制可以摆脱这种窘境。所谓解释机制就是抽出意义的那种机制（这里说的解释机制并不是指乐曲的演奏者，而是指欣赏该曲时从中得出意义的那位听众的心理机制）。当第一次听到某首乐曲时，解释机制会发现它的意义的很多重要方面。这似乎证明了意义寓于乐曲自身之内、并可从中直接读出这样一种观念，不过这只是问题的一部分。音乐解释机制是通过建立一个多维的认知结构------该乐曲的一个心智表示------而工作的。这个多维认知结构努力寻找与那些为以前的经验编码的多维心智结构之间的联系，借此同早就存在的信息结合起来。随着这个过程的进行，整个意义就逐渐显露出来了。事实上，当一个人感到他已经悟透一段乐曲的核心意义时，可能若干年的时间已经过去了。这似乎又有利于相反的观点：音乐的意义向四方延伸，解释机制的作用就是逐步收集它。

真理毋庸置疑是位于这两者之间的某个地方：无论音乐的意义还是语言的意义，在某种程度上都是可定位的，在某种程度上又都是向四方延伸的。用第六章的话来说，我们可以说乐曲和文章都既是显明意义的触发器，又是运载者。可以通过一块刻有铭文的古碑的例子，来给意义的这种两重论作个生动的说明：碑文的意义部分地贮存在世界各地的图书馆和学者们的大脑之内，但它显然也隐含在古碑自身之内。

于是，要刻划``句法''性质和``语义''性质（在刚刚提到那种意义下）之间的差别就还有一条路子，即：句法性质毫不含糊地存在于所考虑的客体之内，而语义性质则依赖于它与潜在地无穷多的客体类之间的关系，因而不是完全可以定位的。原则上说，句法性质中没有什么秘密或隐藏的东西，然而这种隐蔽性却是语义性质的本质。这就是我区别视觉形式的``句法''方面和``语义''方面的理由。

美、真与形式

美又怎么样呢？按照上面的思想，它当然不是句法性质。那么它能是语义性质吗？比方说，美是某幅具体的画所具有的性质吗？让我们干脆就限于考虑一个观察者。每个人都有过这样的经验：这个时候觉得某个东西很美，那个时候又会觉得它很丑------而其它时候又可能觉得它很平常。那么，美是不是一个随时间变化的属性呢？但是倒过来也可以说，随时间变化的是那位观察者。那么，当确定了一个在特定的时间观赏一幅特定的画的特定的人之后，是否就有理由断定美是一种确确实实的、或者存在或者不存在的性质呢？还是说，有关美的问题仍然还有某些定义不清的、不可捉摸的东西？

随着环境的变化，在每个人心里都可能会引发解释器的不同层次。这些形形色色的解释器抽出不同的意义，建立不同的联系，并一般地对所有深入的方面给出不同的评估。所以，看起来，美这个概念极难把握。正因为如此，我才在《的确该赞美螃蟹》中把美与真联系在一起，我们已经看到，``真''乃是整个元数学中最不好捉摸的概念之一。

说谎者悖论的神经基质

我打算以关于真理的那个中心问题------即说谎者悖论------的某些想法，作为本章的结尾。我认为，塔斯基在TNT内部对说谎者悖论的复制，指出了深入理解汉语中的说谎者悖论的一个途径。塔斯基所发现的东西是：该悖论的塔斯基形式有两个不同层次。在一个层次上，它是谈论自身的句子，如果它假，那就会真，如果它真，就又会假。在另一个层次------我喜欢把它叫作算术基质------上，它是一个谈论整数的句子，它为真当且仅当它为假。

这样一来，由于某种原因，这后一种说法就比前一种更为令人费解。由于前一种说法的自指性，有人就干脆贬斥为``没有意义''。但你却不能这样贬斥一个谈论整数的悖论陈述，谈论整数的陈述根本不能既真又假。

我的感觉是：塔斯基对说谎者悖论的转换指导了我们去寻找该悖论汉语形式的基质。在其算术形式中，高层的意义建立在低层的算术层次之上。也许与之类似，我们见到的那个自指句子（``本句子是假的''）只是一个双层次统一体的顶层。那么较低的层次又是什么呢？支持语言的结构又是什么呢？是大脑。所以必须给说谎者悖论找一个神经基质------一个较低的、由互相冲突的物理事件------也就是按其性质而言不能同时出现的事件------组成的层次。如果这种物理基质存在，那么我们之所以无法理清说谎者句子的来龙去脉，就是由于我们的大脑在试图作一件力所不及的事情。

那么，这些互相冲突的物理事件的实质又是什么呢？也许当你听到说谎者句子时，你的大脑就建立了该句子的一种``编码''------相互作用的符号的一种内部布局。然后，它就试图把句子归结为``真''、``假''两类。这种分类行为必定迫使某些符号按一种特殊的方式发生相互作用（也许遇到任何一个句子时都会发生这种过程）。如果碰巧这种分类行动会从物理上搞乱该句子的编码动作------一种轻易不会发生的事情------那就算碰上麻烦了，因为它相当于试图要一台唱机播放毁坏它自己的唱片。我们用物理术语而不是用神经学的术语描述了这种冲突。如果这种分析到现在为止还是有效的，那么当我们了解了一些大脑中来自神经元及其发射的``符号''建构，又了解了一些给句子编码的方式时，余下的讨论就可以继续下去了。

说谎者悖论的神经基质的这幅草图（至少对我）暗示了：说谎者悖论的汉语形式的这种解决，很可能类似于该悖论的塔斯基形式的解决。这种解决包括放弃``大脑总是能够为真理概念提供一个完全精确的描写''这样一种观念。这种解决的新颖性，在于它暗示着给出真理的全面模型是不可能的，而这是由于一个来自物理方面的理由，即：做这样一个模型从物理上需要在一个大脑中出现不相容的两个事件。

\begin{enumerate}
\item
  \phantomsection\label{Chapter17.xhtml_B_1}
\item
  \phantomsection\label{Chapter17.xhtml_B_2}
\item
  \phantomsection\label{Chapter17.xhtml_B_3}
\item
  \phantomsection\label{Chapter17.xhtml_B_4}
\item
  \phantomsection\label{Chapter17.xhtml_B_5}
\item
  \phantomsection\label{Chapter17.xhtml_B_6}
\item
  \phantomsection\label{Chapter17.xhtml_B_7}
\item
  \phantomsection\label{Chapter17.xhtml_B_8}
\item
  \phantomsection\label{Chapter17.xhtml_B_9}
\item
  \phantomsection\label{Chapter17.xhtml_B_10}
\item
  \phantomsection\label{Chapter17.xhtml_B_11}
\end{enumerate}

\phantomsection\label{Dialog18.xhtml}{}

\chapter[施德鲁，人设计的玩具]{\texorpdfstring{施德鲁，人设计的玩具\footnote{}}{施德鲁，人设计的玩具}}

有一天，伊她·娥英女士漫步到麻省理工学院人工智能实验室，碰上了才气横溢的青年计算机程序施德鲁。施德鲁正急着找人来试验试验新近开发出来的一个人------``格拉德·维维诺诺博士''。施德鲁向伊她说明：这位格拉德·维维诺诺博士在他的专业领域里是颇为聪明的，他的专业是分析关于一种``玩具世界''的对话，而这种``玩具世界''由各种不同形状、不同大小、不同颜色的积木组成。这些积木摆在一张桌子上，可以捡起来，挪到别的位置。伊她女士一听就着了迷，立刻在施德鲁的鍵盘上噼哩啪啦地打起字来。格拉德·维维诺诺博士答应站在她身后随时提供一些说明。

\pandocbounded{\includegraphics[keepaspectratio]{./media/OEBPS/Images/img_110.png}}

图110．与施德鲁对话的情景之一

``拣起一个大红方块。''{[}选自特里·维诺格拉德，《理解自然语言》第8页{]}

01．\dialogname{伊她·娥英：}拣起一个大红方块。{\small {[}参见图110{]}}

\dialogname{施德鲁：}行。

\dialogname{格拉德·维维诺诺博士：}施德鲁执行命令时候照例回答：``行！''。为拣起这个红方块，须先找出一个地方，以便移开上面的绿方块。

02．\dialogname{伊她·娥英：}抓住那个方锥。

\dialogname{施德鲁：}我不明白你指的是哪个方锥。

\dialogname{格拉德·维维诺诺博士：}屏幕上有三个方锥，施德鲁懂得``那个''所指的是说话人心目中的一个对象。

03．\dialogname{伊她·娥英：}找一块比你手里的那块大点的积木，把它搁到盒子里。{\small {[}参见图111{]}}

\pandocbounded{\includegraphics[keepaspectratio]{./media/OEBPS/Images/img_111.png}}

图111．与施德鲁对话的情景之二

``找一块比你手里的那块大点的积木，把它搁到盒子里。''{[}选自维诺格拉德，《理解自然语言》第9页{]}

\dialogname{格拉德·维维诺诺博士：}施德鲁先要弄清``那块''、``它''各指什么，还要推定``手里''有什么。

\dialogname{施德鲁：}``它''指大点的那块积木。

\dialogname{格拉德·维维诺诺博士：}施德鲁能用启发规则确定``它''所指的是``手里的那块''还是``找''出来的那块。

\dialogname{施德鲁：}行。

\dialogname{格拉德·维维诺诺博士：}施德鲁须先放下红方块，才能拣起蓝方块。

04．\dialogname{伊她·娥英：}盒子里都有什么？

\dialogname{施德鲁：}蓝方锥和蓝方块。

05．\dialogname{伊她·娥英：}方锥下边是什么？

\dialogname{格拉德·维维诺诺博士：}以前说``方锥''时，施德鲁说``我不明白你说的是哪个方锥''。现在，他认为你所说的就是刚才谈及的那个方锥。

\dialogname{施德鲁：}盒子底。

06．\dialogname{伊她·娥英：}有几个方块没在盒子里？

\dialogname{格拉德·维维诺诺博士：}施德鲁能掌握各种不同的否定词语的含义，例如``没''，``不''等等。

\dialogname{施德鲁：}四个。

07．\dialogname{伊她·娥英：}这四个方块里有没有比我刚说让你拣起的那块大的？

\dialogname{格拉德·维维诺诺博士：}``这四个方块''指上一问题中谈及的方块；``比\ldots\ldots 大的''也应指方块。名词短语中可能含有复杂的数量词，如``至少一块''、``三块以上''、``恰好七块''等。为确定``我刚才让你拣起的那块''，需要用到已发生事件的记录。此外，``我刚才说让你拣起的''和``那块''的句法关系是间接的。施德鲁的句法分析程序能以一种具有普遍性的方式处置这种复杂的句法关系。

\dialogname{施德鲁：}有，就是那个红方块。

08．\dialogname{伊她·娥英：}它下边有东西吗？

\dialogname{施德鲁：}有，桌面。

\dialogname{格拉德·维维诺诺博士：}施德鲁用启发规则以确定对于人有关的信息，即使问题中并未明确提出要求。

09．\dialogname{伊她·娥英：}桌子能不能拣木块？

\dialogname{施德鲁：}不能。

\dialogname{格拉德·维维诺诺博士：}只有动物才能拣东西，而桌子不是动物。施德鲁的语义规则具备这些知识。

10．\dialogname{伊她·娥英：}方锥能不能搁在方块上？

\dialogname{施德鲁：}能。

\dialogname{格拉德·维维诺诺博士：}施德鲁的推理系统发现了一个例证，所以他知道这是可能的。

11．\dialogname{伊她·娥英：}方锥上能不能搁方块？

\dialogname{施德鲁：}不知道。

\dialogname{格拉德·维维诺诺博士：}这个问题本身是有意义的，但是在屏幕上并没有出现例证。施德鲁无法看到他自己的程序里的细节，尽管他的能力终究是由这些细节所确定的。

12．\dialogname{伊她·娥英：}把两个方锥摞起来。

\dialogname{施德鲁：}我摞不起来。

\dialogname{格拉德·维维诺诺博士：}施德鲁尝试这样做，但失败了。

13．\dialogname{伊她·娥英：}蓝方锥是松木的。

\dialogname{施德鲁：}明白了。

\dialogname{格拉德·维维诺诺博士：}施德鲁知道对象的颜色、形状，大小和位置。质料则是对象的未知性质。告诉它以后，它就记录在数据库中。

14．\dialogname{伊她·娥英：}除了红的木块和方锥下边的木块以外，都是松木的。

\dialogname{施德鲁：}明白了。

\dialogname{格拉德·维维诺诺博士：}施德鲁对这句话的理解，导致它建立一个新的推理过程，以确定某个木块是否是松木的。

15．\dialogname{伊她·娥英：}盒子是不是松木的？

\dialogname{施德鲁：}不是。

\dialogname{格拉德·维维诺诺博士：}盒子中有个方锥，盒子是方锥下面的东西。

16．\dialogname{伊她·娥英：}盒子里有松木的东西吗？

\dialogname{施德鲁：}有两个：蓝方块和蓝方锥。

17．\dialogname{伊她·娥英：}你把两个红方块跟一个绿色的方块或是方锥摞成一摞，好吗？{\small {[}参见图112{]}}

\pandocbounded{\includegraphics[keepaspectratio]{./media/OEBPS/Images/img_112.png}}

图112．与施德鲁对话的情景之三

``你把两个红方块跟一个绿色的方块或是方锥摞成一摞，好吗？''{[}选自维诺格拉德，《理解自然语言》第12页{]}

\dialogname{施德鲁：}行。

\dialogname{格拉德·维维诺诺博士：}对``跟''、``或是''这类逻辑联词的处理同时考虑到了语法和语义。一个命令即使形式上是一个问题也会被执行。这个命令执行过程如下：把一个绿方块放在大红方块上（他选中的是上面没有东西的绿方块），把小红方块上面的小方锥移开，再把小红方块摞到绿方块上面。

18．\dialogname{伊她·娥英：}哪个方块是放在桌面上的？

\dialogname{施德鲁：}红方锥底下的大绿方块。

\dialogname{格拉德·维维诺诺博士：}一般说来回答``大绿方块''已经满意了，但就这个问题而言，最好把它限定得清楚一些。

19．\dialogname{伊她·娥英：}有没有放在方锥后边的大块积木？

\dialogname{施德鲁：}有三块：一个大红方块，一个大绿方块，一个大蓝方块。

20．\dialogname{伊她·娥英：}往那个上面放着方锥的绿方块上再搁上一个小方块。

\dialogname{施德鲁：}行。

21．\dialogname{伊她·娥英：}把最小的方锥搁到它上面。

\dialogname{格拉德·维维诺诺博士：}施德鲁先要比较所有方锥的大小，再采取行动。

\dialogname{施德鲁：}行。

22．\dialogname{伊她·娥英：}有多少东西是放在绿方块上面的？

\dialogname{施德鲁：}你说``放在绿方块上面''，我不明白这个``上面''指什么意思。你是指：

1．与它的表面接触，还是

2．在它上方的任何地方？

\dialogname{格拉德·维维诺诺博士：}前面那一句话里，``上面''的解释未发生困难，是因为从上下文中可以澄清它的含义。在这一句话里，解释歧义的启发规则无法做出清楚的定义。施德鲁只好向你要求进一步的信息。类似的问题在解释``它''、``那些''等词的时候也可能遇到。

23．\dialogname{伊她·娥英：}你把绿方锥放在小方块上面以前动过哪个方锥没有？

\dialogname{施德鲁：}动过绿的。

\dialogname{格拉德·维维诺诺博士：}施德鲁能从它的记忆中找出现在说的是什么事件，并能把它和别的事件前后联系起来。对于本句的分析涉及像``没有''这样的词的种种不同用法。

24．\dialogname{伊她·娥英：}你是什么时候动的它？

\dialogname{施德鲁：}我把两个红方块跟一个绿方块摞起来的时候。

\dialogname{格拉德·维维诺诺博士：}时间可以用当时所作的工作来描述。

25．\dialogname{伊她·娥英：}为什么要动它？

\dialogname{格拉德·维维诺诺博士：}施德鲁记得前一句对话，因此能理解这类问题。

\dialogname{施德鲁：}把它拿到一边去。

\dialogname{格拉德·维维诺诺博士：}施德鲁能从原有的子目标树找出其工作的每一步骤的目的。

26．\dialogname{伊她·娥英：}这又是为什么？

\dialogname{施德鲁：}好把它下边的红方块抓起来。

\dialogname{格拉德·维维诺诺博士：}``这''指的是前一答句中的动作。

27．\dialogname{伊她·娥英：}为什么要抓那个红方块？

\dialogname{格拉德·维维诺诺博士：}``那个''指的是前一句里所说的红方块。

\dialogname{施德鲁：}好把它放在一个大绿方块上。

28．\dialogname{伊她·娥英：}这又是为什么？

\dialogname{施德鲁：}好把两个红方块和一个大绿方块摞起来。

29．\dialogname{伊她·娥英：}这又是为什么？

\dialogname{施德鲁：}你让我这样做的。

\dialogname{格拉德·维维诺诺博士：}施德鲁记得他做了什么，而不是对它的要求的原话。

30．\dialogname{伊她·娥英：}这件事你是怎样做的？

\dialogname{格拉德·维维诺诺博士：}``这件事''指的是前一问句里所说的事。

\dialogname{施德鲁：}把一个大红方块放在桌上，松手，再把一个红方块放在绿方块上，松手。

\dialogname{格拉德·维维诺诺博士：}本句描述了为实现一个目标所建立的最高层子目标。

31．\dialogname{伊她·娥英：}你干这些事的时候动了多少木块？

\dialogname{格拉德·维维诺诺博士：}``干这些事的时候''指的是一段时间。

\dialogname{施德鲁：}四块。

32．\dialogname{伊她·娥英：}你抓起小红方块以前，它上面有什么？

\dialogname{施德鲁：}绿方锥。

\dialogname{格拉德·维维诺诺博士：}回顾以往的动作可以推算过去的场面。

33．\dialogname{伊她·娥英：}那时候盒子左边有五块积木。

\dialogname{格拉德·维维诺诺博士：}``那时候''是指上一句所说的时间。

\dialogname{施德鲁：}没有，只有四块：两个红方块跟两个绿方块。

\dialogname{格拉德·维维诺诺博士：}前边说到质料，施德鲁马上接受了。但他知道他具有关于木块位置的全部信息，所以他要与它的知识进行核对，并表示同意与否。

34．\dialogname{伊她·娥英：}把蓝方锥搁在盒子里的蓝方块上头有问题吗？

\dialogname{格拉德·维维诺诺博士：}这句话句法上是有歧义的，因为它可以分析为``把蓝方锥搁在盒子里''的``蓝方块''上头有某种东西与否。施德鲁会排除这种句法歧义。

\dialogname{施德鲁：}没问题。

35．\dialogname{伊她·娥英：}搁上去吧。

\dialogname{施德鲁：}行。

36．\dialogname{伊她·娥英：}一个塔\ldots\ldots{}

\dialogname{施德鲁：}对不起，我不明白什么叫``塔''。

\dialogname{格拉德·维维诺诺博士：}施德鲁读到一个不知道的字，就会打断你的话。

37．\dialogname{伊她·娥英：}``塔''是用两个绿方块和一个方锥做成的。

\dialogname{格拉德·维维诺诺博士：}你用了引号，施德鲁知道你在解释一个生词，于是就用适当的形式把这个词放在词汇表里。

\dialogname{施德鲁：}明白了。

38．\dialogname{伊她·娥英：}现在有没有塔？

\dialogname{施德鲁：}没有。

\dialogname{格拉德·维维诺诺博士：}新定义可以用于辨认塔。

39．\dialogname{伊她·娥英：}你做一个吧。

\dialogname{施德鲁：}行。

\dialogname{格拉德·维维诺诺博士：}新定义也可以用于判定建造的计划。这里``一个''指的就是``塔''。于是施德鲁就可以出如下的计划：把两个大绿方块摞起来，并把红方锥放在上面。

40．\dialogname{伊她·娥英：}把最大的那个木块叫做``大胖子''。

\dialogname{施德鲁：}行。

\dialogname{格拉德·维维诺诺博士：}可以给木块取名字。

41．\dialogname{伊她·娥英：}你动过大胖子吗？

\dialogname{施德鲁：}动过。

\dialogname{格拉德·维维诺诺博士：}``动过''当然是指从对话开始以来。

42．\dialogname{伊她·娥英：}你怎么把它丢了？

\dialogname{施德鲁：}我想你说的``它''指的是大胖子。但是我不能解释一件没有发生过的事情。

\dialogname{格拉德·维维诺诺博士：}施德鲁从它的记录中查不到丢木块的事件。它有专门的程序对付这类``别有用心''的问题。

43．\dialogname{伊她·娥英：}红方锥的右边有什么？

\dialogname{施德鲁：}有四样东西：盒子，小红方块，蓝方锥和叫做大胖子的那块积木。

\dialogname{格拉德·维维诺诺博士：}所取的名字也可以用于指称事物。

44．\dialogname{伊她·娥英：}谢谢你！

\dialogname{施德鲁：}不客气！

\dialogname{格拉德·维维诺诺博士：}施德鲁的语法知识实际上是一个程序，这个程序使他能处理惯用语。

\begin{enumerate}
\item
  \phantomsection\label{Dialog18.xhtml_B_1}
\end{enumerate}

\phantomsection\label{Chapter18.xhtml}{}

第十八章

\chapter{人工智能：回顾}

图灵

在1950年，阿兰·图灵写下了一篇关于人工智能的既有远见又有争议的文章，题目是``计算机器与智能''，发表在《心智》杂志上\footnote{}。我将会谈到这篇文章，但我想先说说图灵这个人。

\pandocbounded{\includegraphics[keepaspectratio]{./media/OEBPS/Images/img_113.jpg}}

图113．阿兰·图灵

在赛跑中获胜（1950年5月）{[}摘自萨拉·图灵，《阿兰·图灵》{]}

阿兰·麦瑟森·图灵1921年生于伦敦。当他还是一个孩子的时候，就充满了好奇心和幽默感。由于他的数学天赋，他进了剑桥，在那里他关于机械和数理逻辑的兴趣相得益彰，导致他写出了一篇关于``可计算的数''的著名论文。在这篇论文中，他提出了图灵机理论，而且证明了停机问题的不可解性。这篇论文发表于1937年。在40年代，他的兴趣从计算机理论转向了计算机的实际制造。他在英国的计算机发展史上发挥了重要的作用。人工智能刚刚出现时，面临的是一片攻击与嘲笑，而图灵当时则是它的一名坚定的捍卫者。那时戴维·钱伯农是他最好的朋友之一（此人后来致力于计算机作曲）。他们俩都是国际象棋迷，并发明了``绕屋跑''下棋法：走一步棋后绕房子跑一圈------如果你跑回来时对手还没有走棋，你就有权再走一步。但他们也合作干了件正经事：发明了第一个下棋程序。他们为它起了个名字，叫``图尔钱''。图灵死得很早，才41岁------似乎是一次化学事故造成的，也有人说是自杀。他的母亲萨拉·图灵为他写了传记。从她所引用的那些人的话中，人们得到的印象是：图灵不是个循规蹈矩的人，有些不善交际，但又非常诚实正直，因此很容易被他人所伤。他喜爱游戏、下棋、孩子和骑自行车，还是个长跑好手。在剑桥上学时，他买了一把旧小提琴，自己学会了演奏。尽管水平不是很高，他却从中得到了很大乐趣。他有点古怪，对于他感兴趣的事情，他会全力以赴去做，做完便又对另外的什么东西感兴趣了。他所探索的领域之一是生物学的形态发生问题。用他母亲的话说，图灵``特别喜欢《匹克威克外传》''，但``对诗歌------除了莎士比亚的------却丝毫不感兴趣''。阿兰·图灵是计算机科学领域中真正的拓荒者之一。

图灵测验

图灵的那篇文章是这样开始的：``我准备考虑一个问题：`机器能思维吗？'\,''。正像他所指出的，这些词都太复杂了，易于挑起争端，因此我们显然应当寻找一种更好的方式来接近这个问题，即给出一种限定其含义的过程，而非一个定义。他认为这种方式就包含在他所谓的``模拟游戏''之中------今天我们称之为``图灵测验''。图灵是这样介绍的：

这种游戏需要三个人来玩：一个男人（A），一个女人（B），以及一个性别不限的询问人（C）。询问人待在一间看不到另外两人的屋子里。游戏的目的是让询问人确定另外两个人谁是男的谁是女的。他只知道这两个人的代号是X和Y，而游戏结束时他要说出``X是A，Y是B''或``X是B，Y是A''。询问人可以这样向A和B提问题：

C：请X告诉我他或她头发的长度。

现在假设X实际是A，那A就必须回答。A在游戏中的目的是设法使C作出错误的判断。这样他的回答就可能是：

``我梳的是短发，最长的一股大约有二十多厘米长''。

为了不让答话的声音给询问人以线索，答案可以写在纸上，要是用打字机就更好了。理想的办法是用电传打字机进行两个屋子之间的通讯。也可以让中间人转述问题和答案。游戏的第三个参加者（B）的目的是帮助询问人。对她来说，最好的策略大概就是提供真实的答案。她可以在答案上加上``我才是那个女的，别听他的话！''或诸如此类的东西。但这于事无补，因为那个男人也可以这样做。

现在我们要问，``如果让一台机器在这场游戏中扮演A的角色，那会发生什么情况呢？''当游戏这样进行时，询问人作出错误判断的可能性是否和参加游戏的是一男一女时一样大？这些问题实际上就是我们原先的那个``机器能思维吗？''的问题。\footnote{}

在讲清楚他的测验的性质之后，图灵接下去对此作了些评论，这些评论即使在今天看来都很复杂，更不必说在他写作的年代了。开始，他给出了询问人和被询问人之间的一小段假想的对话：\footnote{}

问：请以福斯大桥为题写一首十四行诗（这座桥位于苏格兰的福斯河口）。

答：这可把我难倒了。我从来也不会作诗。

问：把34957和70764相加。

答：（停了大约30秒然后给出答案）105621。

问：你会下国际象棋吗？

答：会。

问：我就剩一个国王在K1，没有别的子了，而你有一个国王在K6，一个车在R1，该你走棋，你怎么走？

答：（停了大约15秒），车到R8，将！

可能没有多少读者会注意到，那个算术问题中，不但反应时间异乎寻常地长，而且连答案都是错的！如果回答者是人，就很容易被解释：只不过是个计算错误而已。但如果回答者是台机器，这就可能有许多不同的解释。下面就是某些解释：

⑴．一个偶然的硬件层次上的运行时失误（即不可再现的意外故障）；

⑵．一个无意中出现的硬件（或程序）故障，它可能造成算术运算中的错误（可再现的）；

⑶．一个由该机器的程序员（或制造者）有意插入的花招，它可能不时造成算术运算中的错误，以此来迷惑询问人；

⑷．一个不可预测的旁效现象：程序由于长期进行艰苦的抽象思维，偶然出了点小错，下次就不会再错了；

⑸．一个由机器自己想出来的玩笑，故意戏弄询问人。

要搞清图灵在这里到底是什么意思，就涉及到和人工智能有关的几乎全部的重大哲学问题。

图灵接下去指出：

这个新问题的优点在于把人的物理能力和心智能力明确地区分开了\ldots\ldots 我们不想因为一台机器不能参加选美大赛而责备它，就像我们不会因为一个人没有飞机速度快而责备他一样。\footnote{}

读这篇文章时的一大乐事就是看图灵怎样展开他的每条思路。通常他总是在某个阶段摆出一种似乎是矛盾的东西，然后通过修正他的概念，再在分析的更深层次上解决。由于这种对问题入木三分的剖析，在计算机有了长足的发展，人工智能也进行了深入的工作的近三十年之后，这篇文章仍放射着耀眼的光辉。在下面的引文中你可以看到思想的这种往复工作方式：

也许有人会指责说这种游戏的设计对机器太不利了。如果那个人想装成机器，很显然他会把事情搞糟。仅在做算术题时又慢又不准确这一点上就足以让他漏馅了。难道说机器就不能完成某种应当被描述成思维、但又与人所具有的完全不同的东西吗？这种反对意见是很有力的，但至少我们可以说，如果能构造出一台能令人满意地玩这种模拟游戏的机器来，那我们就不必为这种意见所烦扰。

有人也许认为机器在玩这种``模拟游戏''时最好的策略可能并不是模拟人的行为。也许会是这样的，但我认为即使如此似乎也不会造成很大的影响。无论如何这里并不打算探讨博弈论，而且我们假定最好的策略就是设法提供那种人会很自然地给出的答案。\footnote{}

当这种测验被提出并讨论之后，图灵写到：

至于原来那个问题------``机器能思维吗？''，我相信它没什么意义，不值得讨论。不过，我相信到本世纪末，词的用法和广大受过教育的人的看法会发生很大的改变，因此那时将可以谈论机器思维而不怕再造成矛盾了。\footnote{}

图灵所预料的反对意见

由于图灵知道这种观点无疑会遇到暴风雨般的反对，他接着就简洁而又带有幽默感地剖析了一系列对``机器能够思维''这种想法的反对意见。我在下面给出他列出的九类反对意见，并采用他自己对它们的描述方式。\footnote{}遗憾的是，由于篇幅所限，这里不能复述他所给出的幽默而又机智的答复了。你也许会乐于自己思考一下这些反对意见，并给出你自己的答复。

⑴．来自神学的反对意见：思维是人的不朽的灵魂所具有的功能。上帝把不朽的灵魂赐给了每个男人和女人，但没有赐给任给其它动物或机器。因此动物和机器都是不能思维的。

⑵．鸵鸟式的反对意见：机器思维的后果将是非常可怕的。让我们希望并且相信它们不能做到吧。

⑶．数学化的反对意见（这基本上就是卢卡斯的论点）。

⑷．基于意识的论点：``除非一台机器能出于自身的思想和情感------而不是据符号的任意组合------写出一首十四行诗或一首协奏曲，否则我们不能同意说机器等于大脑------这就是说它不但要能写出来，还要能知道它已经写出来了。没有什么机器能够真正产生下列感受（而不仅仅是发出人为设置的信号------那是简单的发明）：在它成功时感到愉快，在它的部件损伤时感到痛苦，听到奉承话就觉得舒服，犯了错误就觉得悲伤，碰到异性会被吸引，得不到它想要的东西就会愤怒或失望。''（这是援引某位杰弗逊教授的话。）

图灵非常希望能对这种严厉的批评作出圆满的答复。因此，他为他的答案提供了不小的篇幅，而且在其中他又提供了一段假想的对话：\footnote{}

询问人：你的十四行诗中的第一行是：``我怎能把你比作夏天''，为什么不是``春天''呢？那样不是更好吗？

被问者：那样就不符合格律了。

询问人：那么``冬天''怎么样？这符合格律。

被问者：对，但谁也不想被比作冬天。

询问人：提到匹克威克先生，是否会令你想到圣诞节？

被问者：可以这么说吧。

询问人：因为圣诞节是在冬天，所以我想你如果把匹克威克先生比作冬天，他是不会计较的。

被问者：我觉得你不是在严肃地讨论问题。提到冬天，我们总是指冬天的一个一般的日子，而不是指像圣诞节这样特殊的日子。

在这段对话之后，图灵问道：``如果一台能写十四行诗的机器在口试时能这样回答问题，那杰弗逊教授还有什么话说呢？''其它的反对意见还有：

⑸．基于各种缺陷的论点：``就算你能造一台机器做你所提到过的那些事，你也绝不可能造出一台能做到X的机器。''有许多与此相关的X被提了出来。我在下面提供其中一部分：善良、机智、优美、友好、有创造性、有幽默感、区别是非、犯错误、坠入情网、爱吃草莓和奶油、使人动心、从经验中学习、用词得当、成为自己思考的对象、像人那样行为多变、做某些真正新的事情。

⑹．洛芙莱丝命妇的反对意见：关于巴比奇的分析机，最详尽的信息来自洛芙莱丝命妇所写的一本传记。她在其中写到：``分析机并不自称能\heititext{创造}任何东西。它能做\heititext{任何我们知道怎样命令它去做的事}''（黑体是她用的）。

⑺．基于神经系统连续性的论点：神经系统肯定不是一台离散状态的机器。一个撞击神经系统的神经脉冲的强度这样的信息若是有一点误差，都可能使输出脉冲的大小产生很大差别。因此可以争辩说，由于这个原因，我们不能指望用一个离散状态系统来模拟神经系统的行为。

⑻．基于行为非形式化的论点：这种观点大概是这样的：``如果每个人都有一个确定的指导规则集来控制他的生活，那他也就无异于机器了。但并不存在这种规则，因此人也就不可能是机器。''

⑼．基于超感官知觉的论点：假设我们要进行那种模拟游戏，而两个被问者一个是能接受心灵感应的人，另一个是数字计算机。询问人可以问这样的问题：``我右手里这张扑克牌是什么花色的？''那个人利用心灵感应或特异功能，在400次询问中能说对130次。而那台机器只能随机地猜测，也许只能说对104次，这样，询问人就能作出正确的判断了。

正如你所能看到的那样，许多反对意见彼此重叠、互相联系。我已经在本书中设法对其中的大多数作出了这样或那样的回答，但一般都不如图灵的那么简洁。

我觉得第9种反对意见是值得注意的。我曾见到图灵这篇文章被重印在某本书里，但第9种反对意见被略去了------我觉得这同样是值得注意的。我自己的观点------和图灵的相反------是认为超感官知觉并不存在。图灵也不想承认真的有超感官知觉，不过后来还是被迫承认了，原因是他那种杰出科学家的诚实精神要求他接受他认为是强有力地支持了超感官知觉的统计证据。我不这样看，尽管我把这个问题当成一个极其复杂而迷人的问题。

``帕里遇上了医生''

在本章前面的对话中，你已经看到了计算机程序和人交谈时的真实情况。还有许多``交谈''能力比施德鲁差的程序。其中有两个很著名，一个是约瑟夫·魏增鲍姆设计的``医生''，另一个是肯尼思·科尔比设计的``帕里''。前者是要模拟一个使用``启发式疗法''的精神病医生------也就是那种只靠问问题来治病的心理医生，而后者是要模拟一个妄想狂患者的信念结构。人们通过计算机网为这两个程序安排了一次有趣的正面较量，下面就是谈话记录中的一小段：\footnote{}

帕里：我不理解你的动机。

医生：你为什么不理解我的动机呢？

帕里：我难道不应该吗？

医生：你为什么现在这样说？

帕里：事实就是如此。

医生：假设事实并不是如此。

帕里：然后呢？

医生：我不知道我是否完全理解你了。

帕里：你总是唠叨同一个问题。

医生：这种相似性对你有什么启发？

帕里：我不太明白你的意思。

医生：你希望你太明白我的意思吗？

帕里：你现在该对我的感情表现出兴趣了。

医生：哦？我对你的感情表现出兴趣了吗？

当然，这些程序是被设计来和人谈话的，因此它们在这里互相对话时表现得很糟糕。它们大多是依赖于对输入的性质进行巧妙的猜测（它们对输入的分析十分肤浅），然后在大量事先准备好的回答中仔细地选出一些，再把它们吐出来。这些答案可能仅仅是部分地预先设计好的。例如，可能是一块模板，其中有些空白可供填写。其实是假定和它们交谈的人会从它们所说的话中读出比实际要多的意义来。事实上，按照魏增鲍姆在其著作《计算机能力和人类理性》中所说的那样，情况也的确是如此。他写道：

``伊莉莎''（即生成``医生''的那个程序）在与它交谈的许多人心中造成了很强烈的错觉，使他们感到它已经理解了他们所说的话。\ldots\ldots 他们常常要求与该系统私下交谈。而且在谈了一段时间之后，尽管我作出了解释，他们仍会坚持说机器真的理解了他们。\footnote{}

对于上面谈到的这种情况，你可能会感到难以置信。的确难以置信，但这是真的。魏增鲍姆对此提供了一种解释：

大多数人对计算机一无所知。因此，除非他们抱有极强的怀疑心理（就像我们在看魔术表演时那样），否则他们如果想解释计算机的那些智能行为，就只能采用他们所能采用的唯一的类比，即类比于他们自己的思维能力模型。这样一来，他们那些过火的想法也就不足为怪了。比如说，当一个人仅具有``伊莉莎''的语言能力时，的确无法设想他能模拟``伊莉莎''的行为。\footnote{}

这就相当于承认，这种程序基本上是虚张声势和连蒙带唬的一种巧妙的混合，是利用了人们的轻信心理。

根据这种稀奇古怪的``伊莉莎效应''，有些人提出应当对图灵测验进行修正，因为显然用极简单的花招就可以使人上当受骗。还有人提议询问人应当由得过诺贝尔奖的科学家来担任。更可取的办法可能是把图灵测验反过来，要求询问者也应当由一台计算机充当。或者说设立两个询问者------一个人和一台计算机------和一个被问者，这两个询问者要设法说出这个被问者是人还是计算机。

更严格来说，我个人认为图灵测验的原始形式是很合理的。至于对那些魏增鲍姆宣称被``伊莉莎''所骗的人们来说，问题在于没有提醒他们保持怀疑精神，用全部智慧来设法确定和他们笔谈的那个``人''是否真的是人。我认为图灵对这个问题的见解是正确的，因此图灵测验会本质上保持不变地存在下去。

人工智能简史

我希望能用下面几页的篇幅，以一种可能是非正统的观点，来讲述某些旨在阐明``智能背后的算法''的努力。我们已经遇到了许多错误和挫折，今后还会是这样。但不管怎么说，我们还是学到了许多东西，而且这的确是一个激动人心的时期。

甚至从帕斯卡和莱布尼茨开始，人们就梦想着用机器来完成智能性的任务了。在十九世纪，布尔和德·摩根发明了``思维定律''------本质上就是命题演算------因此就迈出了通往人工智能软件的第一步。查尔斯·巴比奇设计了第一台``计算机械''------现代计算机硬件的前身，也就是人工智能硬件的前身。我们可以把人工智能出现的时刻定义成机械装置完成了某些以前只能用人脑完成的任务之时。现在很难回过头去想象，那些人初次看到齿轮能完成大数的加法和乘法时会有何感想。或许他们在看见``思想''从物理的硬件中涌流而出时，体验到了一种敬畏之感。无论怎样说，我们确实知道，在近一个世纪之后，当第一台电子计算机构成的时候，其设计者们是体验到了在面对另一种``思维生物''时的一种敬畏、神秘之感。在那种情况下到底在多大程度上发生了真正的思维，这是许多困惑的根源所在。甚至在几十年后的今天，这个问题仍是刺激与苦恼的一个丰富的源泉。

有趣的是，今天实际上已经没有人再有敬畏感了------尽管和当时那些令人后脊梁骨上冒凉气的东西相比，现在计算机所完成的工作已经复杂得令人不可思议了。``大电脑''这个一度使人激动不已的名词现在只不过是一个过时的陈词滥调，成了旧时代的一个可笑的遗迹。我们这样快就对它感到厌倦了，这多少有点令人悲哀。

有一个与人工智能的进展相关的``\heititext{定理}''：一旦某些心智功能被程序化了，人们很快就不再把它看作``真正的思维''的一种本质成分。智能所固有的核心永远是存在于那些尚未程序化的东西之中。这个``\heititext{定理}''是拉里·泰斯勒首先向我提出的，因此我称它``\heititext{泰斯勒定理}''\heititext{：``人工智能是尚未做到的东西''。}

下面我来对人工智能作一个有选择的概述，以展示若干个集中了研究者努力的领域，其中每个领域似乎都以其自身的方式需要智慧的精华。对其中的某些领域，我又根据所使用的方法或所特别关注的问题作出了更细致的划分。

\begin{itemize}
\tightlist
\item
  机器翻译
\item
  \begin{itemize}
  \tightlist
  \item
    直接的（词序重新排列过的查字典过程）
  \item
    间接的（通过某种作中介的内部语言）
  \end{itemize}
\item
  博弈
\item
  \begin{itemize}
  \tightlist
  \item
    国际象棋
  \item
    \begin{itemize}
    \tightlist
    \item
      盲目的超前搜索
    \item
      带修剪机制的启发式超前搜索
    \item
      无超前搜索
    \end{itemize}
  \item
    跳棋
  \item
    围棋
  \item
    五子棋
  \item
    桥牌（叫牌和打牌）
  \item
    扑克
  \item
    三子棋的各种变形
  \item
    等等
  \end{itemize}
\item
  数学各分支中的定理证明
\item
  \begin{itemize}
  \tightlist
  \item
    符号逻辑
  \item
    \begin{itemize}
    \tightlist
    \item
      ``归结法''定理证明
    \end{itemize}
  \item
    初等几何
  \end{itemize}
\item
  数学表达式的符号处理
\item
  \begin{itemize}
  \tightlist
  \item
    符号积分
  \item
    代数式化简
  \item
    无穷级数求和
  \end{itemize}
\item
  视觉
\item
  \begin{itemize}
  \tightlist
  \item
    印刷品：
  \item
    \begin{itemize}
    \tightlist
    \item
      识别取自某小集合中的单个手写字符（如数字）
    \item
      阅读各种字体的西文文字材料
    \item
      阅读手写的西文段落
    \item
      阅读印刷体的中文或日文
    \item
      阅读手写体的中文或日文
    \end{itemize}
  \item
    图像：
  \item
    \begin{itemize}
    \tightlist
    \item
      在照片上寻找事先说明的物体
    \item
      把一幅图像分解成独立的物体
    \item
      在一幅图像中辨识出独立的物体
    \item
      根据人提供的简图识别物体
    \item
      识别人的面容
    \end{itemize}
  \end{itemize}
\item
  听觉
\item
  \begin{itemize}
  \tightlist
  \item
    在有限的词汇范围内理解口头表达的词（如十个数字的名称）
  \item
    固定的领域中理解连续的话语
  \item
    \begin{itemize}
    \tightlist
    \item
      找到音素间的界限
    \item
      辨识音素
    \item
      找到词素间的界限
    \item
      辨识词素
    \item
      整合出完整的词和句子
    \end{itemize}
  \end{itemize}
\item
  自然语言理解
\item
  \begin{itemize}
  \tightlist
  \item
    在特定领域中回答问题
  \item
    复杂句的语法分析
  \item
    对文字材料进行分段
  \item
    使用关于现实世界的知识以理解一段话
  \item
    解决有歧义的指称关系
  \end{itemize}
\item
  自然语言生成
\item
  \begin{itemize}
  \tightlist
  \item
    抽象的诗（如俳句）
  \item
    随机的句子、段落或更长的文字片段
  \item
    从内部知识表示中产生输出
  \end{itemize}
\item
  构造有独创性的思想或艺术作品
\item
  \begin{itemize}
  \tightlist
  \item
    写诗（如俳句）
  \item
    写小说
  \item
    计算机绘画
  \item
    音乐创作
  \item
    \begin{itemize}
    \tightlist
    \item
      无调性的
    \item
      有调性的
    \end{itemize}
  \end{itemize}
\item
  类比思维
\item
  \begin{itemize}
  \tightlist
  \item
    几何形状（``智力测验''）
  \item
    根据一个数学领域中的证明构造另一个相关领域中的证明
  \end{itemize}
\item
  学习
\item
  \begin{itemize}
  \tightlist
  \item
    参数调整
  \item
    概念形成
  \end{itemize}
\end{itemize}

机器翻译

在我下面要进行的讨论中，我将有所选择，略去上面列出的一些课题。不过，把它们列出来也是必要的，否则这份清单就不准确了。前面的几个课题是按历史次序排列的，每一个课题都没有达到预期的目的。例如，在机器翻译中所遇到的障碍使许多人大吃一惊。他们原以为这大概是个简单明确的任务，要做得尽善尽美的确很难，但基本上实现应该是很容易的。然而实践证明，翻译远不只是查字典和把词重新排列所能完成的。难点也不仅在于缺乏关于习惯用语的知识。事实上，翻译涉及到要为所讨论的领域建立一个心智模型，然后去处理那个模型中的符号。一个程序如果在读入文章段落时不使用一个领域模型，那它不久就会在遇到意义不清或多重意义时毫无希望地陷入困境之中。甚至一个人------他和计算机相比占有很大的优势，因为他对这个世界具有深入的理解------给他一段用他不懂的语言写成的文章和一本字典的时候，也会发现他几乎不可能把这段文字翻译成他自己的语言。因此------这在回顾时就不足为奇了------人工智能的第一个难题立即就涉及到了人工智能的核心问题。

计算机弈棋

计算机下国际象棋也被证明比开始时的直觉估计所设想的要困难得多。在这里，实践再次说明，人在心里表示一个棋局的方式，要远远复杂于仅仅知道哪个棋子在哪个位置上，以及下棋的规则。这涉及到感知若干相关的棋子所形成的结构，以及``启发式''的知识（或者说是估计），因此与高层组块有关。虽然不像正式规则那样严格，但启发式规则为了解棋盘上进行的事件提供了捷径，而这是正式规则做不到的。人们对这一点在一开始就有所认识，但只是过低地估计了对棋局的这种直观的、组块化的理解在人的棋艺中所起的巨大作用。当时的预想是，一个程序只要具有某些基本的启发式规则，再加上计算机在对局时进行盲目的超前搜索时的速度和精度，通过对每种可能的步骤进行分析，就能轻而易举地击败第一流的人类棋手------这个预想甚至在许多人紧张地工作了二十五年之后的今天，仍然远远没有成为现实。

时至今日，人们在从不同的角度处理下棋问题，其中最新颖的想法之一就是认为超前搜索是劳而无功的。反之，我们只需看看当前的局势是怎样的，并使用某些启发式规则生成一个计划，然后寻找一个能推进这个特定计划的棋步。当然，下棋计划的形成规则中必须要包括一些启发式知识，它们在某种意义下是超前搜索的``压缩''形式。这就是说，在许多盘棋中的超前搜索经验的等价物被``凝固''成另一种形式，而这种形式表面上并不需要超前搜索。在某种意义下这是个文字游戏。但如果这种``压缩的''知识在提供回答时比实际的超前搜索更有效------即使它偶尔会误入歧途也罢------那我们就已经有所收获了。像这样把知识蒸馏成更有效的形式，正是智能的长处所在，因此在下棋时更少地进行超前搜索，可能是进行研究时的一条行之有效的路径。尤其吸引人的想法是设计这样一个程序，它自己能把从超前搜索中得到的知识变换成``压缩的''规则------但这是个艰巨的任务。

塞缪尔的跳棋程序

事实上，在阿瑟·塞缪尔所设计的令人赞赏的跳棋程序中，已经采用了这样一种方法。塞缪尔所用的技巧是同时使用``动态''（超前搜索）和``静态''（非超前搜索）两种方式来评价任何给定的棋局。

静态方法采用一个由若干适用于任意棋局的特征量所构成的简单函数，这样就可以立即实际计算出来。而动态估值方法需要构造一棵``树''，其中表示了所有后面可能出现的棋步、对它们的反应、对这些反应的反应，如此等等（如图38所示）。在静态估值函数中有一些可以改变的参数，对它们进行改变的结果就会提供一组静态估值函数的可能的不同形式。塞缪尔的策略是以一种渐进的方式为这些参数选择越来越好的值。

该过程是这样完成的：每当程序要对一个棋局进行估值时，它既静态进行也动态进行。通过超前搜索所得到的答案------让我们称它D------被用来确定下一步要走的棋。而静态估值的结果S被用于一种更巧妙的目的：在每一步中，都对那些可变参数进行轻微调整，以使得S尽可能地接近D。这样做的结果是使通过对树的动态搜索所得到的知识被部分编码于静态估值函数的参数值里了。简而言之，基本想法就是把复杂的动态估值方法``压缩''到简单得多也有效得多的静态估值函数中来。

这里有一种良好的递归效应。这是因为对任意一个棋局的动态估值都需要进行有限步数的超前搜索------比如说七步。这样一来，走了七步之后所遇到的那些棋局本身也需要以某种方式进行估值。但当程序对这些局势进行估值时，它显然不能再超前搜索七步了，否则它将不得不超前搜索十四步，然后是二十一步，如此等等------这是一种无穷回归。因此，它在超前搜索七步后就必须依赖于静态估值了。这样，按照塞缪尔的办法，一种复杂的反馈就出现了。只要程序不断地设法把超前搜索估值``压缩''到一个较简单的静态方法中，这种方法转过来就会在动态的超前搜索估值过程中发挥关键性的作用。于是两种方法就密切地联系起来了，而且每种方法都以递归的方式从另一种方法的改进中受益。

塞缪尔的跳棋程序棋艺高强，基本达到了世界上最好的棋手的水平。如果是这样的话，那为什么不把同一技术应用于国际象棋呢？为了研究计算机下国际象棋的可行性，在1961年召集了一个国际性的委员会，其中包括荷兰的国际象棋特级大师兼数学家麦克斯·尤伟，委员会所得到的结论是前景暗淡：塞缪尔的技术如想应用于国际象棋，则难度大约是应用于跳棋的一百万倍，这似乎等于宣布此路不通。

这个跳棋程序所拥有的高超棋艺不能被解释成``已经达到了智能''，但对它也不能估计过低。在这个程序中结合了若干方面的知识，包括跳棋是什么、如何思考跳棋问题、以及怎样编程序。有些人可能会觉得它所表现出来的实际上是塞缪尔本人的下棋本领。但这是错的，原因至少有两个。其一是棋艺高超的棋手在选择棋步时所依赖的心智过程，往往连他们本人也无法全部理解------他们在使用自己的直觉。现在还没人知道怎样才能揭示自己的所有直觉。我们在内省过程所能采用的最好方法是使用``感受''或``元直觉''------这是一种关于直觉的直觉------作为向导，并且设法描述我们认为我们自己的直觉是什么样子的。但这只不过是为直觉方法的真实的复杂性提供了一种粗略的近似而已。因此实际上塞缪尔肯定无法把他个人的下棋方法映射到他的程序中去。说塞缪尔的程序的棋术不应与塞缪尔本人的棋术相混淆的另一个理由是：塞缪尔的跳棋还不如他的程序下得好------它能赢他。这一点也不会造成悖论------就好像事实上如果给一台计算机编上计算π值的程序，那么它一定会比给它编这个程序的人算得快。

程序何时才有独创性？

这个程序超过其设计者的论题联系着人工智能中的``创造性''问题。如果一个人工智能程序想出了一个主意、或一个下棋时的策略，而这是其设计者所从未想到的------那该归功于谁呢？关于这一点，已经出现了各种有趣的实例，其中有一些是微不足道的，但也有一些涉及到很深的层次。最著名的例子之一是一个为初等欧几里得几何中的定理寻找证明的程序。这个程序的最初构思来自马尔文·明斯基，实际设计者是IBM公司的赫伯特·吉伦特。有一天，这个程序为一个基本几何定理找到了一个精妙绝伦的证明------这个定理就是所谓``ponsasinorum''，即``驴桥''。它得了这么个诨名，是因为头脑简单的人很难通过它。

\pandocbounded{\includegraphics[keepaspectratio]{./media/OEBPS/Images/img_114.png}}

图114．``驴桥''的证明

（由帕普斯{[}约公元300年{]}和吉伦特的程序{[}约公元1960年{]}所发现）。问题：证明等腰三角形二底角相等。解：由于该三角形是等腰的，故AP和AP′等长，因此三角形PAP′和三角形P′AP是全等的（边、边、边）。这就说是对应角是相等的。具体来说，两个底角是相等的。

这个定理说：一个等腰三角形的两个底角是相等的。它的标准证明要求作底边的高线，把三角形分成对称的两半。这个程序发现的高明方法（见图114）不需要辅助线。这种方法是把这个三角形和它的镜像看成两个不同的三角形。然后，证明它们是全等的，并据此指出二底角在这一全等关系中彼此对应，这就证明了它们相等。

这个绝妙的证明使程序的设计者和其他人大为兴奋，有些人在它的行为中看到了创造才能的证据。实际上，在公元300年几何学家帕普斯已经发现了这个证明，但这一事实并没有使这项成就黯然失色。不管怎么说，下述问题依然存在，即``功劳该归于谁？''。这能算是智能行为吗？还是说这个证明本来是深藏在人（吉伦特）的心里，计算机只不过是使它浮到表面上来了？后面这个问题差不多打中了要害。我们可以把它反过来：证明是本来深藏在程序之中，还是接近程序的表面？这也就是要问，要使多大劲才能看清程序为什么要做它所做的事？程序的发现是否能归因于其中某个简单的机制，或机制的简单组合？也许，这里存在着一种复杂的相互作用，即使我们听到了对它的解释，也不会减少我们对它的出现所产生的敬畏之感？

我们似乎完全可以这样说：如果我们能把程序的行为归因于其中某些容易被追踪的操作，那在某种意义上，程序所做的只不过是揭示出那些在本质上是埋藏在------虽然并不很深------程序设计者头脑中的想法。相反地，如果对程序的追踪无助于说明为什么这个特定的发现得以产生，那我们也许应当开始把程序的``头脑''与其设计者的头脑加以区别了。程序的构造应归功于人，但不能说在他的头脑里已经有了程序所产生的那些想法。在这种情况下，人可以被说成是``元作者''------即作者的作者，而程序是（普通的）作者。

在吉伦特及其几何机器这种特定的情况下，尽管吉伦特可能并没有重新发现帕珀斯的证明，产生这一证明的机制离程序表面还是太近了。这使得我们难于毫不犹豫地认为该程序本身就有资格被称为一个几何学家。如果这个程序一而再、再而三地提出机智得令人吃惊的新证明，而且其中的每一个看起来都是基于天才的思想火花，而不是基于某种标准的方法，那时我们会坦然地把这个程序称为一个几何学家------但这种情况尚未出现。

计算机音乐是谁创作的？

作者与元作者之间的区别在计算机作曲的情况下表现得尤为突出。在作曲过程中，一个程序看上去拥有不同层次上的自主权。在以贝尔实验室的麦克斯·马修斯为``元作者''的一段曲子中就表现了其中的一种。他把两首进行曲的乐谱送入了计算机，一首是《约翰尼回家乡》，另一首是《英国掷弹兵》，并要求计算机生成一首新的乐曲------它开始时是``约翰尼''，但要慢慢变成``掷弹兵''，当乐曲演奏到一半的时候，``约翰尼''完全消失了，我们只能听到``掷弹兵''\ldots\ldots 然后这一过程又反了过来，当乐曲结束时，又和开始时一样，只剩下``约翰尼''了。用马修斯的话说，这是

\ldots\ldots 一个令人厌恶的音乐实验，但也并非毫无趣味，尤其是在节奏的转换上。``掷弹兵''本来是2/4拍子、F大调，而``约翰尼''本来是6/8拍子、E小调。从2/4拍子变成6/8拍子可以很清楚地被鉴别出来，虽然音乐家演奏起来很困难。从F大调转到E小调，这就要求在音阶中改变两个音符，结果听起来很刺耳，因而一个更小一点的转调无疑会是个更好的选择。\footnote{}

这样生成的曲子听起来有点奇异之感，虽然有些地方是浮华而混乱的。

是计算机在作曲吗？最好还是不要提这个问题，但对它也不能完全置之不理。要提供一个答案是很难的。所用的算法是确定的、简单的，也不难理解。其中并没有涉及复杂的或难以理解的计算，没有使用``学习''程序，没有出现随机过程，机器是以一种完全机械化的、简单明确的方式在运行。但是，其结果是一个并没有被作曲家完全计划好其细节的声音序列，尽管这一段的总体结构是被完全、精确地说明了的。这样，作曲家就常常吃惊地、而且是又惊又喜地看到了他的想法的实现细节。计算机只是在这种意义下作曲。我们称这一过程为算法化的作曲，但我们马上要再一次强调这些算法是能一目了然的。\footnote{}

这就是马修斯对那个他认为最好``废问''的问题的答复。但尽管他拒不承认，许多人还是发现简单地把这段曲子说成是``计算机作的''更方便。我觉得这种说法完全歪曲了实际情况。这个程序中没有任何结构类似于人脑中的``符号''，在任何意义下也不能说它在``思考''它所做的事。把这样一段音乐的创作归功于计算机，就像把本书的创作归功于产生它的计算机编辑照排系统一样。

这就引出了一个问题，这个问题不在人工智能领域中，但离得也并不远。这就是：如果你在一段文字中看到``我''这个字，你认为它指的是什么呢？举例来说，如果你在一辆脏汽车上发现``洗我''这样两个字，那这个``我''是指谁呢？也许这是某个可怜的孩子的呼声，他非常想洗个澡，所以就随便在一个地方乱画出了这两个字？也许是这辆车需要洗刷？也许是这句话想洗一次淋浴？也许，还有可能是要求把我们语言中那些乌七八糟的东西清洗干净？你可以把这个游戏继续做下去。在这里，这两个字是一个玩笑，是设想你在某个层次上会假装相信是这辆汽车自己写了这两个字，要求把它洗刷一下；而在另一个层次上，你又能清楚地看出这是一个孩子的恶作剧，他看着别人产生了误解，觉得很好玩。事实上，这种游戏是基于在不适当的层次上看``我''这个字的。

这种歧义现象在这本书里也已经出现过了，先是在《对位藏头诗》中，后来是在对哥德尔串G（及其相关物）的讨论之中。对不可播放的唱片的说明是：``我不能在唱机X上播放''，对不可证明的判断的说明是：``我不能在形式系统X中被证明''。让我们以后一句话为例。在任何别的场合，当你碰到一句包含有代词``我''的话语，你是否还可能自然地认为它不是指说这句话的人，而是指这句话本身？我猜这种可能性很小。``我''这个词如果出现在莎士比亚的一首十四行诗里，那它指的绝不是印在纸上的这十四行诗句，而是指幕后那个有血有肉的人，这个人是不在场的。

我们通常把一句话中的``我''追溯到多远呢？其答案在我看来，是我们要找到一个有感知力的主体，认为他（它）就是说这句话的人。但什么是一个有感知力的主体呢？是指我们可以方便地拟人化的东西。在魏增鲍姆的``医生''程序中，是否存在一个人格呢？如果存在，那是谁的？关于这个问题的一场小型论战近来正在《科学》杂志上激烈地进行着。

这又把我们带回到原来的论题上：到底是``谁''创作了计算机音乐？在大多数情况下，这些曲子背后的驱动力是人的智慧，而计算机只不过是一种带有或多或少的创造性的工具，被用来实现由人所提出的想法。完成这种工作的程序和我们人类毫无共同之处。它只不过是一个头脑简单的软件，毫无灵活性，不知道自己在做些什么，缺乏自我意识。但是如果有一天人们真的开发出了具有上述属性的程序，而且它们源源不断地创作出新的曲子来，到那时我就会认为应当把功劳区分开：部分功劳归于程序设计者------因为他创造了这样奇妙的程序，部分功劳归于程序本身------因为它良好的音乐修养。而在我看来，只有当这种程序的内部结构是基于某种类似于我们头脑中的``符号''及其触发模式，并以此来处理复杂的概念意义时，这种情况才可能发生。这种内部结构的存在将赋予程序某些性质，使得我们在一定程度上感到可以说它们和我们是一样的。但就算到了那时，我觉得说``这支曲子是一台计算机作的''还是有点别扭。

定理证明和问题分解

现在让我们转回到人工智能的历史上来。人们在开始时企图程序化的智能活动之一是定理证明。从概念上看，这和为一台计算机编程序，让它在WJU系统中为WU找一个推导是没有什么区别的，唯一的不同是所涉及的形式系统往往比WJU系统复杂得多。它们是谓词演算的各种变形，而谓词演算是在命题演算中加入量词后得到的扩充形式。事实上，谓词演算中的大部分规则都已包括在TNT之中了。编写这种程序的关键是要灌输进一种方向感，这样程序才可能不盲目地到处乱撞，而是只在``有关的''途径上工作------所谓``有关的途径''，是指那些经过合理的检查，看来是引向所要得到的符号串的途径。

在本书中我们还没有大量地涉及过这类问题。你怎样才能确切地知道你是在朝着一条定理前进，而不是在毫无意义地空转呢？这是我希望能通过WU问题加以描述的东西之一。当然，不可能存在确定的答案，而这正是那些限制性定理的内容。因为如果你总能知道应走哪条路，你也就能构造一个算法来证明所需的定理了，这样就违反了\heititext{丘奇定理}。这样的算法是不存在的（我希望读者自己来考虑为什么从\heititext{丘奇定理}可以严格地推出这一结论）。但是，这并不意味着不可能产生出任何关于哪条路有希望、哪条路没希望的直觉来。事实上，最好的程序都有很复杂的启发式知识，使它们在进行谓词演算的演绎推理时，其速度能和训练有素的人相媲美。

定理证明的关键是要运用下列事实：你有一个总目标------也就是你所要生成的那个符号串------可以用来指导你的局部行动。已经发展起来的一种技术就是把总目标转化成推导的局部策略。这种技术被称为``问题分解''。它是基于这样一个想法之上的：每当你面对一个遥远的总目标时，通常总存在一些``子目标''，达到了它们就有助于达到总目标。因此，如果我们把给定的问题分解成一系列新的子问题，然后再把它们分解成子子问题，如此以递归的方式进行下去，我们最终总能得到一些非常简单的目标，而我们有把握通过若干步达到它们，或者说至少看起来好像是这样的\ldots\ldots{}

\pandocbounded{\includegraphics[keepaspectratio]{./media/OEBPS/Images/img_115.png}}

图115．芝诺为了从A到B所建立的无穷目标树

问题分解使芝诺陷入了困境。你可能还记得，芝诺从A走到B（把B看成目标）所用的办法是把这个问题``分解''成两个子问题：首先走一半，然后再走另一半。这样就把两个子目标``压入''------在第五章中的意义下------你的``目标栈''中了。然后二者中的每一个又都依次被两个子子目标所取代------如此下去，直至无穷。你最后得到的是一个无穷大的目标栈，而不是单个目标（图115）。从你的栈中弹出无穷多个目标将被证明是难以做到的事------当然，这只是芝诺的观点。

关于问题分解中的无穷递归，另一个例子是在对话《和声小迷宫》中，当阿基里斯希望一个无类型愿望得到批准的时候。对这个愿望的批准要延期到元怪物许可时才能进行，但为了得到表示许可的许可，她必须去召唤元元怪物------如此等等。尽管目标栈是无穷大的，阿基里斯还是实现了他的愿望。问题分解胜利了！

尽管我对问题分解进行了嘲笑，它对把总目标转换成局部目标来说仍不失为一种有效的技术。它在特定的环境中可以大显身手，如象棋残局，在这种情况下超前搜索技术常常表现不佳，甚至把搜索长度增加到荒谬的程度------比如十五步或更多------也无济于事。这是因为超前搜索技术不是基于``规划''的，它根本没有目标，只是盲目探索大量的可能性。有一个目标就使你能构成一种达到那个目标的策略，这是与机械的超前搜索完全不同的理念。当然，在超前搜索技术中，一个局面是否合乎需要是通过估值函数来度量的，而这也就间接体现了一系列目标，主要是不能让对方将死。但那样未免太不直接了。优秀的棋手在和超前搜索的下棋程序对弈后通常所得到的印象是：他们的对手太不善于形成计划或者策略了。

山迪和肉骨头

问题分解的办法并非总是有效的。在许多情况下它会翻船。让我们以一个简单问题为例。假设有一条狗，名叫山迪，它的一个人类朋友刚刚扔过来一块它最爱吃的肉骨头，但那块骨头飞过栅栏落在邻居的院子里了。山迪能隔着栅栏看到那块骨头，就躺在那边的草地上------多么有诱惑力呀！在离骨头大约十几米远的栅栏上有一道打开的门。它会怎么办？有些狗会一直向栅栏冲过去，然后站在它前面叫起来，另一些狗会蹿向那道开着的门，然后绕到那块可爱的骨头那里去。这两类狗都可以说是在运用问题分解技术，但是，这个问题是以不同的方式表示在它们头脑里，而这就造成了全部差别。那个叫着的狗把子问题看成（1）跑到栅栏前，（2）通过它，（3）跑向骨头------但第二个子问题太难了，所以它就叫了起来。另一类狗把子问题看成（1）跑到门口，（2）通过门，（3）跑向骨头。请注意一切都取决于如何表示``问题空间''------也就是说，把什么看成是对问题的``简单化''（即朝向总目标的行动），把什么看成是对问题的``复杂化''（即背离总目标的行动）。

改变问题空间

有些狗开始时试图直接跑向骨头，当它们遇到栅栏后，某个东西在它们脑子里卡嗒一响，它们于是就改变路线，跑向了那道门。这些狗认识到，某种开始看来似乎会增加初始情景和期望情景之间的距离的东西------即离开骨头，跑向门口------实际上会缩短这一距离。在开始的时候，它们混淆了物理上的距离和问题中的距离。根据定义，任何离开骨头的行动似乎都是坏事。但后来------通过某种方式------它们认识到可以对它们的看法做修改，重新看待什么能使它们``接近''那块骨头。在一个适当选择的抽象空间中，跑向门才是使狗接近骨头的一个途径！在每个时刻，狗都在``接近''------在这个新的意义下------那块骨头。这样，问题分解是否有用，就取决于你在头脑里如何表示你的问题了。在一个空间中有可能被看成是后退的行为，在另一个空间中可能被看成一个革命性的进步。

在日常生活中，我们往往会遇到并解决和上述狗与骨头问题相类似的问题。例如，假设一天下午我决定乘车去南边一百公里的地方，但我现在正在办公室，而且是骑自行车来上班的，那我必须先进行许多看起来是``南辕北辙''的行动，然后才能真正乘车向南行驶。我得先离开办公室，这就是说，往东走几米，然后沿楼里走廊转向北，然后再向西。然后我骑车回家。在这段路中间可能任何方向前进的情况都有，这样我才到了家。此后，我要经过一串短距离移动才能最后进入汽车中，然后上路。当然，我不是马上就往南行驶------我坐的车可能一会儿向北，一会儿向西或向东，这取决于汽车的线路。

我们对这一切丝毫也没有感到矛盾，在这样做的时候甚至没觉得有什么可笑之处。这个空间已经在我头脑里扎下了根，而其中物理上的逆行被感知成直接奔向目标的行动。因此，当我往北走的时候根本没感到任何有讽刺意味的东西。公路和走廊之类的东西起着中间渠道的作用，我没经过什么思想斗争就接受了它们，因此在对感知情景方式的选择中，有一部分只需要接受强加的东西就行了。但栅栏前的狗有时很难做到这一点，尤其是当骨头离得这样近------就在眼皮底下，看上去又是这般美妙的时候。而且如果问题空间以比物理空间更抽象一些的形式出现的话，人也常常像那只光会叫的狗一样缺乏远见。

在某种意义下，所有问题都可以说是抽象化了的狗与骨头问题。许多问题都不是针对物理空间，而是针对某种概念空间的。当你认识到直接指向目标的行动在那个空间中会使你碰上某种抽象的``栅栏''，你可能会在下列两种办法中选择一种：（1）试着以某种随机的方式离开目标，希望你会发现一个隐蔽的``门''，你可以通过它跑到骨头跟前去；（2）设法发现一个新的问题表示``空间''，在其中没有抽象的栅栏把你和你的目标隔开------然后你就可以在这个新空间中直接跑向目标了。第一种方法似乎太慢，而第二种方法似乎太困难、太复杂。然而，那些涉及到重建问题空间的答案多半是一种突然闪现的洞见，而不是一系列缓慢、审慎的思维过程的产物。也许这些直觉的火花就是来自智能的核心------不必说，它们的来源是我们满怀戒意的大脑所严格保守的一个秘密。

在任何情况下，麻烦都不在于问题分解本身可能导致错误。它可以说是个很好的技术。真正的难题比这还要深刻：你怎样为问题选择一个良好的内部表示呢？你是在一个什么样的``空间''中观察它呢？在你已经选好的空间中，你用哪种行动来缩短你和你的目标之间的``距离''？这可以用数学语言表示成一个在状态之间寻找适当的``度量''（即距离函数）的问题。你要找到这样一个度量，在这种度量之下你和你的目标之间的距离很短。

现在既然选择一个内部表示这件事本身又是一种问题------而且也是最复杂的问题------你可能会想到用问题分解技术转过来对付它！为了做到这一点，你必须用某种方式表示许多不同的抽象空间，这是一项极其复杂的工作。我不知道是否已经有人在这些方向上进行过尝试。这可能仅仅是一种理论上有趣并吸引人的建议，实际上完全不现实。不管怎样说，人工智能急需这样一种程序：它们能够``后退几步''，看看正在发生什么事，并根据这种观察重新调整方向，以完成手边的任务。不难编写一个擅长于某项单一工作的程序------这项工作当由人来完成时似乎是需要智能的，但这完全不同于编写一个智能程序！这两者之间的差别就如同黄蜂（见第十一章）------它们的先天本能行为看上去会使人误以为它们具有很高的智能------和观察黄蜂的人之间的差别一样。

再谈W方式和J方式

一个智能程序大概要具有相当好的通用性，能解决各种各样的问题。它应当能学着做各种不同的事情，并在此过程中积累经验。它应当能依据一组规则来工作，而同时又能在适当的时候后退几步，判断一下依据这组规则工作是否会有利于达到它的某些总目标。它应当能在需要时决定不再继续工作于一个给定框架中，并构造一个新的规则框架，在里面工作一段时间。

这一讨论中可能有许多东西使你回想起WU谜题的性质。例如，远离问题的目标会使你回想起通过构造越来越长的串远离WU，因为你希望这样能以某种间接的方式构成WU。如果你是``一条幼稚的狗''，你可能在串长超过两个字符后会觉得你在远离你的``WU骨头''。如果你是``一条老练的狗''，这种加长规则的使用就会有一种间接的理由，就像为了得到你的WU骨头而奔向门口一样。

在前面的讨论和WU谜题之间的另一个联系是两种操作方式的区分，这导致了对WU谜题本质的认识：机方式（J方式）和惟方式（W方式）。在前一种方式下，你被嵌入某个固定的框架之中；在后一种方式下，你总可以向后退几步以了解事物的概貌。了解事物的概貌就相当于选择一种表示方式，在其中可以进行工作；而根据系统的规则进行工作就相当于在那个选定的框架中试验问题分解技术。哈代对拉玛奴衍风格------特别是他甘愿修改自己的假设------的评论，描绘了在创造性思维活动中J方式和W方式之间的这种相互作用。

黄蜂斯费克斯在J方式下工作得很好，但它根本不能选择它的框架，甚至不能对它的J方式进行最细微的修改。即使同一种情况在它的系统中一而再、再而三地发生，它也不能发现。因为要发现这种情况，就得跳出系统之外，哪怕只是稍微出去一点也行。但它根本不注意这种重复中的同一性。当我们把这种想法（即不注意某些特定的重复事件的一致性）应用于我们自身时，将得到有趣的结果。在我们的生活中或许也存在一些一再出现的重复情景，对此我们每次都用同一个笨拙的方式去应付，原因是我们没看到事物的概貌，也就无法感知它们之间的同一性。有这种可能吗？这又引回到一个前面出现过的问题上，即``什么是同一性？''。它不久就会在我们讨论模式识别时成为一个人工智能的主题。

把人工智能用于数学

从人工智能的观点来看，数学在某种意义下是个非常值得研究的领域。每个数学家都觉得在数学中的思想之间存在着某种度量------即整个数学就是一张由成果构成的网络，其间存在着非常多的关联。在这张网中，有些思想之间很接近，而另一些则要通过复杂的途径才能相互联系。有时数学中两个定理相互接近是因为从其中一个可以容易地证出另一个来。还有些时候两个想法相互接近是因为它们彼引相似，甚至是同构的。在数学领域中，``接近''这个词有两种不同的意义。可能还会有其它一些意义。我们这种数学的``接近感''中是否存在客观性或普遍性，或者它是否在很大程度上只是历史发展中的一个偶然现象，这都是很难说的。处于不同数学分支中的某些定理在我们看来很难联系起来，我们可能会说它们是不相关的------但后来可能会发现某种关系，这使得我们不得不改变我们的想法。如果我们能把我们这种高度发达的数学接近感------不妨称之为一种``数学家内心的度量''------灌输到一个程序之中，我们或许就能生成一个初等的``人工数学家''。但这同时取决于我们是否能传达一种``简洁感''或``自然感''，这是另一块巨大的绊脚石。

在一些人工智能项目中，人们已经开始正视这些问题了。在麻省理工学院，有人开发了一组名为``麦克西玛''的程序，其目的就是帮助数学家对复杂的数学表达式进行符号处理。这组程序在某种意义下能感觉到``该往哪里走''------某种``复杂性梯度''引导它把我们普遍认为是复杂的表达式变成简单的。在麦克西玛中有一个程序叫``赛因''，它可完成函数的符号积分，它已被普遍承认在某些范围内比人强。它依赖于许多不同的技能，而这些技能对智能来说一般是必不可少的：大量的知识、问题分解技术、许多启发式规则以及某些特殊技巧。

另一个程序是斯坦福大学的道格拉斯·利内特编写的，它的目的是在非常初等的数学中发明概念并发现事实。它在开始工作的时候已经有了关于集合的概念，还有一些事先灌输进去的关于什么东西是``有趣的''的观念。它从这些东西出发，``发明''了计数的概念，然后是加法的概念，然后是乘法，然后------除了其它各种各样的玩意儿------是素数的概念。它走得如此之远，以至于重新发现了哥德巴赫猜想！当然，这些``发现''人们在成百------甚至上千------年以前就已经完成了。或许这可以部分地被解释为利内特用许多规则所表达的``有趣''感已经被他的二十世纪的数学训练所影响了，不过这还是给人留下很深的印象。这个程序似乎在完成这番精彩的表演之后就精疲力尽了。它的一个有趣的性质是它不能发展或改善它自己关于什么是``有趣的''的感觉。这个困难似乎属于它的上一个层次------也可能是其上的若干个层次。

人工智能的关键：知识表示

在上面举出的例子之中，有许多是为了强调：一个领域的表示方式，在很大程度上，决定了那个领域是怎样被``理解''的。一个程序如果仅仅能以预先确定的次序打印出TNT中的定理，那它还不能算是理解了数论。一个程序如果像利内特那个带有额外知识层面的程序那样，那就可以说是对数论有了初步的了解。或许，只有当一个程序将数学知识植根于广博的现实经验之中的时候，我们才会认为它是以与我们同样的方式在``理解''。正是这个``知识表示''问题成了人工智能的关键所在。

在开始的时候，人们认为知识是像句子那样一``包''一``包''地存在的，而且把知识注入系统的最好方式就是找到一种简单的方式把事实翻译成被动的小数据包。那样每个事实可能只不过是一条数据，可以被使用它的程序所存取。下棋程序证实了这种想法，在这些程序里，棋盘局势被编码在某种矩阵或列表之中，高效率地存在存储器里，可以把它们提取出来并用子程序加以处理。

心理学家早就知道，人类是以某种更为复杂的方式来存储事实的。但这一点只是近来才被人工智能研究工作者所重新发现。他们目前正面临着``组块化''知识的问题和过程性与描述性知识的区别问题。而后一个问题，正像我们在第十一章中所看到的，涉及到``内省可达的''和``内省不可达的''这样两类知识间的区别。

认为所有知识都应能被编码成被动的数据，这种朴素设想实际上是与关于计算机设计的大部分基本事实相矛盾的。这就是说，关于怎样加、怎样减、怎样乘等等的知识不是编码在数据中并存在存储器内的。事实上，这些知识不是放在存储器中的某个地方，而是体现在硬件的线路里。一个袖珍计算器并没有在其存储器中存放加法的做法，这个知识是编码在它的``内脏''中的。如果有人说：``请指给我看在这个机器中加法的做法存在哪里！''，那么这种地方是找不到的。

但在人工智能中大量的工作都是与这样一些系统有关的：它们大部分知识都存放在特定的地点------也就是采用描述性的方式。当然，某些知识必须被嵌入在程序中，因为否则的话所得到的就根本不是一个程序，仅仅是一部百科全书了。问题在于怎样才能把知识分成程序和数据。不管怎么说，把程序和数据区分开也并非总是轻而易举的。我希望这一点在第十六章中已经说得够明白的了。但在一个系统的开发过程中，如果程序员直观地感到某些项应当是数据（或程序），那将对该系统的结构产生相当大的影响，因为人编程序的时候总是倾向于把那些像数据的对象和那些像程序的对象区分开。

必须指出，原则上把信息编码成数据结构或过程的任何方式都是同样好的，也就是说如果你不太在意效率的话，你能用一种方式所做的一切也都能用另一种方式来做。但是，可以提出一些理由来说明一种方式似乎肯定比另一种优越。例如，请考虑下面主张只使用过程性表示的论点：``一旦你想要把相当复杂的特性编码在数据中，你就不得不开发出相当于一种新的语言或新的形式系统的东西。因此实际上你的数据结构变得像程序，而你的一些程序成了它们的解释程序，那你还不如一开始就直接把这些信息表示成过程的形式，这就用不着外层的解释程序了。''

DNA和蛋白质能对我们有所启发

上述论点听起来很有道理，但如果在解释时稍加发挥，就会变成一个要求取消DNA和RNA的论点。为什么要把遗传信息编码在DNA中呢？如果把它们直接编码在蛋白质里，你不就不仅能取消了一层解释工作，而且事实上是取消了两层解释工作吗？答案是：我们发现，使同一个信息为了不同的用途具有多种不同的形式，往往是极其有用的。把遗传信息存储在模块化的、像数据似的DNA形式中，其优点之一就是两个个体的基因可以容易地被重组，以形成一个新的遗传型。如果信息仅仅在蛋白质里，那就很难做到这一点了。把信息存在DNA中的另一个理由是：这样可以很容易地复制并翻译到蛋白质中。当不再需要它的时候，它占不了多大地方，一旦需要了，它又能被当作一块模板。不存在能从一个蛋白质中复制出另一个来的机制，因为它们曲折的三级结构将使复制过程变得极其复杂。反过来说，能把遗传信息表示成像酶那样的三维结构，这几乎又是绝对必要的。因为分子的识别和制造在本质上都是三维的操作。可见，这种主张使用纯过程性表示的论点，在细胞这一环境中是很不恰当的。这说明如果能在过程性和描述性表示之间进行相互转换，将带来许多优点。对人工智能来说恐怕也是如此。

这个问题是弗兰西斯·克里克在一次关于与外星人通讯的会议上提出的：

我们知道在地球上存在两种分子，其中一种便于复制（DNA），而另一种长于行动（蛋白质）。是否可能设计一个系统，其中一个分子可以身兼二任呢？还是说也许存在来自系统分析的很强的证据，说明（如果有这种证据的话）把任务分成两部分有很大的优点？对这个问题我不知该怎样回答。\footnote{}

知识的模块性

在知识表示中出现的另一个问题是模块性。要插入新知识有多难？要修改旧知识有多难？书籍的模块化程序有多高？这都要依情况而定。如果从一本高度结构化的------即含有许多交叉相关的------书中删去一章，这本书余下的部分可能会变得实际上无法看懂了。这就像想从一个蜘蛛网上拉出一根丝来一样------你这样做会毁掉整张网的。而在另一方面，某些书是非常模块化的，即各章相互独立。

请考虑一个使用TNT的公理和推理规则的简单的定理生成程序。这个程序的``知识''有两个方面。它隐含在公理和规则中，而又显含在它已经生成的定理中。你既可以把它看成模块化的，也可以把它看成分布在系统之中------因而是完全非模块化的，这取决于你以哪种方式看待这些知识。例如，假设你已经编写了这样一个程序，但忘了把TNT的公理1放在公理表中。当程序已经进行了数千次推导之后，你发现了你的疏忽，并把这个新公理插了进去。你转眼之间就能完成这次插入，这个事实表明系统中的隐含知识是模块化的，但这个新公理对系统的显含知识的贡献将要过很长一段时间才能反映出来------要等到它的影响已经向外``扩散''之后，就像当瓶子被打碎之后，香水的气味会慢慢扩散到整个房间中一样。在这种意义下，新知识要经过很长时间才能被体现出来。进一步来说，如果你想退回去用公理1的否定来取代它，你就不能仅仅这样做完就完了。你必须删除所有在其推导中涉及了公理1的定理。显然这个系统的显含知识不像它的隐含知识那样模块化。

如果我们能学会怎样模块化地移植知识，那将是很有用的。那时，要想教某个人学法语，我们只需打开他的脑袋，在他的神经结构上做个手术------然后他就知道怎样讲法语了。当然，这只不过是大白天作美梦而已。

知识表示的另一个方面与人希望以何种方式使用这条知识有关。当一条条信息到达时是否要进行推理？是否要不断地在新信息和旧信息之间进行对照和比较？例如，在一个下棋程序中，如果你要生成一个超前搜索树，那么一种能以最小的冗余量对棋局进行编码的表示法，就比以多种不同的方式重复地存储信息要好。但如果你想让你的程序通过寻找模式，以及把它们与已有模式相比较，来``理解''一个局势，那么把同样的信息以不同的形式重复表示多次将是更有用的。

在一个逻辑形式系统中表示知识

在什么是知识表示和加工的最佳方式这一问题上，存在着各种不同的思想流派。其中影响很大的一派主张用类似于TNT所用的那种形式化表示法------即采用命题联词和量词。这种表示法中的基本运算无疑是形式化了的演绎推理。逻辑演绎可以通过推理规则来完成，就像在TNT中那样。对系统中某个特定想法的询问就设立了一个目标，其形式是一个要推导出的串。例如，``无朋是一个定理吗？''然后自动推理机制开始在这个目标的引导下使用各种问题分解方法进行工作。

例如，假设已知命题``所有形式算术系统都是不完全的''，而对程序的询问是：``《数学原理》是完全的吗？''。通过扫描已知事实的表------通常称作``数据库''------该程序可能会注意到，如果它能证实《数学原理》是个形式算术系统，那它就能回答上述问题。因此命题``《数学原理》是个形式算术系统吗？''就可能被当作一个子目标来设立，然后再进行问题分解。如果它能找到更进一步的事实来帮助它证实（或否定）这些目标或子目标，它将对这些东西进行处理------如此递归地进行下去。这个过程被称之为``反向链接''，因为它是从目标开始反向工作，大体上是朝着已知的方向推进的。如果我们用图示法表示这些主目标、子目标、子子目标等等，我们就会得到一个树形结构，因为主目标可能包含若干个不同的子目标，每个子目标又可能包含若干个子子目标，如此等等。

要注意这个方法并不保证一定能解决问题，因为在系统中可能无法证实《数学原理》是一个形式算术系统。但这并不意味着这个目标或子目标是个假判断------只能说它们不能从系统当前所拥有的知识中推导出来。在这种情况下，系统可能会打印出``我不知道''或类似的话。一些问题被存而不论，这个事实显然类似于那种使某些著名的形式系统身受其害的不完全性。

演绎式认识之别于类比式认识

这种方法提供了对所表示领域的一种``演绎式认识''，即可以从已知事实中推导出正确的逻辑结论。但是，它忽略了人类某些发现相似性和对情境进行比较的能力------即忽略了那种被称为``类比式认识''的东西------而这正是人类智能的一个关键性的方面。这不是说类比思维过程不能被硬塞到这种模子中来，而是说用这种形式化系统无法自然地对它们进行把握。近来，面向逻辑的系统不像另一些系统那样吃香了，原因之一就是后者允许更自然地实现复杂的比较。

一旦你认识到知识表示绝不仅仅是一些数字的存储，那么``计算机过目不忘''的说法便是一个容易攻破的迷信了。``存在存储器中''的东西不一定就是一个程序``知道''的东西的同义语，因为即使一条给定知识是编码在一个复杂系统中的某个地方，如果没有过程或规则或其它类型的数据处理器可以得到它，那它也会是无法使用的。在这种情况下，你可以说这条知识已经被``遗忘''了，因为通向它的途径已经暂时地或永久地失去了。这样一个计算机程序可能会在高层``忘掉''某些它在低层实际还``记得''的东西。这是那些反复出现的层次间差别中的一种，而我们或许能从中进一步了解我们自己。当一个人在遗忘时，这往往是意味着失去了一个高层指针------而不是说某些信息被删除或毁坏了。这就进一步强调了，保持通向你所存储的经验的路径是极端重要的，因为你无法事先知道你将要在什么情况下------或以什么角度------从存储器中取出某些东西来。

从计算机俳句到RTN语法

人脑中知识表示的复杂性初次给我造成强烈的印象，是在我编写一个程序，使它随机地产生自然语言语句的时候。我是以一种很有趣的方式想到这个课题的。我曾在收音机中听到过所谓``计算机俳句''的几个例子，那里面有些东西深深地震动了我。让计算机生成某些在日常情况下会被当作艺术创作的东西，这具有极大的幽默感，又充满了神秘色彩。我既对其幽默的一面忍俊不禁，同时又被编写能从事创造活动的程序------这甚至有点自相矛盾------的神秘感所强烈吸引着。因此我就着手写了一个甚至比俳句程序更幽默、更神秘地自相矛盾的程序。

最初，我想把语法设计得具有灵活性和递归性，这样人们就不会觉得这个程序只是在某些模板中进行填空。正在那时我偶然在《科学美国人》上看到维克托·扬弗的一篇文章，他在这篇文章中描述了一个简单而又灵活的语法，用它可以产生各种各样能在某些儿童读物中发现的那类句子。我对我从那篇文章中得到的一些想法进行了修改，并编制了一组过程，它们构成了一个递归迁移网语法，就像我们在第五章中描述过的那样。在这个语法中，一个句子里词的选择是由一个过程所确定的，这个过程开始时先选择------随机地------该句子的整体结构，然后这个决策生成过程逐渐渗透到结构的低层，直至达到词的层次和字的层次。对英语来说，在词的层次以下还有不少事要做，例如动词的词尾变化和名词的单复数变形。不规则的动词和名词开始时也是以规则的形式出现的，然后如果它们与一张表中的表目相匹配，那时再把它们替换成适当的（不规则的）形式。一旦每个词都达到了其最终形式，就把它打印出来。这个程序就像是让一群猴子在打字机上随机地敲打键盘，只不过这种操作是同时在语言结构的若干层次上进行的操作------不仅仅是在字的层次上操作。

在程序开发的早期，我使用的是一组毫无意义的词汇------这是故意的，因为我想达到幽默的效果。程序生成了一大堆胡说八道的句子，其中一些具有很复杂的结构，而另一些则很简短。下面就是其中的一部分：

一支只会笨拙地大笑的雄铅笔定会嘎嘎大叫。程序必定不会总是在记忆里把姑娘嘎嘣嘎嘣地嚼吧？那个吐痰时笨了吧唧的十位数毛病可能会坍塌。肯定把一个突如其来的男人认作亲戚的那位蛋糕必然会一个劲儿地甩掉。

程序应该兴高采烈地运转。

这可尊敬的机器不得总是粘这位天文学家。

呜呼，真该把那位姑娘挟走的程序为剧院写乐师。商业式的关系嘎嘎叫。

那位总能嘎嘎大叫的幸运姑娘将永远不会没错嘎嘎叫了。

这个游戏嘎嘎叫。教授将写臭豆腐。有个毛病完蛋了。人拿了那位滑跤的盒子。

这些句子的效果带有很强烈的超现实主义色彩，而且有时会令人想起俳句------如最后那个例子中那四个连续的短句就是如此。它开始看起来很有趣，而且有一定吸引力，但不久就变得令人厌烦了。在读了几页输出之后，人们会感到程序运行于其中的那个空间的界限，而在那以后，再看到该空间中的随机点------尽管每一个都是``新''的------就不觉得有任何新鲜之处了。这在我看来是个一般规律：你若对某个东西感到厌烦，这往往不是由于你已经穷尽了它的所有可能的行为，而是由于你已经发现了包含其全部行为的那个空间的界限。只不过一个人的行为空间太复杂了，所以才总能使别人感到惊奇。但我的程序可并非如此。我认识到，要达到我那个产生出真正幽默的输出的目标，在程序中还需要编进去更多精妙的东西。但在这种情况下，``精妙''是指什么呢？显然那样荒谬地把词排列起来是太不精妙了，我需要找到一种方法来保证词确实是按照世界的真实情况来使用的。关于知识表示的思想就是从这里开始产生的。

从RTN到ATN

我所接受的想法就是把每个词------名词、动词、介词、等等------按若干不同的``语义维度''进行分类。这样，每个词就成了不同的类中的一个成员，此外还有一些超类------即类所组成的类（这使人想起了乌兰姆的话）。原则上说，这种聚集过程可以继续进行至任意多层，但我只设立了两层。在任意给定时刻，现在词的选择在语义上都是受限制的了。原因是要求在被构造的短语中的各部分都必须相互``符合''。这样做的想法，举个例子来说，就是某种行动必须只能由动物来完成，或只有某种特定的抽象概念才能对事件产生影响，如此等等。要确定哪些范畴是合理的，以及是否每个范畴最好都被想象成一个类或一个超类，这种判定过程是十分复杂的。所有的词都被依照若干不同的维度作了标记。有些词------``的''、``之中''等------具有若干不同的表目，对应于它们的不同用法。这样一来，输出就变得可理解多了------它们因此而以一种新的方式引人发笑。

一个小型图灵测验

下面，我要复述九句话，它们是从我那个程序的最新版本所产生的许多页输出中精心选择出来的。和它们混在一起的是三句由人（出于严肃的目的）所写下的句子。但哪三句是呢？

⑴．冲口而出会被认为是一种动力学反射作用中语义材料（复义的）和某种语义对话产物的交互替代物。

⑵．不如设想一``串''由阿赖耶识实验的许多愚伯们组成的小径，在那里后继世系乍看上去是一种发生了历时性倒错的过渡状态的情形。

⑶．设想一下最终作为产品（认识的种种条件？）而出现的某种链强度可能性，而该产品又不像汉堡包似的将其全部囊括于其中的情形。

⑷．尽管作了种种努力，这一回答，如果你想知道的话，其时已为东方所支持了。正因如此，这位大使即将持有的态度会使某种失误得以延缓发生。

⑸．当然，及至动乱发生之前，这位大使一直在积羽群轻、积微逐渐地宠惯这伙暴徒。

⑹．一如众所笃信的那样，完善的自由已导致这些态度直至这样一种程度，即和平将无法由这一命令最终所不可逆转地引起的后果提纯出来，而这种不可逆转到了这种地步，以至于和平有时正以无限小的惊人度导致拒不让步的态度。

⑺．换言之，按照智者的观点，各城邦中的攻势已被东方狡猾地接受。当然，东方已被这些城邦以异乎寻常的粗暴分裂了。东方支持着过去一直为人类所支持的种种努力。

⑻．然而应该承认，这种谬误在统治集团中的根源将会被其敌人预料到。同样，个人主义者们到时也将会证明拒不让步并未能终止这一攻势。

⑼．毋须说，在必然会确保了保密状态的那场动乱中，这些回答并不能分裂东方。当然，这一事实本身也就说明了这些国家总是在探测自由。

⑽．虽然人道主义者们获得了一项诺贝尔奖，然而除此之外，农奴也获得了。

⑾．这是一种为冲突所苦的国家中的农奴们所经常持有的态度。

⑿．此外，诺贝尔奖是可以得到的。同样地，尽管有这种后果，只要有人能获得诺贝尔奖，未来终将可以由一位妇女得到这种奖。

人写的句子是第1到第3句。它们是从当代的《艺术语言》杂志中摘下来的\footnote{}，而且------据我所知------是绝对严肃地致力于在有文化而且神志清醒的人之间彼此交流某些东西。它们出现在这里时离开了原先的上下文，但这并不会造成太大的误会，因为它们所处的上下文读起来和它们也差不太多。

其它句子都是我的程序所生成的。选择第10到第12句是为了说明有时这些句子完全是清楚易懂的；选第7到第9句表现了输出中的典型情况，它们漂浮于有意义和无意义之间的那个稀奇古怪而又引人入胜的阴曹地府之中；而第4到第6句则是远远超出了``意义''的。出于一种宽宏大度的心情，人们可能会说它们本身构成了一种纯粹的``语言对象''，就像是某种用词而不是用石头刻成的抽象派雕塑。或者，人们也可以说它们纯粹是冒充成有理智的胡言乱语。

我对于词汇的选择仍以产生幽默效果为目的。输出的风格特征是很难加以刻划的。尽管其中许多是``有意义''的------至少在单句层次上是如此，但人们还是不免要得到这样一种感觉，即产生输出的那个源泉根本不理解它在说些什么，而且在说这些话时根本没有目的。尤其是人们会觉得在这些词背后完全缺乏视觉表象。当我看着这些句子从行式打印机中源源而出的时候，我体验到了非常复杂的情感。输出的那些傻话使我觉得很好笑。我同时又对自己的成就感到骄傲，而且总是设法对朋友们把它描述成类似于给出一些规则，根据这些规则我们可以一挥而就地生成有意义的阿拉伯故事------这不免有些夸大其辞，但这样想使我很愉快。最后，我深深地感到激动，因为我看到这个极其复杂的机器正根据规则使其中的那些长长的符号列车进行``转轨''，而这些长长的符号列车就是某种类似于我自己头脑中的思想的东西\ldots\ldots 似乎是这样。

\pandocbounded{\includegraphics[keepaspectratio]{./media/OEBPS/Images/img_116.png}}

图116．一个用阿拉伯语写成的有意义的故事

摘自埃·卡迪比与莫·西哲儿梅西{[}A. Khatibi; M. Sjelmassi{]}，《伊斯兰书法大观》{[}The Splendour of Islamic Calligraphy{]}，纽约，Rizzoli，1976年版

关于思维的想象

当然，我不能自欺欺人地以为在这些句子背后存在着意识------远非如此。在所有的人中，我最清楚为什么这个程序和真正的思维有着天壤之别。\heititext{泰斯勒定理}在这里完全适用：一旦这个层次上的语言能力被机械化了，它就显然不再是智能的一部分了。但这种强烈的体验给我留下了一个想象：我隐隐约约地感到``真正的''思维是由大脑中更长、也更复杂的符号列车所组成的------许多列车同时沿许多彼此平行或交叉的轨道行驶，它们的车厢被推过去或拉过来，被挂上去或摘下来，被无数个神经转轨器从一条轨道上变换到另一条轨道上\ldots\ldots{}

这是一个不可捉摸的想象，我无法用语言来表达它。它仅仅是个想象而已。但想象、直觉和动机在头脑中是紧密地混在一起的，而我对这个想象的极大兴趣不断地鞭策我更加深入地去思考思维到底是什么。我已经设法在本书的其它部分表达了一些从这个原始想象中派生出来的子想象------尤其是在《前奏曲，蚂蚁赋格》之中。

当我以十几年后的观点回顾这个程序的时候，现在我脑海里最强烈的感想是：在它所说的话背后缺乏形象感。这个程序不知道什么是农奴、什么是人，或任何其它事物。这些词只是空洞的形式符号，就像pq系统中的p和q那样空洞------或许比它们还要空洞。我的程序的优势来自这样一个事实：当人们读一段文字的时候，他们很自然地趋向于为每个词填满其全部意思------就像这些意思必然地附在形成该词的这些笔划上一样。我的程序可以被看作一个形式系统，它的``定理''------即输出的句子------具有既定的解释（至少对使用这种语言的人来说是这样）。但和pq系统不同，在这样解释了之后，这些``定理''并非都是真陈述。其中有许多是假的，还有许多是无意义的。

pq系统以其微不足道的方式反映了世界的一个小角落。但当我的程序运行时，在它的内部并未反映世界是如何运行的，它只反映了它所必须遵循的那一小组语义限制。为了能反映对世界的理解，我必须得为每个概念包裹上一层又一层关于世界的知识。这样做将是和我以前的做法不同的另一种努力。我并非没想过设法做到这一点------但我一直不知该怎么干。

高层语法\ldots\ldots{}

事实上，我常常在思考我是否能写一个ATN语法（或某种其它的句子生成程序），它能只产生关于世界的``真''词句。这样一种语法能为每个词赋予名副其实的意义，就像在pq系统和TNT中所发生的那样。设计一种语言，使假陈述成为不合语法的，这个想法古已有之。这可以上溯到1633年的约翰·艾莫斯·夸美纽斯。这个想法很有吸引力，因为这样一来在你的语法中就嵌入了一块试金石：只需把你想要知道的陈述写下来，然后检查它是否合乎语法\ldots\ldots 实际上夸美纽斯走得更远，因为在他的语言里，假陈述不仅是不合语法的------它们根本就不可表示！

把这一思想引向另一个方向，你可以设想一个能生成随机公案的高层语法。为什么不可能呢？这种语法将等价于一个形式系统，其中的定理就是公案。而如果你有这样一个程序，你难道不能设法只生成``正宗的''公案吗？我的朋友玛莎·梅瑞狄丝热衷于这个``人工主义''的想法，因此她就着手编制一个公案书写程序。通过开始一段时间的努力，生成了下面这个古怪的``准公案''：

一个小的年轻禅师想要一个小的白色异形钵。``我们怎样才能不经过学习就能知道和理解呢？''这个年轻禅师问一个大的困惑着的禅师。这个困惑着的禅师带着一个小的红色的石钵从一座褐色的硬山走向一座白色的软山。这个困惑着的禅师看到了一个红色的软棚屋。这个困惑着的禅师想要这个棚屋。``菩提达摩为何要来中国？''这个困惑着的禅师问一个大的顿悟了的弟子。``桃子是大的''，这个弟子回答这个困惑着的禅师。``我们怎样才能不经过学习就能知道和理解呢？''这个困惑着的禅师问一个大的年老的禅师。这个年老的禅师来自一个白色的石G0025。这个年老的禅师消失了。

你个人的关于公案正宗性的判定过程可能已经对此做出了判决，而不需要所谓``几何编码''或``禅宗串技艺''。如果代词的缺乏、句法的单一和语言的``现代化''还没有引起你的怀疑的话，临近结尾处那个奇怪的``G0025''也一定使你起了疑心。它是什么东西？它是一个奇怪的毛病------是一个故障的外部表现，这个故障使得原来应当打印出一个代表对象的词的地方被打印成了这个``节点''（事实上是Lisp中的一个原子）在程序中的内部名称，在这个节点中存储着关于这个特定对象的全部信息。因此我们在这里就有了一个``窗户''，透过它可以看到禅宗思想的低层------这一层本来是应当看不见的。不幸的是，我们无法在人的禅宗思想上开出来这样一个明晰地通向低层的窗户。

上面公案中的那一串行为的序列，尽管带有一些任意性，还是来自一个叫做``多级瀑布''的递归Lisp过程，这个过程构造出一个彼此通过不明确的因果方式相互联结的行动链。尽管这个公案生成程序对世界的理解程度显然没有高得令人惊讶，这项工作在使其输出看起来更加纯正这一点上还是取得了进步。

音乐的语法？

下一个问题是音乐。你开始可能会猜想这个领域很适合于被整理在一个ATN语法或其它类似的程序中。因为（把这个朴素的想法继续推进）语言必须从与外部世界的联系中获得意义，而音乐却纯粹是形式的。在乐曲的音响中不存在与``外在的''东西的关联，只存在纯粹的句法------一个音符接着一个音符、一个和弦接着一个和弦、一个小节接着一个小节、一个乐句接着一个乐句\ldots\ldots{}

但请等等。在这段分析中有什么地方不大对头。为什么某些音乐要比其它音乐更深刻、更优美呢？这是因为音乐中的形式是具有表达力的------可以表达给我们头脑中某些奇特的潜意识区域。乐曲中的音响并不涉及农奴或者城邦，但它们触发了我们心灵最深处的模糊情感，在这个意义下，音乐的意义的确依赖于从符号到世界中的事物间的那些错综复杂的联系------在这种情况下，那些事物就是我们心中秘密的软件结构。不，伟大的音乐将不会出于像ATN语法这样简单的形式系统。伪音乐，像伪神话一样，是可能被产生出来的------对人来说这样做将是个有价值的尝试------但音乐中意义的秘密远不是靠纯粹的句法所能挖掘到的。

我在这里要明确一点：原则上说，ATN语法拥有任何程序化形式系统所具有的全部能力，因此如果音乐的意义归根结底是能以某种方式来加以把握的（我相信是这样的），那么一定能在一个ATN语法中被把握。这是真的。但在那种情况下，我仍然认为，这个语法的定义将不仅包含音乐结构，而且包含一个观察者的整个心智结构。这个语法将是全部思维的语法------而不仅仅是音乐的语法。

维诺格拉德的程序``施徳鲁''

哪种程序会被人承认为------即使是勉强地------具有一定的``理解力''？它要怎样做才能不使你直觉地感到其中``一无所有''？

在1968-1970年，特里·维诺格拉德（化名格拉德·维维诺诺博士）还是麻省理工学院的一个博士研究生，在研究同语言以及理解有关的课题。在当时的麻省理工学院，许多人工智能研究都涉及到所谓``积木世界''------一个相对简单的领域，很适于在其中处理涉及到视觉和计算机掌握语言的问题。积木世界中包括一张桌子，上面有各式各样的类似于儿童玩的那种积木------方的、长的、三角的、等等，各自有不同的颜色（作为另一种``积木世界''，请看图117：马格里特的画《心算》。我觉得这幅画的标题和内容再相配不过了）。在麻省理工学院的积木世界中，视觉问题是非常棘手的：一台计算机怎样才能通过对一个涉及到多块积木的场景进行电视扫描，指出其中有哪些种积木，以及它们之间的相互关系？某些积木可能叠在别的积木的上面，某些可能在另一些的前面，还有一些阴影，如此等等。

\pandocbounded{\includegraphics[keepaspectratio]{./media/OEBPS/Images/img_117.jpg}}

图117．心算

马格里特作（1931）

不过，维诺格拉德的工作是与视觉问题无关的。他开始时可以假设积木世界已经被良好地表示在计算机的存储器中了，而他所面临的任务是怎样让计算机完成下列工作：

⑴．理解用自然语言表达的关于情境的问题；

⑵．用自然语言回答关于情境的问题；

⑶．理解用自然语言表达的处理积木的请求；

⑷．把每个请求分解成由它所能执行的操作组成的序列；

⑸．理解它所做的事情以及它这样做的原因；

⑹．用自然语言叙述它的行动及其原因。

一个貌似合理的想法是把整个程序分解成模块化的子程序，每个模块处理上述不同问题，然后，当这些模块都分别开发好了之后，再把它们适当地组合起来。但维诺格拉德发现，这种开发独立模块的策略会造成一些根本性的问题。他发展了一种激进的方法，向那些认为智能可以被划分成独立或半独立的部分的理论进行挑战。他的程序施德鲁并没有把问题分解成概念上清楚的若干组成部分。对句子做语法分析、生成内部表示、对其内部所表示的世界进行推理、回答问题等等，这一系列操作都以一种过程性的知识表示形式深深地、内在地纠缠在一起了。有些批评家曾指责他的程序过于错综复杂，以致于根本没有表示任何关于语言的``理论''，也没有以任何方式对我们理解思维过程有所贡献。在我看来这种主张简直是错得无以复加了。像施德鲁这样的杰作可能并不同构于我们的所做所为------事实上不管你采用什么方法，也无法认为施德鲁已经达到了``符号层''------但是在它的构造和分析过程中，却可以为我们认识智能的工作方式提供大量的材料。

施德鲁的结构

事实上，施德鲁不能分解成若干个独立的过程，其中每一个包含关于世界的某些知识。施德鲁中过程之间的相互依赖性很强，因此不能明确地相互分开。这个程序就像一个无法解开的复杂绳结，但你不能解开它这一事实并不意味着你不能理解它。即使这个绳结从物理上看错综复杂，对它仍可能有一个漂亮的几何描述。我们可以回到《一首无的奉献》中的那个比喻上去，把它比拟为从``自然的''角度观察一个果园。

维诺格拉德已经为施德鲁写下了清楚的说明。我要在这里引用他在一篇文章中的话，这篇文章收在尚克和科尔比所编的书中：

在这个模型中隐含的一个基本观点是：所有语言的使用都可以被看成是一种激活听话人心中的过程的手段。我们可以把任何言语表达都想象成一个程序------这个程序可以间接地导致一组操作，后者在听话人的认知系统中被执行。这种``程序的编写''是间接的，这就是说我们是在和一个有智能的解释者打交道，他有可能采用一组和说话人的意图极其不同的行为。这些行为的准确形式取决于他关于世界的知识，以及他对于与他谈话的人的预期，等等。在本程序中，我们为这个发生在机器人心中的解释过程提供了一个简单的形式。机器人所解释的每句话都被转换成一组用Planner语言表示的指令。然后再执行这样构造出来的程序，以达到预想的效果。\footnote{}

Planner为问题分解提供了方便

这里提到的Planner语言是一种人工智能语言，它的基本特点是在其中设置了一些问题分解所需要的操作------比如说构造一棵包含子目标、子子目标等等的树的递归过程。这就是说，程序员不必每次不厌其烦地说明这棵树，它可以通过所谓``GOAL{[}目标{]}语句''自动地被建立起来。读Planner程序的人不会看到这样的操作被明确地提到，用行话来说，他们是``对用户透明的''。如果树中的一条路径没能导致想要达到的目标，那么这个Planner程序将``回溯''并尝试另一个途径。只要提到Planner，就必须涉及``回溯''这个咒语。

维诺格拉德的程序出色地利用了Planner的这些特点------更准确地说，他用的是Microplanner------Planner方案的一个部分实现。但是，在过去的几年中，以发展人工智能为己任的人们已经认识到，Planner中所用的那种自动回溯具有无法避免的缺点，可能无法引导他们到达预想的目标。因此他们就从这条道上退了回来，想去试试完成人工智能的其它途径了。

让我们再听一些维诺格拉德关于施德鲁的解释：

每个词的定义都是一个程序，它将在分析过程中的适当时刻被调用，而且能对所涉及的句子和当前的外界情境进行任意运算。\footnote{}

维诺格拉德所举的例子中有下面这个：

``那个''（英语的`the'）的不同可能意义对应于不同的过程，它们对上下文中的各种事实进行检查，然后规定一些行动，例如``在数据库中唯一地确定满足下列描述的对象''，或``断定所描述的对象在说话者看来肯定是唯一的''。在程序中包括了许多启发规则，它们被用来确定上下文中的哪些部分是有关的。\footnote{}

这个与``那个''这个词有关的问题的复杂程度是惊人的。保险的说法可能是：要想写一个程序来掌握最常用的五个词------汉语中是``的''、``了''、``是''、``一''、``不''；英语中是``the''、``of''、``and''、``a''、``to''------将等价于解决整个人工智能问题，因而也就相当于懂得了智能和意识到底是什么。稍微扯远一点，在汉语中五个最常用的名词------根据北京语言学院语言教学研究所编的《现代汉语频率词典》------是``上''、``人''、``里''、``年''、``天''（就是以这个次序）。关于这一点，令人惊奇的是大多数人想不到他们是这样用词的。如果你问问你的朋友，他们没准会猜想是``东西''、``问题''、``事情''、``同志''、``国家''这些词。另外，既然我们在考虑频率，非常值得一提的是在英语中最常用的12个字母------根据默根泰拉尔的统计------依次是``ETAOINSHRDLU''（这就是``伊她·娥英-施德鲁''的来历）。

施德鲁的一个有趣性质正好打击了那种认为计算机只善于处理数字的成见，这就是维诺格拉德所指出的一个事实：{``我们的系统不是以数字的形式接受数目的，而且只被教会了数到十。''}\footnote{}尽管施德鲁是以数学为基础的，但它本身对数学却一窍不通！像马姨一样，施德鲁对于支撑着它的低级层次一无所知。它的知识在很大程度上是过程型的（参见前面对话的第11节中``格拉德·维维诺诺博士''的解释）。

如果把以过程的方式嵌入在施德鲁中的知识和我那个句子生成程序中的知识作一个对比，那将是很有趣的。我那个程序中全部句法知识都以过程的形式嵌入在扩充迁移网中，由Algol语言书写，但语义知识------关于语义类属关系的信息------却是静态的：它被装在每个词后面的一个短数字表中。有那么一些词，例如动词``是''、``有''等等，是完全表示在Algol过程中的，但它们是例外。与此相反，在施德鲁中所有的词都被表示成程序。这种情况说明，尽管数据和程序在理论上是等价的，在实际应用中选这个不选那个仍会造成很大的差别。

句法和语义

这里再引一段维诺格拉德的话：

我们的程序在运行时不是先对一个句子做句法分析，然后再做语义分析，最后使用演绎推理生成答案。在理解一句话的过程中，这三个活动始终是同时进行的。当一段句法结构开始形成的时候，一个语义程序被调用，以检查这个结构是否有意义，而由此产生的答案可以为句法分析提供指导。在确定它是否有意义的过程中，这个语义子程序可以调用演绎过程，并提出一些关于现实世界的问题。以对话中的第34句为例（``把蓝方锥搁在盒子里的蓝方块上头有问题吗？''），语法分析器先是认为``把蓝方锥搁在盒子里''是一个动词短语，这时语义分析就开始了，即检查在数据库中是否提到过下面这个事件：蓝方块把蓝方锥搁在盒子里。由于找不到这种事件，语法分析器又把``把蓝方锥搁在盒子里的蓝方块上头''理解为一个动词短语，即``问能不能做这样的事情''\ldots\ldots 这样在不同类的分析之间就有着不间断的相互作用，即一个的结果会对另一个发生影响。\footnote{}

在自然语言中句法和语义如此紧密地纠缠在一起，这是十分有趣的。在前一章关于``形式''这个难以捉摸的概念的讨论之中，我们已经把这个概念分成了两个范畴：句法形式------这是可以用一个预知可终止的判定过程来检查的------和语义形式------这是无法做上述检查的。但在这里，维诺格拉德告诫我们------至少在使用``句法''和``语义''的通常意义时是如此------它们在自然语言中是彼此混在一起的。一个句子的外部形式，即它由基本符号所组成的结构，是无法清楚地划分成句法和语义两部分的。这对语言学来说是极其重要的一点。

下面是维诺格拉德对施德鲁所作的一些最后的说明：

\pandocbounded{\includegraphics[keepaspectratio]{./media/OEBPS/Images/img_118.png}}

图118．对``一个支撑着方锥的红立方体''的过程化表示

引自罗杰·尚克和肯尼思·科尔比，《思维的计算机模型和语言》，172页

让我们看一下该系统对``一个支撑着方锥的红立方体''这样一个简单的描述是怎样处理的。这个描述将使用下列概念：``方块''、``红的''、``方锥''和``等边的''------它们都是系统中所含有的关于世界的范畴。其结果可以被表示成像图118所示的那种流程图。注意这是一个程序，它要寻找满足该描述的对象。然后可以把这个程序结合在一个命令中，这个命令可以对这个对象进行某种处理，对它提出某个问题，或者，如果它出现在一个语句中，它将成为程序的一部分，该程序的生成是为了表示特定的意义，以备后来使用。注意，在事先说明了第一条指令------找到方块X1------只查找一个特定对象的情况下，这段程序也可被用来检查该对象是否满足这一描述。

这个程序刚一看上去似乎结构太复杂，因为我们不喜欢把一个简单语句想象成明显地含有循环、条件测试和其它的程序细节。解决的办法是提供一种内部语言，在其基本语句中包括适当的循环和测试，用这种语言来表示这一过程，就像前面那个描述一样简单了。图118所描述的程序可以被写成类似于下面所示的Planner程序：

\begin{longtable}[]{@{}ll@{}}
\toprule\noalign{}
\endhead
\bottomrule\noalign{}
\endlastfoot
(GOAL(IS ? X1 BLOCK)) & 目标：? X1是方形的 \\
(GOAL(COLOR-OF ? X1 RED)) & 目标：? X1是红色的 \\
(GOAL(EQUIDIMENSIONAL ? X1)) & 目标：? X1是等边的 \\
(GOAL(IS ? X2 PYRAMID)) & 目标：? X2是方锥体 \\
(GOAL(SUPPORT? X1 ? X2)) & 目标：? X1支撑着? X2 \\
\end{longtable}

流程图中的循环已隐含在Planner的回溯控制结构之中。这个描述的评价过程就是沿上述序列向下推进，直到某个目标失败为止，这时系统自动退回到前一个作出决定的地点，去尝试另一种可能性。一旦一个新对象名或``变量''（变量名前面都有个``?''）------如``? X1''或``? X2''------出现，就可作出一个决定。变量是由模式匹配程序来使用的。如果它们已经被指定代表某个特定的项，程序将检查该目标在这个项上是否成立。如不成立，它将检查所有可能满足该目标的项，检查的方法是先挑选一个，然后一旦在这个地方出现了回溯，再去选下一个。这样，甚至检测和选择之间的区别都变成隐含的了。\footnote{}

在这个程序的设计过程中，一个重大的战略决策是不把自然语言全部翻译成Lisp语言，而是只翻译一半------翻译成Planner语言。这样（因为Planner解释程序本身是用Lisp写的），一个新的中间层------Planner------就被插入在顶层语言（自然语言）和底层语言（机器语言）之间了。这样，根据自然语言生成了Planner程序并将其送到Planner解释程序之后，施德鲁的高层就可以被解放出来，去处理新的任务了。

这种决定会不断地涌现：一个系统应当有多少层？在哪一层上该设置哪种``智能''？设置多少？这几个是人工智能在今天所面临的最困难的问题中的一部分。由于我们对自然智能所知甚少，所以很难说清人工智能系统的哪一层应该完成一个任务的哪一部分。

希望以上内容能帮助你管窥一下本章前面的对话场景背后的东西。在下一章中，我们将看到人工智能中一些新鲜的、推测性的想法。

\begin{enumerate}
\item
  \phantomsection\label{Chapter18.xhtml_B_1}
\item
  \phantomsection\label{Chapter18.xhtml_B_2}
\item
  \phantomsection\label{Chapter18.xhtml_B_3}
\item
  \phantomsection\label{Chapter18.xhtml_B_4}
\item
  \phantomsection\label{Chapter18.xhtml_B_5}
\item
  \phantomsection\label{Chapter18.xhtml_B_6}
\item
  \phantomsection\label{Chapter18.xhtml_B_7}
\item
  \phantomsection\label{Chapter18.xhtml_B_8}
\item
  \phantomsection\label{Chapter18.xhtml_B_9}
\item
  \phantomsection\label{Chapter18.xhtml_B_10}
\item
  \phantomsection\label{Chapter18.xhtml_B_11}
\item
  \phantomsection\label{Chapter18.xhtml_B_12}
\item
  \phantomsection\label{Chapter18.xhtml_B_13}
\item
  \phantomsection\label{Chapter18.xhtml_B_14}
\item
  \phantomsection\label{Chapter18.xhtml_B_15}
\item
  \phantomsection\label{Chapter18.xhtml_B_16}
\item
  \phantomsection\label{Chapter18.xhtml_B_17}
\item
  \phantomsection\label{Chapter18.xhtml_B_18}
\item
  \phantomsection\label{Chapter18.xhtml_B_19}
\item
  \phantomsection\label{Chapter18.xhtml_B_20}
\item
  \phantomsection\label{Chapter18.xhtml_B_21}
\end{enumerate}

\phantomsection\label{Dialog19.xhtml}{}

\chapter{对实}

螃蟹邀请了几个朋友来看星期六下午电视中播放的橄榄球赛。阿基里斯已经到了，可还要等等乌龟和他的朋友树懒。

\dialogname{阿基里斯：}驭骑在那辆与众不同的独轮车上的两个人，会是我们那两位朋友吗？

（树懒和乌龟下了车走进来）

\dialogname{螃蟹：}哦，朋友们，你们光临我真高兴。请让我介绍一下我最爱戴的老朋友，树懒先生------这是阿基里斯。我敢肯定你认识乌龟。

\dialogname{树懒：}就我记忆所及，这是我有生以来第一次结识一位双目独眼龙。很高兴见到你，阿基里斯。关于双目独眼类的生物，我已经耳闻不少了。

\dialogname{阿基里斯：}我也是。我可以了解了解你那别致的交通工具吗？

\dialogname{乌龟：}你是问我们那辆双轴独轮车？谈不上什么别致。只不过是供两个人以同一速度从A到B的一种交通工具。

\dialogname{树懒：}它是由生产``双头翘''翘翘板的公司制造的。

\dialogname{阿基里斯：}噢，原来是这样。车上那个球形把手是什么东西？

\dialogname{树懒：}那是调挡装置。

\dialogname{阿基里斯：}哦！那它有多少种速度？

\dialogname{乌龟：}一种，包括回动。许多别的型号的比这更少，可这辆是特型的。

\dialogname{阿基里斯：}一看就是辆挺不错的双轴独轮车。噢，老蟹，我早就要告诉你，昨晚你的管弦乐队的演出我非常喜欢。

\dialogname{螃蟹：}谢谢，阿基你当时也在吗？小懒？

\dialogname{树懒：}不，我不巧没能到场。我那时正参加一场混合单打乒乓球锦标赛。由于我所在的队陷入了无对手的争夺冠军的决赛，所以特别扣人心弦。

\dialogname{阿基里斯：}你获得什么奖品了吗？

\dialogname{树懒：}当然------是一条一面镀金、一面镀银的铜制莫比乌斯带。

\dialogname{螃蟹：}祝贺你，小懒。

\dialogname{树懒：}谢谢，哎，跟我讲讲那个音乐会吧。

\dialogname{螃蟹：}那是场最美妙的演出。我们演奏了一些巴赫孪生兄弟的作品------

\dialogname{树懒：}是著名的约翰和塞巴斯第安吗？

\dialogname{螃蟹：}反正都是一回事。有一部作品让我想到了你，小懒------那是一首用两只左手弹的了不起的钢琴协奏曲。其中倒数第二个乐章（也是仅有的一个乐章）是一支单声部的赋格。你简直没法儿想象它有多复杂。我们还演奏了贝多芬的第九交响曲作为压轴节目。演出结束时，听众们全体起立，每个人都用一只手鼓掌，真是掌声雷动啊。

\dialogname{树懒：}哦，我没看上真是太可惜了。不过，不知你录了音没有？我家里有一台高保真的录音机------这是你所能买到的最好的双声道单喇叭的音响系统。

\dialogname{螃蟹：}我敢肯定你能搞到那录音。好啦，朋友们，球赛就要开始啦。

\dialogname{阿基里斯：}今天谁跟谁比赛，老蟹？

\dialogname{螃蟹：}一定是主队对客队。哦。不------那是上周。我想这周是本地队对外地队。

\dialogname{阿基里斯：}我为本地队加油。我一向这么做。

\dialogname{树懒：}啊，那你可太地方主义了。我从不为本地队加油。越是远在十万八千里的队，我越是要给他们加油。

\dialogname{阿基里斯：}噢，那么你是住在十万八千里啦？我听说那儿很迷人，不过我还是不想去那儿。毕竟太遥远了。

\dialogname{树懒：}而且奇怪的是，你无论走哪条路，都不会更近些。

\dialogname{乌龟：}那是种对我路子的地方。

\dialogname{螃蟹：}球赛就要开始了。我看得打开电视了。

（他走到一个有屏幕的巨大的箱子前，箱子下面有个仪表盘，复杂得跟喷气式飞机上的一样。他轻轻按了一下一个球形开关，屏幕上出现了色彩鲜亮、图像清晰的橄榄球场。）

\dialogname{解说员：}下午好，球迷们。一年一度的本地队和外地队在绿茵场上的较量又拉开战幕了。今天下午断断续续地下了点小雨，场地多少有点湿。不过，尽管天公不作美，这场比赛肯定会非常精彩，特别是本地队有两名著名的八分位，他们是提里趿拉和回文斯。大家知道，根据规则，美式足球------也就是这种橄榄球------每个球队里应该只有一名四分位指挥进攻，不过只要比赛精彩，弄成两名八分位也没有关系，不是吗？

\dialogname{乌龟：}天哪！这种比赛我还从没见识过。

\dialogname{解说员：}现在丕里帕拉为本地队开球。球飞在空中！外地队的法兰孙接到了皮球，他抱球往回跑------跑到了20码线的地方，到了25码线，到30码线，他在32码的地方摔倒了。是本地队的哞呕一个擒抱，把他摔倒的。

\dialogname{螃蟹：}好漂亮的反攻！你们看到没有，他差一点叫悟诡给擒抱住------可他还是挣脱掉了。

\dialogname{乌龟：}叫谁擒抱？

\dialogname{螃蟹：}悟诡！

\dialogname{乌龟：}怎么跟我的名字就差一点儿？

\dialogname{树懒：}别犯傻了，老蟹。没这么回事，悟诡没有擒抱法兰孙。没有必要用这种``差一点''发生的事儿蒙可怜的阿基里斯（和我们别人）。事实就是事实------没有什么``差一点''啦、``要是''啦、``还''啦、``可是''啦这些东西。

\dialogname{解说员：}这是重放的镜头。注意79号队员悟诡，他从边路包抄过来，威胁法兰孙，打算擒抱他！

\dialogname{树懒：}``打算''！哼！

\dialogname{阿基里斯：}这一招儿可太漂亮啦！要是没有重放镜头我们怎么能看得到！

\dialogname{解说员：}外地队第一次进攻，攻进十码。现在球在闹得乐手中，他把皮球传给了牛韦------这是个后传球------牛韦绕到右面，把球掷给了法兰孙------这是个双后传球，诸位！------现在法兰孙把球递给华果树，华在争球线后面十二码的地方达阵得分。这个三后传球丢了十二码！

\dialogname{树懒：}太棒了！真是一场扣人心弦的比赛！

\dialogname{阿基里斯：}哦，小懒，我还以为你一直在为外地队加油呢。这场球他们丢了十二码。

\dialogname{树懒：}他们丢了十二码？啊------只要他们表演得出色，管它输还是赢呢！我们再看一遍吧。

（\ldots\ldots 就这样，上半场比赛结束了。在第三节比赛就要结束的时候，本地队决定性的时刻来到了。这时他们落后八分。这是他们第三次进攻，还一码未进，他们这时急需继续进攻的机会）

\dialogname{解说员：}球被抛到提里趿拉手里，他退后几步，找人传球，晃过了悟诡。这时回文斯在最右边，附近没有别的队员。提里趿拉看到了他，来了一个低传，把球传到回文斯手里。回文斯接住球以后，就------\kaitismall{（可以听到人群中发出的惋惜声。）}------唉，他跨出了边线！各位观众，这对本地队可是个沉重打击。要是回文斯没有跨出边线，他会一直跑到底线区，达阵得分的。现在让我们看一看虚拟的重放镜头。

（屏幕上出现了同刚才一样的阵容。）

球被抛到提里趿拉手里，他退后几步，找人传球，晃过了悟诡。这时回文斯在最右边，附近没有别的队员。提里趿拉看到了他，来了一个低传，把球传到回文斯手里。回文斯接住球以后，就------\kaitismall{（可以听到人群中发出的紧张的惊叹声。）}------他差一点跨出边线！不过他仍在界内，通向得分区的路上完全没有对方队员。回文斯一往直前，达阵得分！\kaitismall{（球场内欢声雷动。）}啊，球迷们，要是回文斯没有跨出边线，这就是将会发生的情形。

\dialogname{阿基里斯：}等等\ldots\ldots 达阵得分这事有还是没有？

\dialogname{螃蟹：}哦，没有。这只是虚拟的重放镜头。这只不过是演示了一下假设而已。

\dialogname{树懒：}这是我听到过的最荒唐的事。我看他们还会发明水泥护耳皮套呢。

\dialogname{乌龟：}虚拟的重放镜头有点少见，是吧？

\dialogname{螃蟹：}要是你有一台虚拟电视，也没有什么特别的。

\dialogname{阿基里斯：}需你电视？是一种非得你亲自操作不可的电视吗？

\dialogname{螃蟹：}不，你理解岔了。这是种可以进入虚拟状态的新型电视，特别适用于观看球赛之类的节目。这是我刚搞到的。

\dialogname{阿基里斯：}它上面为什么有这么多的球形开关和叫人眼花缭乱的仪表盘？

\dialogname{螃蟹：}你可以用它们来调你所需要的频道。有很多频道都在广播虚拟的节目，你能很方便地从中选择。

\dialogname{阿基里斯：}你能把你说的表演给我们看看吗？恐怕我还不太明白``广播虚拟的节目''是什么意思。

\dialogname{螃蟹：}哦，其实很简单。你自己也能弄明白。我现在去厨房做点油炸土豆片，我知道小懒很爱吃。

\dialogname{树懒：}噢！去吧，老蟹！油炸土豆片是最对我口味的食品。

\dialogname{螃蟹：}你们别的人呢？

\dialogname{乌龟：}我也想吃点。

\dialogname{阿基里斯：}我也想。可等等------在你去厨房之前，能吿诉我使用你的虚拟电视有什么窍门吗？

\dialogname{螃蟹：}没什么特别之处，尽管接着往下看，什么时候有某种几乎要成却没成的事，什么时候你希望某件事情不是这样发生，而是换一个样子，你就摆弄摆弄那些开关，看看会发生什么事。你也许会调出些稀奇古怪的频道来，但你不会把它们弄坏的。

（他进了厨房）

\dialogname{阿基里斯：}我不明白他是什么意思。嗯，好吧，我们还是回到这场比赛上来吧，我真叫它迷住了。

\dialogname{解说员：}这是外地队第四次进攻，本地队为守方。外地队采取踢悬空球的战术，提里趿拉位置靠近得分线。牛韦将球踢回------这是个高飘球，皮球朝着提里趿拉落下去------

\dialogname{阿基里斯：}接球，提里趿拉！给外地队一点颜色看看。

\dialogname{解说员：}------球落到了水坑里------噗哧！球反弹的方向很奇怪！斯普鲁克猛地一个飞身接球！反弹起的皮球几乎擦着提里趿拉的身边滑开了------这个球被判为失球。裁判判凶猛的斯普鲁克为外地队得分，本地队落后七码。这对本地队是很不利的。这只能说是老天不想成全他们。

\dialogname{阿基里斯：}噢，不！要是没有这场雨\ldots\ldots{}\kaitismall{（绝望地绞着双手。）}

\dialogname{树懒：}又一个该死的假设！你们怎么都陷进自己幻想的荒诞世界里去了！我要是你们，我就不会违背现实。我的格言是``避免虚拟的谵语''。即使别人给我一百块------不，一百一十二块------油炸土豆片，我也不会放弃我的格言。

\dialogname{阿基里斯：}哎，这倒提醒了我。没准儿摆弄摆弄这些球形开关，就能碰巧找到不下雨的虚拟重放镜头，那样一来就没有泥水坑了，皮球也不会向稀奇古怪的方向弹了，提里趿拉也就不会失球。我想知道\ldots\ldots{}\kaitismall{（走到虚拟电视机前，盯着它瞧了一会儿。）}可我一点也不知道这些球形开关都是管什么的。\kaitismall{（随便拧了几下。）}

\dialogname{解说员：}这是外地队第四次进攻，本地队为守方。外地队采取踢悬空球的战术，提里趿拉位置靠近得分线。牛韦将球踢回------这是个高飘球，皮球朝着提里趿拉落下去------

\dialogname{阿基里斯：}接球，提里趿拉！给外地队一点颜色看看。

\dialogname{解说员：}球落到了水坑里------噗哧！球反弹到他手里。斯普鲁克猛地从他身后飞起抢球，却被他有效地阻挡住了，他避开凶猛的斯普鲁京，在他前面是一片无人之境。注意，他到了50码，到了40码，30码了，到20码，10码------达阵得分，本地队得分！\kaitismall{（本地队一方欢声如雷）}球迷们，要是橄榄球不是橄榄形的，而是西瓜形的，这就是将会出现的情况！不过，在现实中，本地队失了球，而外地队攻到了本地队的七码线上。这真是皮球无知，不解人心。

\dialogname{阿基里斯：}这你怎么想，小懒？

（阿基斯冲着树懒得意洋洋地笑了一声，可后者却完全没有注意，他正全神贯注地瞧着螃蟹端着一只大盘子走了进来，那盘子上盛着一百一十二块------不------一百块又大又馋人的油炸土豆片，以及为每个人准备的餐巾。）

\dialogname{螃蟹：}你们觉得我的虚拟电视怎么样？

\dialogname{树懒：}坦率地说，老蟹，真叫人失望。它好像出了毛病。至少有一半时间里它是在瞎捣乱。这要是我的，我会马上不要了，把它给一个像你这样的人，不过，这当然不是我的电视。

\dialogname{阿基里斯：}这是件很古怪的玩艺儿。我想让它重演在另一种天气下的比赛，可这玩艺儿似乎自有一套想法！它不去改变天气，而是把橄榄球变得不是橄榄形的，却是西瓜形的！那么请告诉我------橄榄球怎么能不是橄榄形的？这真是自相矛盾。太荒诞了！

\dialogname{螃蟹：}真没劲！我还以为你们会发现什么有趣的虚拟呢。要是把橄榄球赛改成棒球赛，刚才那场比赛的结局会怎么样，你们想看看吗？

\dialogname{乌龟：}噢，这主意太伟大啦！

（螃蟹拧了拧两个球形开关，然后退了几步。）

\dialogname{解说员：}现在有四个人出局，还------

\dialogname{阿基里斯：}四人出局？

\dialogname{解说员：}对，球迷们------四人出局。你一旦把橄榄球赛改成棒球赛，就得接受一些既成事实！我还是要说，四人出局，外地队是守队，本地队进攻。现在提里趿拉打垒。外地队采取触击战术。牛韦现在举手投球------他投了一个上升球。球朝着提里趿拉直飞过来------

\dialogname{阿基里斯：}狠扣它一个，提里趿拉！给外地队一点颜色看看！

\dialogname{解说员：}------但是球好像沾了水，在空中划了一道奇怪的弧线。斯普鲁克猛地飞身抱球！那球看去险些擦过提里趿拉的球棒，从他身边弹开了------这被判为飞球。裁判判凶猛的斯普鲁克为外地队得分，从而结束了第七局的比赛。这对本地队很不利。要是把橄榄球赛改成棒球赛，这就是刚才的比赛将出现的结局。

\dialogname{树懒：}哼！你们干脆把比赛挪到月球上去算啦。

\dialogname{螃蟹：}说干就干！只要在这儿拧一下，在那儿拧一下\ldots\ldots{}

（屏幕上出现了一块孤零零的火山口样的场地，双方队员都穿着太空服定在原地不动。但很快，两队队员同时动起来，队员们跳得很高，有时都高过了其他队员的头顶。球被抛向空中，高到几乎看不见了。然后再缓慢地落到某个身穿太空服的队员手里，球落下的地方离它被抛起的地方几乎有四分之一公里远）

\dialogname{解说员：}朋友们，你们现在看到的是虚拟的比赛在月球上进行时的情景。现在让我们先抽空儿告诉大家一点重要的商业行情。提供这一信息的好心人专酿古鲁皮啤酒------这是最对我口味的啤酒------待会儿我们再回到球赛上去。

\dialogname{树懒：}要是我不这么懒的话，我一定要亲自把这台破电视奉还给卖主！可是，唉，我已经命中注定要做一个懒惰的树懒了\ldots\ldots{}

（嚼了一大口油炸土豆片。）

\dialogname{乌龟：}嘿，这可是个了不起的设想，老蟹。我可以提出一种假设情况吗？

\dialogname{螃蟹：}当然可以。

\dialogname{乌龟：}要是空间变成四维的，刚才那场比赛会怎么样？

\dialogname{螃蟹：}哦，这可够复杂的，龟兄，不过我看我还是能够把这种构想输进仪器里面。请稍候。

（他走上前去，看来是第一次充分使用了那台虚拟电视机的控制盘，几乎每只球形开关都被拧了两三次。他仔细地检查了各种控制器的数据，然后退回来，脸上带着一副满意的神情。）

我看这下行了。

\dialogname{解说员：}现在让我们来看看虚拟的重放镜头。

（屏幕上出现了许多混乱地绞在一起的管。一会儿变大，一会儿变小，有时又似乎在旋转。随后又变成某种奇特的蘑菇状的东西，到最后又回到管的形状。就在那些图像一会儿这样、一会儿那样，变成各种稀奇古怪的图形时，解说员又开始解说了。）

提里趿拉退后几步。他看到了场外十码处的回文斯，把球向右向外抛给了回文斯------看来很成功！回文斯现在到了35码平面，40码，他在场地上自己一方43码的平面上被擒抱了，三维的球迷们，要是球赛在四维空间里举行，这就是将会发生的情况。

\dialogname{阿基里斯：}你对它做了些什么，老蟹，你是怎么摆弄控制盘上那些各式各样的开关的？

\dialogname{螃蟹：}我选择了适当的虚拟频道。喏，有这么多的虚拟频道同时广播，我只是把它调到广播我们所要求的那种虚拟的频道上。

\dialogname{阿基里斯：}随便什么电视机都能这么来吗？

\dialogname{螃蟹：}不，大多数电视机无法接收虚拟频道。这需要一套特殊的、极难制作的电路。

\dialogname{树懒：}你怎么知道哪一个频道播放哪一种节目？看报纸上的电视节目预告吗？

\dialogname{螃蟹：}我不需知道每种频道的呼叫代号，而是用编码的方法在这些开关上调出我想要的虚拟情景。用行话讲，这就叫作``通过对实参数调用某一频道''。总是有大量频道在播放各种可能设想的情景。那些内容``相近似''的各频道，其呼叫代号也彼此相近似。

\dialogname{乌龟：}为什么刚才在我们第一次看虚拟的重放镜头时，你不用动那些开关呢？

\dialogname{螃蟹：}这是因为我只是把它调到一个跟事实频道很相近的频道中，跟事实相差很小。所以它只是同事实偶有不同。把它完全调到事实频道几乎是不可能的。不过这没关系，因为事实总是很乏味的。所有的重放镜头都是事实的本来面目，这你能想象吗？多没劲！

\dialogname{树懒：}可我觉得虚拟电视这种构想整个儿就没劲。不过，要是有证据说明你这台机器可以搞出点有趣的对实节目的话，我没准儿会改变对它的看法。比方说，要是加法不满足交换律，刚才那场比赛会怎么样？

\dialogname{螃蟹：}噢，老天，噢，老天爷！我担心对于这种型号的机器来说，这个要求所需要的推理过程太复杂了。无奈我现在还没有一台在这方面很完备的超拟电视机。有了这样一台超拟电视机，就可以满足你各种花样的要求了。

\dialogname{树懒：}嗬。

\dialogname{螃蟹：}不过请看------这差不多也能做到。你们想不想瞧瞧要是13不是个素数，刚才那场比赛会怎么样？

\dialogname{树懒：}多谢！不必了。这毫无意义！我要是刚才那场比赛的话，一定会叫你们这样随心所欲地播来放去弄得厌烦透顶啦。哼，你们这群头脑不清的家伙在这里没完没了地改换概念！接着往下看球赛吧！

\dialogname{阿基里斯：}你从哪儿搞来的这么一套虚拟电视机，老蟹？

\dialogname{螃蟹：}信不信由你：小懒和我有天晚上去了趟乡村市场，这电视机就是我们在一次抽彩中得的头奖。一般地说，对这种无聊玩艺儿我一向不感兴趣，可当时我被某种疯狂的冲动抓住了，于是就买了一张彩票。

\dialogname{阿基里斯：}你呢，小懒？

\dialogname{树懒：}我承认我也买了一张，只不过是为了迁就迁就老蟹。

\dialogname{螃蟹：}等宣布了中奖号码，我才大吃一惊，我中彩啦！

\dialogname{阿基里斯：}不可思议！在我认识的人里，还从没有谁中过彩！

\dialogname{螃蟹：}当时我也是目瞪口呆，觉得难以置信。

\dialogname{树懒：}关于中彩的事，你还有什么别的话要讲吗，老蟹？

\dialogname{螃蟹：}没有了。哦------我那张彩票号码是129。他们宣布的中彩号码是128，只差一点。

\dialogname{树懒：}所以，你们明白了吧，他事实上根本没中，

\dialogname{阿基里斯：}不过他差一点中了\ldots\ldots{}

\dialogname{螃蟹：}我更倾向于说我中了，因为我是那么接近\ldots\ldots 要是我的号码再减去个一，我就会中了。

\dialogname{树懒：}但可惜的是，老蟹，在这种情况下``差一点''和``差十万八千里''没什么两样儿。

\dialogname{乌龟：}没错。哎，你呢，小懒？你的号码是多少？

\dialogname{树懒：}我的号码是256------是128的两倍。太接近了，是吧！可我不明白那些市场官员的脑袋为什么死不开窍。他们拒绝给我完全应该归我的奖品。有个家伙声称他应该得这个奖，因为他的号码是128。我觉得我的号码比他的更接近，可对那些官僚们你真是一点辙都没有。

\dialogname{阿基里斯：}我整个儿一个糊涂，要是你根本就没有赢得虚拟电视机，老蟹，我们怎么能今天一下午坐在这儿看这台电视？看来我们自己好像也是呆在某个可能会存在的虚拟世界里，只要环境稍有些变化\ldots\ldots{}

\dialogname{解说员：}各位观众，要是螃蟹赢得了虚拟电视机的话，这就是这个下午在螃蟹家将会发生的事情。但是因为他并没有中彩，所以这四个朋友事实上只是观看了一场本地队被打得落花流水的比赛，用这种方式度过了一个愉快的下午。本地队在这场比赛中以0比128失利------也许是0比256？哦，对于一场在五维空间中的冥王星上举行的西瓜球比赛，这是无关紧要的。

\phantomsection\label{Chapter19.xhtml}{}

第十九章

\chapter{人工智能：展望}

``差一点儿出现''的情景与虚拟情景

我的一个朋友读过《对实》之后对我说：``我叔叔差一点儿就是美国总统了！''\,``真的？''我问。``敢情，''他答道，``他当时是108号鱼雷艇艇长。''（约翰·肯尼迪曾任109号鱼雷艇艇长。）

这就是《对实》的全部意义所在。在日常思维中，从我们面临的情况、我们拥有的想法、客观发生的事件之中，我们总会构造出一些精神上的``变奏曲''，让某些性质基本保持原样，而让另一些性质发生``滑动''，从而形成一些``虚拟的''------我叫它们``对实的''情景。哪些性质是我们允许滑动的？哪些是我们根本不会考虑让其滑动的？在某种较深的直觉层次上，哪些事件被感知成是与确实发生的事件紧密相关的？有哪些事情会被我们看作``差一点儿''发生，或者``本当''发生了，但它们事实上没有发生？当我们听人讲故事时，哪些可能的结局会自然而然地出现在我们的脑海里？为什么我们会觉得某些虚拟事件比其它一些``虚拟得少点''？归根结底，没发生的事件显然就是没发生。在这里不存在一个``未发生度''。对``差一点儿''出现的情况也是一样。有时某个人会伤心地说``它差一点儿就发生了''，而另一个时候他又会欣慰地说同一句话。但这种``差一点儿''是存在于思想中的，并非存在于外界事实里。

假设你驾车行驶在乡间公路上，迎面碰上一群蜜蜂。你不仅会及时注意到它们，而且整个情景会立即在脑海中产生一群``重放的镜头''。典型地，你会想``幸亏我没开车窗''------或不幸是其反面：``糟糕，车窗没关！''；还可能是``幸亏我没有骑自行车来！''；``我怎么没早来五秒钟呢？！''；另一些奇特但可能出现的想法是``如果这是一头牛，我就已经完蛋了！''；``我敢打赌，这些蜜蜂更愿意碰上一棵玫瑰。''；甚至还有更怪的念头：``这些蜜蜂要都是钞票就好！''；``幸亏这些蜜蜂不是水泥的！''；``倒霉，为什么蜜蜂是一群而不是一只呢？''；``我是我，而不是蜜蜂，这可是万幸！''哪些滑动是自然的？哪些不是？为什么？

在最近的一期《纽约客》杂志上，重新引用了《费城迎客报》中的一段话：\footnote{}

如果列奥纳多·达·芬奇是个女的，西斯廷教堂的天花板就不会被画上画了。

《纽约客》评论道：

如果米开朗基罗是个连体双胞胎，那他完成这幅作品的时间就会减少一半。

《纽约客》评论的要点不是说这种虚拟事件是``假的''，而是说抱这种念头------想改动一个特定人物的性别或数目------的人，未免有点不近情理。但具有讽刺意味的是，在同一期的一篇新书介绍中，下面的句子被毫不脸红地印出来了：

我想他（非利鲁·弗兰克教授）如果仍健在，一定会对这两本书表示非常欣赏的。\footnote{}

可怜的弗兰克教授已不在人世了。一个人读到写于他去世之后的书，这无论如何是说不通的。那为什么没人对这个一本正经的说法进行嘲弄呢？在某种难以捉摸的意义下，这个例子中参数的改动不像前一个例子那样违反我们关于``可能性''的观念。有某种东西使得我们在这个例子里面比前一个例子更容易想象``其它一切都是照常不变的''。但这是为什么呢？我们要怎样对事件和人物进行分类，才能进一步知道哪些滑动是``有意义的''，哪些是``胡闹''呢？

不妨看看在下列表达之间进行滑动是多么自然吧：开始是一个没什么价值的陈述：``我不懂俄文''，而后是一个内容丰富一些的条件句``我要懂俄文就好了''，下一个是一个充满情感的主观想象：``我希望我已经懂俄文了''，最后是一个内容丰富的虚拟语句``如果我懂俄文，我就能读契科夫和莱蒙托夫的原文著作了。''如果心智在一个否定句中除了晦暗的障碍以外什么都看不到，那该是多么平淡乏味啊！一个活生生的心智总能发现通往充满了可能性的世界的窗口。

我相信，那种``差一点儿''发生的情况和潜意识中构造的主观想象，代表着洞察力的某种最丰富的潜在来源。靠这种洞察力，我们可以了解人类是怎样组织他们对世界的感知，并形成范畴的。这个观点的一个强有力的支持者是语言学家兼翻译家乔治·斯坦纳。他在他的著作《巴别塔之后》中写道：

假设、想象、条件、虚拟的句法和偶然的句法可能是人类语言的发生中心\ldots\ldots{}{[}它们{]}不光是造成哲学和语言学中的一些麻烦，与它们有关的也不像有人感觉的那样仅限于将来时态。或许该把它们划归到更大的集合中，如``假设推测''或``可能的选择''。这些带``如果''的命题为人类的情感机制提供了基础\ldots\ldots{}

我们能够------同时也需要------对世界进行否定或``不定''，以便以别的方式来设想和表述\ldots\ldots 我们需要找一个词来描述语言的这种威力和这种意向：即谈论``其它情况''的能力\ldots\ldots 或许该叫它``替代''，以此来定义``不同于真实的情况''、虚拟命题、想象、愿望和遁词，我们用这些东西掌管着我们的头脑，并以它们为工具为我们自己构造了一个变动不居、充满幻想的生物环境和社会环境\ldots\ldots{}

最后，斯坦纳为虚拟事件唱了一首虚拟颂歌：

从我们已有的知识推测，假如没有这种幻想的、虚拟的、反决定论的语言工具，没有产生并存储在大脑皮质的``多余''部分中的语义能力，并以此来想象和表达那些在单调沉闷的生物性衰退和死亡之外的各种可能性，人类是几乎不可能生存至今的。\footnote{}

这种``主观想象出的世界''的构造过程产生得如此随便、自然，以致于我们几乎没注意到自己在做些什么。我们从我们的幻想中选出一个世界，它在某种内在精神意义下接近于真实世界。我们把真实的东西和我们感觉``差不多''真实的东西相互比较。这样，我们就得到了某种难以捉摸的对实在的看法。对话中的树懒就是实在的变奏之中的一个滑稽可笑的例子------一种没有进入主观想象能力的智能动物（或至少它自己宣称没有这种能力------但你可能已经注意到，在它的话中充满了虚拟语句！）试想，如果我们没有这种滑出真实世界、进入``假如''世界之中的创造性能力，那我们的精神生活将是何等的贫乏啊！从人类思维过程研究的观点看，这种滑动是十分吸引人的，因为它大多发生于完全缺乏意识指引的情况下，这意味着对什么可以滑动、什么不能滑动的观察将为我们认识潜意识提供一个良好的窗口。

认识这种精神测度的性质的一种方式是``以其人之道，还治其人之身''。在对话中就是这样干的，即要求我们的``主观想象力''设想一个世界，其中和我们的预期相比恰好让主观想象力这个概念滑动了。在对话中，第一个虚拟的重放镜头------即回文斯没有出界------是一个非常自然的想象。事实上，这个想法来自一个人对我说过的一番平常而又凑巧的话，这个人在一场球赛中坐在我身旁。出于某种原因，他的话触动了我，使我思考是什么使得这种滑动如此自然，而另一些滑动------如改变比分------则不然。从这些想法中，我进一步考虑另一些或许是更难滑动的特征，例如天气（这反映在对话中）、比赛的种类（也在对话中）、甚至一些更离奇的变化（也在对话中）。我注意到，在某一情景中显得荒唐可笑的滑动在另一个情景中可能就很容易发生。例如，有时你会很自然地考虑如果球的形状变了会怎么样（如你在玩一个气没打足的篮球时），而另一些时候这个念头则绝不会进入你的脑海（如你在看一场橄榄球赛的电视转播时）。

稳定性的分层

当时我觉得（现在依然如此），某些事件（或环境）特征的可滑动性，取决于我们在一个什么样的嵌套环境集合中感知此事件（或环境）的发生。在这里，从数学中借用``常量''、``参量''和``变量''这几个词可能是有用的。数学家、物理学家和其他一些人在进行运算时常常会说``c是个常量，p是个参量，v是个变量''。他们的意思是这三者都是可变的（包括常量），但是，存在一种可变性的层次差别。在这些符号所表示的情景中，c构成一种全局条件，p构成一种全局性较弱的条件------它可以在c保持固定时变动，最后，v可以在c和p保持固定时变来变去。在v固定时设想让c和p变化，这是没有意义的，因为c和p构成了使v有意义的环境。例如，设想一个牙科医生，他有一排病人，而每个病人又都有一排牙齿。一个完全有意义的事是让病人保持固定而让他的牙齿变动------但让牙齿保持固定而让病人变动则是根本没有意义的。（虽然有时变动牙科医生还是挺有意义的\ldots\ldots）

我们逐层构造了一个情景的心智表示。其中最低层构成了环境的最深刻的方面------有时这个层次低得根本无法变动。例如，我们世界的三维性是个非常根深蒂固的观念，我们中的大多数人都从没想过让它在头脑中滑动。它是一个恒定的常量。然后是这样一些层次：它们虽然不是永久固定的，但却构成了情景中暂时固定的方面。这些方面可以被称为背景假设------在心灵深处，你知道这些方面是可能改变的，但在大多数时候，你总是不假思索地把它们当作不变的方面。这些仍可以被称为``常量''。例如，当你去看橄榄球赛的时候，比赛规则就属于这种常量。再往后就是``参量''：你认为它们有更大的可变性，但暂时把它们当作常量。在一场橄榄球赛中，参量可能包括天气、参赛球队，如此等等。可能------或很可能------存在着多层参量。最后，我们到达了你关于情景的心智表示中最不稳定的那些方面------变量。这就是像回文斯的出界之类的东西，它们在脑海中很``散漫''，你不会介意让它们在短时间内偏离其真实值。

框架和嵌套的环境

``框架''这个词正流行于当今的人工智能研究中。它可以被定义为``环境的可计算的实例化''。这个词的普及应归功于马尔文·明斯基，他发表了不少关于框架的见解。在框架语言中，可以说情景的心智表示包含着彼此嵌套的框架。情景中的各种成分都有它自己的框架。我心里对嵌套的框架有个想象，用语言表达出来可能会是很有意思的。设想有若干个柜子，每一个柜子上可以装若干个抽屉（就像中药铺里的那种柜子）。一旦你选择了一个柜子，你就有了一个框架，而且在这个柜子上的装抽屉的洞中可以放进``子框架''。子框架本身又是一个这样的柜子。怎么可能把整个柜子插到另一个柜子中的一个抽屉槽里去呢？很容易：你可以把它压缩一下，使它变形，因为这毕竟是在想象中，不是真要这样干。在外面的这个框架中，可能有若干不同的抽屉槽需要填充，然后你可能要用抽屉（或子框架）来填充某些装在里面的柜子。这个过程可以递归地进行下去。

对一个柜子进行挤压和扭曲，使得它能被装进任何形状的槽中，这种生动的超现实主义想象或许是极其重要的，因为这说明当你迫使某概念进入一个环境时，它也被该环境所挤压和扭曲了。这样，如果你考虑的人是个橄榄球队员，那么此时你关于``人''的概念变成了什么样的呢？显然它已经被扭曲了，是被整个环境强加给你的。你已经把``人''这个框架插在框架``橄榄球赛''的一个槽里了。用框架进行知识表示的理论依赖于这样一个想法：世界是由半封闭的子系统构成的，其中每个子系统都能充当其它子系统的环境，并且在此过程中不会受到或造成太大的破坏。

关于框架的主要想法之一，就是让每个框架带有自己的一组预设值。与此相应的想象是给每个柜子的每个抽屉槽中都预先装上一个称为``缺席选择''的抽屉，但并不将其装紧。如果我对你说``画一段河岸''，你将会激发出一幅具有各种特征的视觉想象，而如果我再添加进一步的说明，例如``在干旱季节''、``长满垂柳''、或``惊涛拍岸''，则你可能会去掉那些特征中的大多数。槽中缺席值的存在保证了递归填充槽的过程具有终点。实际上你是在说``我自己将只填充外面的三层槽，其它的就用缺席项''。和缺席预设值并存，一个框架还包含关于其使用限制的知识，以及在超出其容许限度时转换到其它框架的启发式规则。

框架的嵌套结构使得你可以``拍摄特写镜头''------以足够近的距离观察对象的细节：你只需选择适当的子框架，然后是子框架的子框架，等等，一直推进到拥有你所需要的细节为止。这就像一本交通地图册，第一页是一张全国地图，然后是分省地图，如果你需要进一步的细节，甚至还会有城市的地图或某些大村镇的地图。你可以设想一本拥有任意详细程度的地图册，一直到一条街、一幢房子，一间屋子，等等。这就像使用一台配有放大倍数不同的透镜的电子显微镜一样，每片透镜都有它自己的用途。重要的是我们能利用所有这些不同的尺度，因为细节常常是不相干的，甚至只能起干扰作用。

由于不同的框架可以被任意地插进其它框架的槽中，这就有发生冲突或``碰撞''的潜在可能。那种用单一的、全局的观点对``常量''、``参量''、``变量''进行分层，以形成一个整齐划一的模式的想法，未免把问题过分简单化了。事实上，每个框架自己都有一个可变性的层次结构。正因为如此，分析我们如何感知一个复杂事件------如一场橄榄球赛，连同它的大量子框架、子子框架等等------才成为一种混乱得令人难以置信的工作。这一大堆框架是怎样相互作用的？如果某个框架说：``这个项是个常量''，而另一个框架说：``不，它是个变量！''，那这种冲突怎样才能得到解决呢？这是一些关于框架理论的更深更难的问题，对此我无法提供答案。关于框架到底是什么东西，以及如何在人工智能程序中实现它，至今仍没有一个完全一致的看法。我准备在下面几节中讨论其中的一些问题，并谈谈视觉模式识别中的一些谜题，我称之为``邦加德问题''。

邦加德问题

``邦加德问题''的一般形式是俄罗斯科学家莫·邦加德在其著作《模式识别》中给出的。一个典型的邦加德问题------在他的一百个问题中的第51号------见图119。

\pandocbounded{\includegraphics[keepaspectratio]{./media/OEBPS/Images/img_119.png}}

图119．邦加德问题51号

摘自莫·邦加德《模式识别》

这些有趣的问题是为模式识别者准备的，不管是人还是机器（也可以扔一个给外星人）。每个问题包括十二幅用方框围住的图像（以后简称为``框''）：六个在左面，形成第Ⅰ组，六个在右面，形成第Ⅱ组。这些框可按下列方式编号：

Ⅰ-A　Ⅰ-B　　　Ⅱ-A　Ⅱ-B

Ⅰ-C　Ⅰ-D　　　Ⅱ-C　Ⅱ-D

Ⅰ-E　Ⅰ-F　　　Ⅱ-E　Ⅱ-F

问题是``第Ⅰ组和第Ⅱ组中的框有何不同？''

一个求解邦加德问题的程序可能会包括若干阶段，原始数据通过它得以逐步转换成描述。初级阶段相对来说灵活性不大，而高级阶段逐渐变得更灵活一些。最终的阶段具有一种我称之为``尝试性''的性质，就是说一幅图像的表达方式总是带有尝试性的。只须出现微小的变化，用后面几个阶段提供的手段就能重新构造出一个高层描述来。下面要提出的想法本身也具有一定的尝试性。我将设法先传达总体想法，回避巨大的困难。然后我将回过头来设法解释那些微妙复杂之处。这样你在阅读时对于整个工作方式的认识也要经过某种修正。但这是体现在讨论的精神之中的。

通过预处理选择微词汇表

假设我们要解某个邦加德问题。题目被送到一台电视摄像机前，这样原始数据就被读进去了。然后对原始数据进行``预处理''，这就是说某些显著特征可以被检测出来。这些特征的``名字''构成了该问题的一个``微词汇表''，而它们是从一个通用的``显著特征词汇表''中抽取出来的。这个显著特征词汇表中的一些典型项是：

线段、曲线、水平的、垂直的、黑的、白的、大的、小的、尖的、圆的\ldots\ldots{}

在预处理的第二阶段，使用了某些关于``基本图形''的知识。如果发现了基本图形，它们的名字也要被包括进来。这样，可能会选出下列的项：

三角形、圆形、矩形、凹入，凸出、直角、顶点、尖端、箭头\ldots\ldots{}

在人脑中，这里差不多就是意识和潜意识的交汇点。以上的讨论主要是为了对此后发生的过程进行描述。

高层描述

至此，这幅图像已经根据某些已知的概念，在一定程度上被``理解''了，而且完成了一些检查工作。对于十二个框中的某一个或几个，已经作出了尝试性的描述。在这种描述中一般只使用下面这些简单的描述词：

在上、在下、在左、在右、在里、在外、靠近、远离、平行、垂直、在一行中，分散的、间隔均匀的、间隔参差不一的、等等。

同时，还可能用到确定的或不确的数量描述词：

1、2、3、4、5\ldots\ldots 许多、很少，等等。

以此可以构成较复杂的描述词，例如：

在右方远处、不太靠近、几乎平行，等等。

\pandocbounded{\includegraphics[keepaspectratio]{./media/OEBPS/Images/img_120.png}}

图120．邦加德问题47号

摘自莫·邦加德《模式识别》

这样，对一个典型框的内容------例如邦加德问题47号（图120）中的I-F------可以有下列不同的描述：

三个图形

或

三个白图形

或

右边有一个圆形

或

两个尖朝上的三角形

或

两个三角形和一个圆形

或

一个大图形和两个小图形

或

一个曲线图形和两个直边图形

或

一个里面和外面有同类图形的圆形

这些描述中的每一个都是通过一个``过滤器''来观察这个框的。脱离环境来看，其中每一个都可能是个有用的描述。但我们发现，在它们所处的特定的邦加德问题的环境中，它们全都是``错的''。换句话说，如果你知道邦加德问题47号中第Ⅰ组和第Ⅱ组的区别，然后给你一个上面那样的描述，作为对某个你未见到的图像的描述，那么这种信息将无法使你断定这个图像应该属于哪一组。在这个环境中，这个框的本质特征是它包括了：

一个圆形包含一个三角形

注意，如果某个人听到了这一描述，他无法据此``重建''原始图形，但他能``识别''某个图形是否具有这种性质。这有点像音乐的风格：你可能会准确无误地辨别出莫扎特的作品，但同时又写不出一首能让别人以为是出于莫扎特之手的曲子来。

\pandocbounded{\includegraphics[keepaspectratio]{./media/OEBPS/Images/img_121.png}}

图121．邦加德问题91号

摘自莫·邦加德《模式识别》

现在考虑一下邦加德问题91号（图121）中的Ⅰ-D框。在邦加德问题91号的环境中，一个繁琐然而是``正确''的描述是：

一个带有三个矩形凹入的圆形。

请注意这种描述的复杂性，其中``带有''这个词起了某种否定作用，意思是那个``圆''实际上并不是个圆：它差一点就是个圆了，只是\ldots\ldots 进一步来说，凹入的也并非标准的矩形。在我们使用语言描述事物的方式之中，存在着许多``宽松''之处。显然，许多信息已经被扔掉了，甚至本应扔掉得更多些。想预先得知哪些该扔、哪些该留，这是很困难的。因此，为了得到一个适当的折中方案，必须采用恰当的方法。当然，如果不得不恢复已经丢掉的信息，那就需要求助于低层的描述了（即模块化程度较弱的描述），就好像人能够不断地观察该图形，以求重组关于它的想法一样。那么，关键就在于想出一些明确的规则来说明如何完成下列任务：

为每个框构造尝试性描述；

把这些描述与两组中其它框的描述相比较；

重新构造描述，用下列手段：

（ⅰ）添加信息，

（ⅱ）去掉信息，

或（ⅲ）从另一个角度看同样的信息；

重复这一过程，直到发现能把两组图形区别开的东西为止。

模板和同一性检测器

尽最大可能使描述``在结构上彼此相似''，这可能会是个好的策略。它们所共同具有的结构将使它们之间的比较变得容易一些。在这种理论中有两个要素与该策略有关，其一是关于``描述模式''，或称``模板''的想法，其二是关于``同一性检测器''的想法。

先谈同一性检测器。它是出现在程序的各个层次上的``特派员''（事实上，不同的层次上可能会有不同种类的同一性检测器）。它不断地在个别的描述之中和不同的描述之间巡视，寻找重复出现的描述词或其它东西。一旦发现了某种同一性，就会触发各种重构操作，这既可能发生在单个描述水平上，也可能同时涉及多个描述。

现在来谈模板。在预处理完成后，首先要做的是设法构造一块模板，或者叫描述模式------一种适用于描述问题中所有框的``统一格式''。其基本想法是：描述常常可以通过某种自然的方式被分解成子描述，然后如果需要的话还可以把它们再变成子子描述。当你碰到属于预处理层次的基本概念时，这个分解过程就终止了。现在重要的是选择描述的分解方式，以反映所有框之间的共同点，否则将会引入一些多余而且无意义的``伪规律''。

\pandocbounded{\includegraphics[keepaspectratio]{./media/OEBPS/Images/img_122.png}}

图122．邦加德问题49号

摘自莫·邦加德《模式识别》

模板是在什么样的信息的基础上被构造出来的呢？最好还是看一个例子。请看邦加德问题49号（图122）。预处理得到的信息是每个框都包含若干个小圈和一个较大的封闭曲线。这个观察是很有价值的，应当被加到模板中去。这样关于模板的第一个设想可能是：

大封闭曲线：\underline{　　}

小圈：\underline{　　}

这很简单：在描述模板中有两个明确的``槽''，需要用子描述来把它们填上。

一个异层结构程序

``封闭曲线''这个词触发了一件有趣的事情。这个程序中最重要的模块之一是一种语义网络------``概念网''，在这个网络中所有已知名词、形容词等等被相互联接起来了，而联接的方式就说明了它们的相互关系。例如，``封闭曲线''紧密地联接于``内部''和``外部''这两个词。概念网充满了关于词间关系的信息，例如谁和谁是相反的、谁和谁是相似的、谁和谁经常同时出现，如此等等。在图123中表示了一张概念网的一小部分，对此，我们不久就会加以解释。但让我们先来追寻在解决第49号问题时所发生的事情。``内部''和``外部''这两个词被激活了，因为它们在网络中与``封闭曲线''相邻。这就提示模板构造者：最好还能为该曲线的内部和外部分别设立槽。因此，根据``尝试性''的精神，模板被尝试着改组成下列形式：

大封闭曲线：\underline{　　}

内部的小圈：\underline{　　}

外部的小圈：\underline{　　}

在找到子描述之后，``内部''和``外部''这两个词将调用有关过程去检查框中那些特定的区域。在邦加德问题49号中的Ⅰ-A框内会发现下列情况：

大封闭曲线：{圆形}

内部的小圈：{三个}

外部的小圈：{三个}

对同一个邦加德问题中的Ⅱ-A框有下列描述：

大封闭曲线：{雪茄形}

内部的小圈：{三个}

外部的小圈：{三个}

现在，一直和其它操作并行活动的同一性检测器发现，``三个''这个概念在所有和小圈有关的槽中同时出现，这就构成了着手进行第二次模板改组操作的一个强有力的理由。注意第一次是由概念网络建议的，而第二次是同一性检测器建议的。至此我们关于第49号问题的模板变成了：

大封闭曲线：\underline{　　}

内部的三个小圈：\underline{　　}

外部的三个小圈：\underline{　　}

现在``三个''的普遍性已经提高了一层------也就是说进入了模板------因此也就值得去探索它在概念网中的邻居。其中一个是``三角形''，这就提示我们小圈组成的三角形可能是重要的。这实际上正好把我们带进了一条死胡同------但事先怎么可能知道呢？这是一条典型的、人所要探索的死胡同，因此如果我们的程序也能发现它，那反而是件好事了！对于Ⅱ-E框来说，将会生成下列描述：

大封闭曲线：{圆形}

内部的三个小圈：{等边三角形}

外部的三个小圈：{等边三角形}

当然这里已经丢掉了大量信息，包括这些三角形的大小、位置和方向，还有许多其它东西。但这恰恰是构造描述------不仅使用原始数据------的特点所在。我们在第十一章所讨论的``汇集''也体现了同样的想法。

概念网络

\pandocbounded{\includegraphics[keepaspectratio]{./media/OEBPS/Images/img_123.png}}

图123．一个求解邦加德问题的程序的概念网络片段

``结点''通过``连线''彼此相联，而连线又可能被联接。通过把连线看成动词，把它所联接的结点看成主语和宾语，你可以从图中抽取出一些句子来

我们不必追寻第49号问题的整个解决过程，只需以此表明单个的描述、模板、同一性检测器和概念网络之间不断出现的反复的相互作用。我们现在应当更仔细地观察概念网络及其功能。图123中所示的一个简化了的网络片段表示了下列思想：

\begin{itemize}
\tightlist
\item
  ``高''和``低''是相反的。
\item
  ``上''和``下''是相反的。
\item
  ``高''和``上''是相似的。
\item
  ``低''和``下''是相似的。
\item
  ``右''和``左''是相反的。
\item
  ``左-右''差别相似于``高-低''差别。
\item
  ``相反''和``相似''是相反的。
\end{itemize}

注意在网络中的一切东西------不论是结点还是连线------均可被描述。在此意义下，网络中任何东西都不比别的东西处于更高的层次。图中还显示了网络的另一部分，它表示了下列思想：

\begin{itemize}
\tightlist
\item
  矩形是多边形。
\item
  三角形是多边形。
\item
  多边形是封闭曲线。
\item
  三角形和矩形的差别在于一个有3条边，另一个有4条边。
\item
  4与3相似。
\item
  圆是封闭曲线。
\item
  封闭曲线具有内部和外部。
\item
  ``内部''和``外部''相反。
\end{itemize}

概念网络必须非常大才行。它似乎只存储静态知识，或者叫描述性知识，而实际上并非如此。它的知识也可以说是过程性的，因为网络中的相邻性起到了引导者或``程序''的作用，可以告诉主程序如何深化它对框中图形的理解。

\pandocbounded{\includegraphics[keepaspectratio]{./media/OEBPS/Images/img_124.png}}

图124．邦加德问题33号

摘自莫·邦加德《模式识别》

例如，某些开始时的预感可能在后来被发现是错的，但其中包含着正确答案的萌芽。在邦加德问题33号（图124）中，开始往往会产生的想法是：第Ⅰ组的框都包含``有角的''图形，第Ⅱ组的框都包含``平滑的''图形。但进一步考察就会发现这是错的。不过这种直感还是有价值的，我们可以把它进一步推向深入，在概念网络中考虑与``有角的''相邻的词。它与``锐角''这个概念相邻，而后者显然是第Ⅰ组区别于第Ⅱ组的特证。这样，概念网络的主要功能之一就是允许开始的错误想法被一点点修改，直至变成正确的。

滑动和尝试性

与在两个紧密相联的词间``滑动''这个概念有关的一个概念，就是把一个给定对象看作自另一个对象变化而来。我们已经提到过一个典型的例子------即``带有三个凹入的圆''，其中事实上根本没有圆形。我们必须具备在适当的时候使概念产生形变的能力。没有什么绝对不变的东西。在另一方面，也不能把事情弄得模棱两可，以致于完全丧失意义。关键在于要懂得该在什么时候以及用什么方式从一个概念滑向另一个概念。

邦加德问题85\textasciitilde87号（图125）提供了一组很有趣的例子，其中的关键所在就是从一个描述滑向另一个描述的问题。邦加德问题85号很简单。让我们假设我们的程序在预处理阶段就能对``线段''进行标识。然后，对线段计数并找到邦加德问题85号中第Ⅰ组和第Ⅱ组的差别，这对它来说是比较容易的。现在来处理邦加德问题86号。它所用的一个一般性启发式规则是``设法找出最近行之有效的想法''。对最近用过的方法进行重复，这在现实世界中是非常普遍的，邦加德并不想把这种启发规则排除掉------事实上幸亏他强化了这一规则。这样我们就带着两个想法（``计数''和``线段''）着手分析第86号问题，而且这两个想法融成了一个：``对线段计数''。但实际上邦加德问题86号的关键是数``线列''而非``线段''，所谓``线列''是指一串首尾相连的线段（一条或多条）。程序发现这一点的可能方式之一是``线列''和``线段''二者均为已知，而且在概念网络中相邻。另一种方式是程序能``发明''\,``线列''这个概念------起码这是一个重要的前提。

\pandocbounded{\includegraphics[keepaspectratio]{./media/OEBPS/Images/img_125.png}}

图125．邦加德问题85\textasciitilde87号

摘自莫·邦加德《模式识别》

下一个是邦加德问题87号，其中``线段''这个概念发挥了进一步的作用。什么时候一条线段会被当作三条线段？（见Ⅰ-A框）程序应具有足够大的灵活性，才能在图形的给定部分的这种不同表示之间来往。一种明智的做法是把旧的表示方式存储起来，而不是把它们忘掉，否则可能还得把它们重建起来，因为谁也不能保证新的表示方式就一定比旧的好。这样，和每个旧表示方式一道，还应当存储一些喜欢它们或不喜欢它们的理由。（似乎开始复杂起来了，不是吗？）

元描述

至此我们将遇到识别过程的另一个至关重要的部分，它与抽象的层次和元描述有关。为此让我们再看一下邦加德问题91号（图121）。在这里应当构造哪一种模板呢？情况实在太复杂，难以决定从何处下手。但这本身就是一条线索！这条线索说明，组与组之间的差别很可能存在于较高的抽象层次上，而不是存在于简单的几何描述之中。这个观察结果揭示程序应该构造``对描述的描述''------即``元描述''。或许在这个第二层上会出现某些共同特征，而且如果我们走运，就会发现足够多的共同点，以导致一个对元描述模板的精确刻划！因此我们在没有模板的情况下仍然往前走，为各种框构造描述。然后，一旦完成了这些描述，我们再描述它们。我们的元描述模板将有哪些种槽呢？或许包括这样一些：

所用的概念：\underline{　　}

重复出现的概念：\underline{　　}

槽名：\underline{　　}

所用的过滤器：\underline{　　}

在元描述中可能还需要许多其它的槽，我们这里不过是举个例子而已。现在，假设我们已经对邦加德问题91号的Ⅰ-E框作了描述。它的（无模板）描述可能会是这个样子：

水平线段

立于该水平线段之上的垂直线段

立于该水平线段之上的垂直线段

立于该水平线段之上的垂直线段

当然许多信息已经被丢掉了：如三段垂直线段等长、等间隔的事实，等等。但上述描述是很有可能产生的。因此元描述可能会是这个样子：

所用的概念：垂直、水平、线段、立于

描述中的重复：三次出现``立于该水平线段之上的垂直线段''

槽名：\underline{　　}

所用的过滤器：\underline{　　}

并非元描述中的所有槽都必须被填上。信息在这一层上可以像在``直接描述''层上一样被丢掉。

现在如果我们对第Ⅰ组中的其它任何框构造描述，那我们每次都会在``描述中的重复''这个槽中填上``三次出现\ldots\ldots''这个词组。同一性检测器会注意到这一点，并会选出``三次''作为第Ⅰ组中的框在一个相当高的抽象层次上的一个显著特征。类似地，用元描述的方法，也会识别出第Ⅱ组的标志是``四次''。

灵活性是重要的

现在你可能会提出异议，说在这种情况下求助于元描述方法就好像用高射炮打蚊子，因为三次和四次这种特征在低层次上会同样容易地被发现，只要我们在构造描述时稍微改变一下方式就行了。是的------但重要的是，沿不同的途径解决这些问题是完全可能的。在程序中可能存在着大量的灵活性，因此，即使程序一时``误入歧途''，这也是不足为怪的。在任何情况下，我都要阐明下列一般原理：一旦预处理器发现了大量差异，并因此使得模板难以构造，这本身就成了一条线索，说明这个问题牵扯到比预处理器所知道的东西抽象层次更高的概念。

集聚与过滤

\pandocbounded{\includegraphics[keepaspectratio]{./media/OEBPS/Images/img_126.png}}

图126．邦加德问题55号

摘自莫·邦加德《模式识别》

现在让我们来看另一个问题：如何扔掉一些信息。这涉及到两个概念，我称它们为``集聚''和``过滤''。``集聚''是指构造这样一些描述：其焦点是框中图像的某一部分，而排除掉其它的部分。``过滤''是指构造这样一些描述：它以某种特定的方式来观察框的内容，而完全不顾其它方面。因此二者是互补的：集聚被用来处理客体（不严格地说，即名词），而过滤被用来处理观念（不严格地说，即形容词）。邦加德问题55号（图126）就是有关集聚的一个例子。在这里，我们的注意力集聚于那个凹入和它旁边的小圆圈，而排除掉框中的其它成分。邦加德问题22号（图127）提供了一个过滤的例子。在这里我们必须把除了图形大小之外的其它特性统统滤掉。而解决邦加德问题58号（图128）的问题则需要把集聚和过滤结合起来。

\pandocbounded{\includegraphics[keepaspectratio]{./media/OEBPS/Images/img_127.png}}

图127．邦加德问题22号

摘自莫·邦加德《模式识别》

\pandocbounded{\includegraphics[keepaspectratio]{./media/OEBPS/Images/img_128.png}}

图128．邦加德问题58号

摘自莫·邦加德《模式识别》

为了找到集聚和过滤的方法，最重要的途径之一就是通过另一种``集聚''：即只检查一个特定的简单框------就是说其中的客体要尽可能少。一种非常有效的办法是把两组中最有代表性的框进行相互比较。但在得到描述之前，怎能说出哪些框是有代表性的呢？实际上检测代表性的一种方式就是寻找一个从预处理器那里得到的特征最少的框。这在一开始就可以确定，因为它不需要先准备好的模板。事实上，这可能是一种发现应构造在模板内的特征的有效方法。在邦加德问题61号（图129）中，使用这种技术可以很快找到答案。

\pandocbounded{\includegraphics[keepaspectratio]{./media/OEBPS/Images/img_129.png}}

图129．邦加德问题61号

摘自莫·邦加德《模式识别》

科学研究和邦加德问题的世界

我们可以把邦加德问题构成的世界看成一个进行``科学研究''的小天地------也就是说，这种研究的目的是辨识该世界中的模式。在寻找模式的过程中，模板被装了又拆，拆了又装；槽的普遍性从一个层次转换到另一个层次；进行了集聚和过滤；如此等等。这里有复杂程度相差悬殊的各种发现活动。按照库恩的理论，某些被称为``范式转换''的罕见事件构成了``常规''科学与``观念革命''之间的分水岭。但这种观念似乎不适用于我们所讨论的问题，因为我们可以看到，在该系统中范式转换是随时随地都在发生的。描述的流动性使得范式转换可能以各种不同的规模发生。

\pandocbounded{\includegraphics[keepaspectratio]{./media/OEBPS/Images/img_130.png}}

图130．邦加德问题70\textasciitilde71号

摘自莫·邦加德《模式识别》

当然，某些发现比另一些具有更大的``革命性''，因为它们的影响更为广泛。例如，有人可能会发现第70和71号问题（图130）在一个足够抽象的层次上是``同一个问题''。关键是要看出二者都包含了二级嵌套和一级嵌套间的对比。这对邦加德问题来说是一个新的发现层次。甚至还有一个更高的层次，即把所有邦加德问题看成一个整体。如果有人从没见过这些问题，那把这些问题想出来本身就是个很好的问题。想出这些问题需要一种革命性的洞察力，但必须指出的是，使这种发现得以做出的思维机制与那种被用来解决一个具体的邦加德问题的思维机制并无不同。

根据同样的道理，起初的科学研究并不能被划分成``常规''阶段和``观念革命''阶段，而是充满了范式转换------它们之间只有大小之分和层次之别。INT和G图的递归图形（图32和图34）为这种想法提供了一个几何模型：它们在每个层次上都具有不连续的跳跃，而不仅仅在最顶层是如此------只是层次越低跳跃幅度越小罢了。

与其它类型思维的联系

为了把这个完整的程序置于环境之中，让我们设想两种把它和认知的其它方面联系起来的方式。不仅它依赖于认知的其它方面，而且后者也反过来依赖于它。先让我说明一下它是怎样依赖于认知的其它方面的。要明白什么时候应忽略某些差别、尝试重新构造一个描述，以及回溯、转换层次等等，都是需要直觉的，而这种直觉可能只源于思维中的一般经验。因此，为程序的这些关键性的方面设计启发式规则可能是很困难的。有时，人关于世界中真实对象的经验对他描述或重新描述某些框的方式会产生微妙的影响。例如，谁知道某人对树木的熟悉能在多大程度上有助于他解决邦加德问题70号呢？我们无法断定对人来说与这些谜题有关的概念子网络是否能很容易地从整个网络中被分离出来。我们更倾向于相信一个人得之于观察和把握实在对象------如梳子、火车、线绳、木块、字母、橡胶带、等等------的直觉在解决这些谜题的过程中起着一种不可见的重要引导作用。

反过来，对真实世界情景的理解肯定极大地依赖于视觉表象和空间直觉，因此，如能以灵活有效的方式表示像邦加德问题中出现的那种模式，一定会对提高思维过程的一般效率有所贡献。

我觉得邦加德的问题都是精心设计出来的，它们具有一种普遍性，即每个问题的正确答案都是唯一的。当然有人会对此提出指责，说我们所谓的``正确''实际上是深刻地依赖于我们人类的，而其它星系中的生物可能会有完全不同的看法。虽然正反两种观点都缺乏具体证据，我还是确信邦加德问题取决于一种关于简单性的观念，而这种观念是不仅限于地球上的人类的。我在前面说过，熟悉梳子、火车、橡胶带等等地球所特有的物品或许是很重要的，但这并不排斥``我们关于简单性的观念具有普遍性''这种想法。原因是起决定作用的并不是这些个别物体，而是事实上它们一起确定了一个广阔的空间。很可能外星人和我们同样面对大量的人造或天然物体，而且也从中获得了大量经验。因此，我相信解邦加德问题所用的技巧位于离``纯粹''智能的核心很近的地方------如果这种所谓``纯粹''智能真的存在。因此，如果要研究发现模式或消息的``固有意义''所需的能力，那从这里着手是很合适的。遗憾的是我们只摘出了他的令人兴奋的题目中的一小部分。我希望读者们能从他的著作中（见参考书目）找到所有问题，并对它们有所了解。

有一些关于视觉模式识别的令人惊异的问题，我们人类似乎已将其完全``吸收''在潜意识之中了。这些问题包括：

\begin{itemize}
\tightlist
\item
  识别面孔（在年龄变化、表情变化、光线变化、距离变化、角度变化等之下的面孔不变性。）
\item
  在森林和山岭中识别足迹------这总以某种方式使我觉得是我们模式识别中最微妙的行为之一------动物也有这种本领
\item
  流利地阅读以成百上千种不同的字型印出的文字材料
\end{itemize}

传送消息的语言、框架和符号

为了处理在模式识别和人工智能程序的其它应用中出现的复杂性，卡尔·海威特提出了所谓``演员''形式。（类似于艾伦·凯等人开发的``Smalltalk''语言）。根据这种方案，程序被写成一组相互作用的``演员''的集合，这些演员可以在彼此之间往返传送精心设计的``消息''。从某个角度看，这有点像一个由相互调用的过程所组成的异层结构集合。主要区别在于：过程间往返传送的参数通常数量很少，而演员之间相互交换的消息却可以任意长、任意复杂。

这种能够交换消息的演员在一定程度上成了有自主性的专职人员------事实上，甚至有点像自主的计算机，而消息就像其中的程序。每个演员都可以用自己特有的方式来解释任何给定的消息，这样，一条消息的意义将依赖于对它进行解释的演员。出现这种情况的原因是：演员都备有一段用来解释消息的程序，因此，有多少个演员就可能有多少个解释程序。当然，可能许多演员的解释程序是相同的，但这事实上将带来更大的好处，就好像在一个细胞中至关重要的是具有大量相同的核糖体，它们漂散在细胞质中，并都以一种共同的方式来解释消息------在这种情况下，消息是由RNA带来的。

一个有趣的问题是考虑一下如何把框架这个概念和演员这个概念结合起来。让我们把那种能生成和解释复杂消息的框架称为``符号''：

框架＋演员＝符号

我们现在已经是在讨论怎样实现在第11章和第12章中提到的那种难以捉摸的``活跃符号''了，因此在本章中，``符号''就是指那样的符号。顺便说一下，哪怕你不能马上看出这一综合是如何完成的，那也无须感到不安。虽然这一点显然是人工智能中最吸引人的发展方向之一，但也并非让人一目了然。进一步我们可以肯定，即便是对于这些概念的最佳综合，我们也将发现它的功能远弱于人脑中的实际符号。在这种意义下，把框架和演员的综合物称为``符号''未免为时过早，不过这提供了观察事物的一种最佳方式。

让我们转回到与传送消息有关的话题上来。每条消息是应当专门传向某个目标符号，还是应被扔到广阔的空间之中------像信使RNA被扔到细胞质中那样，去寻找接收它的核糖体？如果消息是有目标的，那每个符号就必须有一个地址，而传给它的消息就应当总是被送到那个地址。另一方面，也可能存在一个消息中心接收站，每条消息只须在那里老老实实等着，直到某个需要它的符号来把它取走。这对应于那种``存局待领''的邮件。或许最好的结论是允许两种类型的消息并存，同时满足不同类型的需要------专送邮件、快件、慢件等等。整个邮政系统为消息传送语言提供了丰富的思想，包括下列奇特的东西：备有回信地址和邮票的信封（发消息的人想尽快得到回信）、邮包（以某种很慢的方式传送非常长的消息）等等。如果你已经了解了邮政系统所包含的思想，电话系统会给你更多的灵感。

酶与人工智能

当然另一个消息传递------事实上，是一般信息加工------的丰富的思想源泉是细胞。细胞中的某些东西------特别是酶------很像演员。每种酶的活动场所都起着过滤器的作用，它只识别某种特殊的``底物''（即消息）。因此事实上每种酶都有一个``地址''。酶已经被``编了程序''（凭借其严整的结构）以完成那些对``消息''的操作，然后再将它释放出来。通过这种方式，一旦一条消息沿化学通道从一个酶被传送到另一个酶，就可以完成很多操作。我们已经描述过那种可能在细胞中发生的精致的反馈机制（通过抑制）。这些机制表明复杂的过程控制可以通过存在于细胞之中的各种消息传递来实现。

酶的最显著特性之一就是它们总是闲待在那里，等着被收到的底物所触发。然后，一旦这种底物来到，酶就会一下子跳起来投入行动，就像一棵捕蝇草一样。这种``一触即发''的程序已经被用在人工智能中，称为``精灵''。在这里重要的是下述想法：应当有许多种类不同的可触发子程序处于待触发状态。在细胞中，所有复杂的分子和细胞器都是通过一些简单步骤构成的。在这些新结构中，常常有些结构本身又是酶，而且它们又参与构造另外一些酶，如此下去。这种酶的递归``多级瀑布''能够对细胞的所作所为产生巨大影响。这种简单的逐步构造过程已经被引入人工智能之中，用于构造有用的子程序。例如，重复作用可以在我们头脑的硬件中形成新电路，因此经常重复的行为会变得被编码于意识水平之下。如果有一种类似的方法能对有效的代码片段进行综合，使它们能和在某个高层``意识''水平上学到的东西产生同样的操作序列，那将是极其有用的。酶的阶梯将为如何做到这一点提供一个模型（杰拉尔德·萨斯曼编写的``海客''程序对短小的子程序进行综合和调试的方式就有点像酶的阶梯）。

邦加德问题求解程序中的同一性检测器就可以用像酶那样的子程序来实现。同一性检测器可能会像酶那样带些随机性地东游西荡，不时在各处碰到小型数据结构。一旦在其两个``活性部位''上装入了同样的数据结构，该同一性检测器就会向程序的其它部分（演员）发出一条消息。只要程序还是顺序执行的，那么使用好几个同一性检测器的副本就没有太大意义，但在一个真正的并行计算机中，控制子程序的副本数，正是对一项操作被完成前的预期等待时间进行控制的一种方式。这就像在一个细胞中，控制了一种酶的副本数，也就控制了相应功能的完成速度。如果能综合生成新的同一性检测器，那就相当于在我们的头脑中使模式检测过程渗透到一个更低的层次。

裂变与聚变

``裂变''和``聚变''是两个有趣而且互补的想法，涉及到符号的相互作用。``裂变''是指一个新符号逐渐脱离它的母符号（即被用来为它做复制模板的那个符号）。``聚变''是指两个（或多个）原来互不相干的符号由于参与了某个``连带激活''，彼此频繁地往返传送消息，紧密地联系在一起，因而以后这个复合体可以被作为一个符号来访问。裂变或多或少是一个不可避免的过程，因为一旦一个新符号从老符号上被``剥下来''，它就变成自主的，而且在它自己的内部结构中反映出它与外部世界的相互作用，因此开始时的一个完全的复制品不久就会变成不完全的，而且此后慢慢会变得越来越不像那个它从其上被``剥下来''的符号。聚变是一个更为复杂的过程。什么时候两个概念才真的变成了一个？是否存在着某个发生聚变的确切时刻？

``连带激活''这个概念打开了一个装满了问题的潘多拉之匣。当我们提到某种``东西''时，我们是否想到了``东方''和``西方''？英文里的``chairman''{[}主席{]}到底与``chair''{[}椅子{]}和``man''{[}男人{]}有多少关系？（有人至今还在指责这种词充满了大男子主义）。一个德国人在说``手套''{[}Handschuhe{]}时是否真的想到了``穿在手上的鞋''（该词的字面意义）？到底这些部分在多大程度上与整体产生共鸣，恐怕是因人而异、因地而异的。

与符号的``聚变''这个概念有关的实际难题是：很难设想出一个通用的算法，能从相互碰撞的符号中构造出有意义的新符号来。这就像两条DNA碰到了一起。你怎样才能从它们中各取出一部分，把它们拼成一条有意义并且能活下去的新DNA，而且使它构成同一物种的一个个体------或一个新物种------的遗传编码？一次DNA片段的随机组合就构成了某种能存活的东西的遗传密码，这种可能性可以说是微乎其微------就像两本书中的词重新随机组合能得到另一本书的可能性一样。想让这种重组DNA在除最低层外的任何层次上仍有意义，这种可能性极其微小，其原因完全是由于在DNA中具有太多的意义层次。对``重组的符号''来说也是一样。

《螃蟹卡农》的渐成过程

我来把我写的对话《螃蟹卡农》当作一个典型事例。其中，有两个想法在我的头脑中发生碰撞，并以一种新方式联结起来了。这时，一种新语词结构突然活生生地出现在我的脑海中。当然我现在仍然可以分别考虑螃蟹卡农音乐和言语对话------它们仍能彼此独立地被激活，但《螃蟹卡农》对话所对应的合成符号也具有它自己所特有的激活方式。为了更详细地描述这个聚变或``符号重组''的概念，我想把我的《螃蟹卡农》的发展作为一个事例来研究。这当然是因为我对此比较熟悉，也还因为这个例子很有趣，而且很典型地说明了一个想法可以被推广到什么程度。我将列举其各个阶段，这些阶段是根据``成熟分裂''的各阶段来命名的，而``成熟分裂''是指那种伴随有``交换''------或称基因重组------发生的细胞分裂，这是进化中的多样性的来源。

\textbf{前期}

我开始时只有一个很简单的想法------对一段音乐，比如说一段卡农，可以用语言来进行模拟。这来自下列观察：通过一种双方共有的抽象形式，一段文字和一段音乐可以联系起来。下一步涉及到设法使这种模糊的预感中的潜在可能变为现实，在这里我想到了卡农中的``声部''可以被映射到对话中的``角色''------这仍然还是个很显然的想法。

然后我集中精力于特殊的卡农，我想起在《音乐的奉献》中有一段螃蟹卡农。那时我刚刚着手写对话，而且其中只有两个角色：阿基里斯和乌龟。由于巴赫的螃蟹卡农有两个声部，这就完全对上号了：阿基里斯应当是一个声部，而乌龟是另一个声部，二者做相同的事，只是一个正行一个逆行。但在这里我遇到了一个问题：这种颠倒过程在什么层次上发生呢？字的层次？词的层次？句子的层次？在一番考虑之后，我觉得``对白''层次可能最为恰当。

现在巴赫的《螃蟹卡农》的``骨架''已经被移植到言语形式中来了，至少在计划中是如此。只剩下一个问题。当两个声部在中间相互交叉时，将会有短期的极度重合现象：这是个讨厌的疵点。怎么办呢？这时发生了一件奇怪的事，一种典型的层次交错的创造性行为：《螃蟹卡农》中的``螃蟹''一词在我脑海中闪现了，这无疑是由于它和``乌龟''这个概念具有某种共同的抽象特性------我立刻明白了：在对话的中心，我应当为去掉那种重合的效果而插入一段特殊的台词，它要出自一个新角色：一只螃蟹！这就是在《螃蟹卡农》的``前期''对螃蟹的设想：处于阿基里斯和乌龟的交叉点。（见图131）

\pandocbounded{\includegraphics[keepaspectratio]{./media/OEBPS/Images/img_131.png}}

图131．对话《螃蟹卡农》模式图

\textbf{中期}

上面是我的《螃蟹卡农》的骨架。然后我就进入了第二个阶段------``中期''------我必须为这个骨架填上肉，这无疑是个艰巨的任务。我对此进行了大量的尝试，逐渐习惯了那种要求相邻的两段台词正读和逆读时都必须有意义的表达方式，而且反复试验去看哪些有双重意义的话能帮助我进行这种形式的写作（例如``周末愉快''）。我开始曾为这篇对话写出过两稿，二者都很有趣，但比较薄弱。我曾有一年多时间中断了本书写作，当我回到《螃蟹卡农》的写作上来时，又有了些新想法。其中之一是在对话中谈到一段巴赫的卡农。原先我的计划是谈到《音乐的奉献》中的《反向增值的卡农》（我称之为《树懒卡农》）。但开始之后我发现这似乎有些无聊，因此我才勉强决定在我的《螃蟹卡农》中，我要谈到巴赫自己的《螃蟹卡农》。实际上这是个重大的转折点，但我当时并没有意识到。

这样如果一个角色要谈论一段巴赫的音乐，那么另一个角色相应地说些完全相同的话岂不是太笨了吗？好，既然艾舍尔在我的头脑和著作中扮演着类似于巴赫的角色，那是否可以通过某种方式稍微对台词作一点修改，使之谈到艾舍尔呢？归根结底，在严谨的卡农艺术中，为了达到典雅和优美，偶尔也难免出现不完全的模拟。在我有了这个念头之后不久，我的脑海里冒出了《白天与黑夜》这幅画（图49）``显然''，我想到，``这是一种形象化的螃蟹卡农，本质上是由两个互补的声部奏着同一个主旋律，一个向左，一个向右，并保持彼此协调！''这里又一次出现了这个思想：同一个``概念骨架''被应用于两种不同的介质上------在这里，是音乐与美术。因此我让乌龟谈论巴赫，让阿基里斯谈论艾舍尔，二者使用平行的语言。这多少是对严格模拟的小小偏离，但显然还保持有螃蟹卡农的精神。

到这时，我开始意识到某种奇妙的事情正在发生：即这段对话变成了自指的，尽管我没有想要这样做！更进一步的是，它是一种间接自指，其中的角色并没有直接谈到他们所在的对话，而是谈到了与其同构的结构（在某种抽象意义上）。换句话说，我的对话现在和哥德尔串G具有共同的``概念骨架''，因此能像中心法则那样对应于G，构成一个``中心螃蟹映射''。这使我非常兴奋，因为从天上掉下来了一个哥德尔、艾舍尔和巴赫的美妙结合。

\textbf{后期}

下一步很令人惊讶。多年前我有一本卡洛林·麦克吉拉弗里关于艾舍尔的镶嵌图案的专著，有一天我翻阅这本书的时候，其中的插图23（本书中的图42）一下子吸引了我的注意力，因为我是以一种前所未有过方式看到它的：一个名副其实的螃蟹卡农------形式和内容都是螃蟹式的！艾舍尔没有给这幅画标上标题，而且他也曾用许多其它动物的形象画过类似的镶嵌图案，因此，这种形式和内容的巧合大概只是我恰好发现的。但不论是否是偶然发现的，这幅无题插图正是我关于本书的主要设想之一的缩影：把形式和内容合为一体。这样，我就高兴地把它命名为《螃蟹卡农》，用它替换了《白天与黑夜》，而且相应地修改了阿基里斯和乌龟的谈话。

至此事情并没有完。我当时已经迷上了分子生物学，一天在书店里仔细翻看沃森的著作时，我在索引中看到了``回文''这个词。当我去读有关内容时，发现了一个奇妙的现象，即DNA中的螃蟹卡农结构。不久我对螃蟹的独白做了相应的修改，使它包含一小段话，表明它把它混淆前进和后退的偏好归因于它的基因。

\textbf{末期}

最后一步是几个月后发生的，当时我在和别人讨论一幅表现DNA中的螃蟹卡农节段的图画（图43），我发现Adenine{[}腺嘌呤{]}、Thymine{[}胸腺嘧啶{]}、Cytosine{[}胞嘧啶{]}中的首字母``A''、``T''、``C''恰巧------说来也怪------和Achilles{[}阿基里斯{]}、Tortoise{[}乌龟{]}、Crab{[}螃蟹{]}的首字母相对应，更有甚者，正像腺嘌呤和胸腺嘧啶在DNA中配对一样，阿基里斯和乌龟在对话中也配成一对。我又考虑了一下，通过另一次层次交叉，看出应当让在DNA中与``C''配对的字母``G''代表``Gene''{[}基因{]}。我又一次回到对话中，对螃蟹的发言作了个小手术以反映这个新发现，这时我已经在DNA的结构和这段对话的结构之间建立了一个对应关系。在这种意义下，DNA可以被说成是对一个表现型进行编码的遗传型：这个表现型就是对话的结构。最后这一步极大地提高了对话的自指程度，并使之具有了更多的涵义，多得我自己都没有料到。

概念骨架和概念映射

上面已经或多或少地对《螃蟹卡农》的渐成过程进行了总结。整个过程可以被看作是思想在不同的抽象层次上不断地相互映射的过程。这就是我所说的``概念映射''，而那种联接两个不同想法的抽象结构就是``概念骨架''。这样，就有一个关于《螃蟹卡农》的抽象观念的概念骨架：

{一个结构中具有两个完成同样功能的部分，只是它们以相反的方向运动。}

这就是一个具体的几何图像，我们的头脑可以像处理一个邦加德模式那样处理它。实际上，当我今天想到《螃蟹卡农》的时候，我把它看成两根相互交叉的带子，在中间被一个``结''联接起来（这就是螃蟹的独白）。这幅栩栩如生的图像在我脑海里一下子就映射到另一幅图像上：两条同源的染色体，中间通过着丝点相互联接。这幅景象直接来自于成熟分裂，如图132所示。

实际上，正是这幅图像启发我用成熟分裂的分期法来描述《螃蟹卡农》的演化过程------当然，这本身又是概念映射的另一个例子。

\pandocbounded{\includegraphics[keepaspectratio]{./media/OEBPS/Images/img_132.png}}

图132．两条同源的染色体

中间通过着丝点相互联接

重组的思想

可以用几种不同的技术实现两个符号的聚变。一种办法是把两种想法并排着列出来（仿佛它们是线性的！），然后审慎地从每个想法中选取一些片段，把它们重新组织在一个新的符号之中。这很强烈地使我们回想起基因的重组。那么，着丝点交换了些什么？又是怎样交换的呢？它们交换的是基因。在符号中与基因相似的是什么呢？如果这些符号中有框架式的槽，那或许交换的就是槽。但交换的是哪些槽呢？为什么是它们？在这里《螃蟹卡农》那样的聚变可以提供某些启示。在把``螃蟹卡农音乐''映射到``对话''上的过程中包括了若干个辅助映射：事实上它们是由它导出的。这就是说，一旦决定要使这两个观念发生聚变，问题就变成：从某个能使类似的部分得以表现的层次上来观察它们，然后逐步在各部分间建立映射关系，如此递归地工作下去，直到发现某个满意的层次为止。例如在这里，当抽象地看《螃蟹卡农》和``对话''时，``声部''和``角色''表现为对应的槽。但这种抽象观点是从哪里来的呢？这是映射问题的关键所在------抽象观点来自何方？你应怎样抽象地去观察特定的概念？

抽象、骨架、类比

我所说的``概念骨架''是指沿某个特定角度对一个概念进行抽象后所得到的一种观点。实际上，我们一直在和概念骨架打交道，只是不常用这个名字罢了。例如，与邦加德问题有关的许多想法就可以用这个术语重新表述。如果我们发现两个或多个想法具有共同的概念骨架，这总是有趣的，而且往往是重要的。一个例子就是在《对实》的开头提到的那个古怪概念的集合：双目独眼龙、双轴独轮车、``双头翘''翘翘板、乒乒球、无对手的决赛、双面莫比乌斯带、``巴赫孪生兄弟''、双左手钢琴协奏曲、单声部赋格、单手鼓掌、双声道单喇叭录音机、两个八分位。所有这些想法都是``同构''的，因为它们都具有下述概念骨架：

把复数事物变成单数，然后错误地复数化。

本书中另外两个具有共同概念骨架的是（1）乌龟对阿基里斯那个求第一个和最后一个部首都是``虫''的词的解答（乌龟的解答就是``虫''，其中两次出现合并为一次），以及（2）帕普斯和吉伦特对\heititext{驴桥定理}的证明，其中一个三角形被重复感知为两个。附带地说，这些滑稽的编造可以被叫作``半对词''。

一个概念骨架像一个常量特征集（区别于参量或变量）------这些特征在虚拟的重放镜头或映射操作中不产生滑动。由于不含有可变的参量或变量，它可能成为若干不同想法的不变内核。它的每个``例''，例如``双轴独轮车''，都具有若干层变量，因此可沿不同途径``滑动''。

虽然``概念骨架''这个名字听起来有些刻板、专横，实际上其中有许多活动余地。概念骨架可能存在于若干不同的抽象层次之上。例如，我们已经指出过，邦加德问题70和71号之间的那种同构就涉及到一种比单独解决其中任何一个问题所需的层次都高的概念骨架。

多重表示

概念骨架不仅必须存在于不同的抽象层次上，而且还必须存在于不同的概念``维度''上。让我们以下面这句话作为一个例子：

``美国副总统是美国政府这部汽车上的备用轮胎。''

我们是怎样理解它的意义的？（其中的幽默成分当然是很重要的，但暂不考虑。）如果在事先没有任何动机的条件下，告诉你说``把美国政府看作一部汽车''，你可能会沿任何对应关系考虑下去：驾驶盘=总统，等等。那什么相当于第一夫人呢？车座上的安全带又是什么？由于被比较的两个对象差异太大，这种对应几乎不可避免地要包括功能方面。因此，你会从你的概念骨架库中取出关于汽车的一部分表示，其中只包括那些与功能有关的因素，而不考虑诸如车子的外型之类的东西。更进一步来说，完全可以在一个足够高的抽象层次上工作，在那里``功能''是在较广的意义下被使用的。这样，对备用轮胎的功能可以有下述两种定义：（1）``破车胎的替换物''和（2）``汽车上某个已失效部件的替换物''，显然在这种情况下我们更愿意要后一个。这只不过是由于汽车和美国政府差别太大，只能在较高的抽象层次上建立二者间的对应关系。

一旦检查完这个特定的句子，对应关系就被迫集中于某一个方面了------但这不管怎么说也不能算是一个麻烦的方式。事实上，在各种概念骨架中你已经有了一个副总统的概念骨架，即``政府中某个已失效的部分的替换物。''因此这种强迫对应可以很容易地生效。但为了作个对照，不妨假定你为``备用轮胎''取出了另一个概念骨架------比如说是描述其外型的。在这种描述中可能包括说一个轮胎是``既圆又鼓的''。显然，这就误入歧途了。（不过也很难说，我的一个朋友曾指出，不少副总统就是圆鼓鼓的！）

出入口

每种特殊的思维类型的主要特点之一，就是看新经验是如何被分类装入记忆中的，因为这就定义了将来把它们取出来时所用的``把柄''。而对事件、对象、思想等等------对所有可以被想到的东西------来说，存在着各种不同的``把柄''。有个事情使我对此深有感触：我常常在伸手打开车内的收音机时惊讶地发现它已经被打开了！之所以发生这种情况，是因为在表示收音机时使用了两种相互独立的方式，一种是``音乐产生器''，另一种是``烦恼解除器''。我知道音乐正在响着，但我又有点心烦，而还没等这两种认识有机会相互作用，我的伸手反应已经被触发了。有一天我刚把收音机留在一家修理部，开车之后想听音乐，就又一次发生了这种伸手反应。真叫怪。关于同一个东西还有许多其它表示方式，例如

\begin{itemize}
\tightlist
\item
  有发亮的银旋钮的东西
\item
  有过热问题的东西
\item
  需要我躺在地上装的东西
\item
  会嗡嗡响的东西
\item
  有滑动刻度盘的东西
\item
  可以用许多方式来表达的东西
\end{itemize}

所有这些都能起``出入口''的作用。尽管它们都从属于我那个表示汽车用收音机的符号，但通过其中一个达到那个符号，并非就打开了所有其它出入口。因此当我伸手把它打开的时候，我不太可能联想起我躺在地上把它往车上装的情景。反之，当我躺在地上拧螺丝的时候，大概也不会想起我用它听《赋格的艺术》的时光。在一个符号的这些方面之间是有一些``隔板''的，它们使得我的思想不至于像在自由联想时那样四处溢流。我的精神隔板是重要的，因为它们为我的思想流提供了容器和渠道。

有一个地方这些隔板非常僵硬：它们把不同语言中对应于同一个事物的词隔离开来了。如果这些隔板不是强有力的，一个会两种语言的人就可能不时在两种语言之间滑来滑去，这将使人感到很不舒服。当然，一个同时学习两种新语言的成年人常常会把二者的词汇搞混------语言间的隔板太薄，很易冲破。那些译员们是很有趣的，因为他们可以说任意一种语言，好像他们的隔板是坚不可摧的，而一旦需要，他们又能取消这些隔板，允许从一种语言进入另一种语言，这样他们才能翻译。斯坦纳的母语有三种，他在《巴别塔之后》中用了几页的篇幅来谈法语、英语和德语在他思维的各个层次上的相互混合，以及他的几种不同的语言是怎样为进入概念提供不同的出入口的。

受迫匹配

看出两个概念在某个抽象层次上具有共同的概念骨架之后，会产生各种各样的后果。通常在第一阶段你把二者都收入``镜头''，而且用高层匹配做指导，设法辨识出对应的子概念来。有时匹配过程可以被递归地扩展到下面若干层，以揭示出深刻的同构关系。有时这一过程停止得较早，从而揭示出了一种相似关系。也有这样一些时候：高层的相似关系具有很大的强制性，以至于即使这种对应没有明显地延续到低层，你也得去造出一个来。这就是``受迫匹配''。

受迫匹配每天都出现在报纸上的政治漫画中。一个政界人物被画成一架飞机、一艘船、一条鱼、蒙娜丽莎；一个政府成了一个人、一只鸟、一座油井；一项条约变为一个公文包、一把宝剑、一罐蚯蚓；如此等等。令人迷惑的是我们怎么如此容易就完成了作者所暗示的对应，而且恰好达到预想的深度。我们不会使这种对应进行得过深或过浅。

把一种东西压到另一种东西的模子中的另一个例子，是我在选择用成熟分裂的词汇来描述我的《螃蟹卡农》的发展时发现的。这个发现分成几个阶段。开始，我注意到《螃蟹卡农》和通过着丝点相联的染色体的图像之间具有共同的概念骨架，这就为受迫匹配提供了灵感。然后我看到包括``生长''、``阶段性''和``重组''的高层相似性。然后我只不过是尽可能把这一类比推向前进。像在邦加德问题的解决过程中那样，``尝试性''起了重大的作用：我在发现一个我认为有说服力的匹配之前忽而前进，忽而后退。

关于概念对应的第三个例子是中心法则映射所提供的。我最初注意到在数理逻辑学家和分子生物学家的发现之间有一种高层相似性，然后就在低层次上追寻它，直到我发现一种很强的类似关系为止。为了进一步强化它，我选择了一种模拟遗传密码的哥德尔配数法。这就是在中心法则映射中唯一的受迫匹配因素。

受迫匹配、类比和比喻是无法轻易区分开的。体育播音员常常使用一些难以分类的生动想象。例如把一场橄榄球赛描述成``主队横扫客队''，很难说你会由此想象出什么样的景象。你是不是给整个球队配上了笤帚？还是给每个队员配上了笤帚？或许都不是。可能出现的情况是，顿时在你脑海中闪现出笤帚横着在扫街上的树叶，然后通过某种神秘的方式，只有其中相关的部分被抽了出来，并转换到球队的表演中去。在你这样做的一刹那，橄榄球队和笤帚之间在多大程度上彼此对应呢？

概括

现在要设法把这些东西集中一下。我已经提出了关于符号的创造、处理和比较的一系列想法，其中多数与某种形式的滑动有关，即认为概念是由一些联系得或松或紧的元素组成的，它们来自于嵌套的环境（框架）的不同层次。那些联系较松的元素可以容易地被去掉或替换，这依赖于不同的环境，能构造出``虚拟的重放镜头''、受迫匹配或是类比。两个符号的聚变可能是来自这样一个过程，其中每个符号都有一部分被去掉，而其余部分留了下来。

创造性与随机性

显然我们现在是在谈论创造性的机械化。但这是不是一种用词上的自相矛盾呢？很有点像，但并不真是。创造性是那种非机械的思维的本质，但每个创造性的行为又的确是机械的------对它也有解释，而且不会比对打嗝的解释更少。创造性的机械基础是我们所看不到的，但确实存在。反过来，一些非机械性的东西现在就已存在于灵活的程序之中了。这些程序不是一设计出来就有创造性的，但是一旦程序对其设计者来说不再是透明的，那时创造的过程就开始了。

一种流行的观念认为随机性是创造性活动中必不可少的一个成分。这或许是真的，但这与创造性的机械化------甚至程序化！------并无关系。世界上充满了随机性，一旦你把其中一部分反映到大脑中来，你大脑的内部结构就吸收了一些这种随机性。因此，符号的触发模式就可能把你引向看上去最具随机性的途径，这只不过是由于你是在和一个古怪、随机的世界打交道而已。因此一个计算机程序也可能这样。随机性是思维的一种固有特性，而不是某种需要``人工授精''的东西，不论是通过骰子、衰变的原子核、随机数表还是别的什么。认为人类的创造性依赖于这样一些带任意性的来源，那真是对人的一种侮辱。

被我们看作随机性的东西常常只不过是通过一个歪曲的``过滤器''来看某个匀称的东西的结果。一个漂亮的例子就是萨尔维亚蒂所提供的两种看π/4的值的方式。虽然说π/4的小数展开式并不是随机的，但人在需要时往往可以把它作为随机的来用：它是``伪随机的''。数学中充满了伪随机性------足够供给所有可能的创造型人才终生享用。

正像在科学研究中随时随地都有着``观念革命''一样，个体思维中也始终贯穿着创造性的行为。它们不仅仅存在于最高层，它们是无处不在的。其中大多数规模很小，而且可能以前曾发生过成千上万次------但它们与那些具有高度创造性的新行为是近亲。当今的计算机程序似乎并未产生大量小创造。它们所做的事大多数还都是很``机械的''。这仅仅证明事实上它们还没能准确地模仿我们的思维方式------但它们正在向这个目标靠拢。

也许使具有高度创造性的思想区别于常规思想的，是某种美、简单、和谐相结合的感觉。事实上，我有一个心爱的``元类比''，在其中我把类比看作和弦。想法很简单：表面上相似的东西常常并没有深入的联系，而有深入联系的东西又常常看上去毫不相干。这很自然地与和弦构成类比：物理上邻近的音符从和弦角度看相距甚远（如E-F-G），而和弦上邻近的音符在物理上又相距甚远（如G-E-B）。具有共同的概念骨架的概念以一种类似于和谐的方式产生共鸣，这种和谐的``概念和弦''如果在一张想象的``概念键盘''上测量，各个键可能离开很远。当然，仅仅把手伸开并按任意一种方式弹下去还是不够的------你可能会弹出一个七度或九度音程来！也许现在这个类比就像个九度和弦------宽倒是宽，但并不和谐。

从所有层次上取出模式

我选用邦加德问题作为本章的核心，这是因为当你研究这些问题时，能体会到我们人类从基因中继承的模式的难以捉摸的性质，这涉及所有的知识表示机制，包括嵌套的环境、概念骨架和概念对应、可滑动性、描述和元描述以及它们的相互作用、符号的裂变和聚变、多种表示（沿不同的维度和不同的抽象层次）、缺席预设值，以及其它方面。

目前可以断言，如果某个程序能从一个领域中取出模式，那么它在另一个领域中一定找不到对我们来说同样显然的另一种模式。你可能还记得我在第一章中提到过这句话的反面，即说计算机可以不顾疲劳地重复工作，而人则做不到这一点。例如，我们看看施德鲁对输入所会作什么反应。如果伊她·娥英想好这样一个句子：``拿起一个大红方块然后再放下''，然后她一次又一次地键入这个句子，那么施德鲁就会愉快地以同样方式一次又一次地做出反应，这正好像是如果有那么个人极有耐心地一次又一次键入``5+5''，他手中的那个忠实的加法器就一定会一次又一次地打印出``10''一样。人可不是这样。如果某个模式一次又一次地出现，他就会把它取出来。施德鲁不具有形成新概念或识别模式的能力：它没有关于多余的东西和总体概貌的意识。

语言的灵活性

施德鲁掌握的语言能力是非常灵活的------但这是在一定的限度之内。施德鲁能明白结构很复杂的句子，也能明白语义具有多义性的句子，只要这能通过对数据库进行检查来解决------但它无法掌握``模糊的''语言。以下述句子为例：``多少个木块彼此重叠构成了一个尖塔形？''我们立刻就能理解它，但如果把它按字面解释将是没有意义的。这也不是因为使用了某种习惯用语。``彼此重叠''是一个不准确的短语，然而它能很好地把预期的图像传送给一个人。没人会由于误解想象出一种荒谬情景：有两个木块，每个都叠在另一个的顶上。

语言的奇妙之处在于：我们很不准确地使用它，但仍能不出毛病。词汇在施德鲁那里好像是``金属的''，而在人们那里则像是``海绵的''或``橡胶的''甚至是``螺钉螺母式的''。假如词汇是螺丝钉和螺丝母，人们就可以把任何一个螺丝钉穿进任何一个螺丝母中：他们只须把前者硬压到后者中去就是了，就像在某幅超现实主义绘画中，一切东西都变软了一样。语言在人的手里几乎变成了液体，尽管它的成分很粗糙。

近来，人工智能在自然语言理解领域的研究已经在某种程度上离开了对隔离状态下的单个语句理解，更多地转向理解简单的儿童故事这样的领域。下面是一个广为传诵的童谣，它叙述了一个没完没了的真实生活中的情景：

一个人乘飞机去旅行，

倒霉，他一下掉到半空中。

万幸，他有一张降落伞，

倒霉，那玩艺儿一点儿不管用。

万幸，下面有个干草垛，

倒霉，一杆钢叉在上面戳。

万幸，他躲开了那杆钢叉，

倒霉，他也错过了干草垛。

它可以被无穷无尽地扩展下去。为表示这个无聊的故事，需要一个极其复杂的基于框架的系统，涉及到连带地激活对应于下列概念的框架：人、飞机、掉出去、降落伞、下落、等等等等。

智能与情感

请考虑下面这个简短但能说明问题的故事：

玛吉紧紧拉着一条线，上面系着她美丽的新气球。忽然，一阵风把它带走了。风把它吹向一棵树。气球在一只树枝上碰破了。玛吉哭个不停。\footnote{}

为了理解这个故事，你必须从字缝里读出许多东西来。例如，玛吉是个女孩子。这是个玩具气球，上面有一条可以让孩子拉着的线。它对大人来说可能并不美丽，但在孩子眼里的确如此。她是在室外。风所带走的``它''是指气球。风没有把玛吉和气球一起吹走，玛吉撒手了。气球碰上尖锐的东西就会破。它们一旦破了就无法弥补。孩子们喜爱气球，而且一旦它破了他们就会很伤心。玛吉看见她的气球破了。孩子们伤心时就会哭。``哭个不停''是指哭的时间很长，哭得很厉害。玛吉哭个不停是因为她的气球破了，她很伤心。

对于表层上所缺少的东西来说，上面写出的或许只是其中一小部分而已。一个程序要想懂得所发生的事情，就必须全部拥有这些知识。你可能会反对说：即使它在某种理智的意义下``理解''了所说的事情，它也绝不会真正理解，除非它也会哭个不停。一台计算机何时才能做到这一点呢？这是一种人本主义的观点，约瑟夫·魏增鲍姆在他的著作《计算机的能力和人类理性》中就涉及到了这一点。我认为这是一个重要的问题，事实上是个极其深刻的问题。遗憾的是，现在许多人工智能工作者出于各种原因不愿意认真考虑这种问题。但从某一方面来说，这些人工智能工作者是对的：现在考虑让计算机会哭未免为时过早，我们必须首先考虑计算机处理语言和其它问题的规则。到时候，我们就会发现我们已经直接面对那个更深刻的问题了。

人工智能任重道远

在人类的行为中，有时完全找不到受规则控制的迹象，好像人类的确不是受规则所控制的。但这是一种错觉------有点像认为晶体和金属具有刻板的内在规律，而液体和花朵则不然。我们将在下一章再回到这个问题上来。

大脑中工作着的逻辑过程本身可能更类似于一系列对象形文字的操作，这种文字抽象地看类似于汉字或玛雅语中的事件描述------唯一的例外是大脑中的逻辑过程的基本元素不仅是字词，而更像是句子或整个故事，在它们之中存在着相互联系，形成一种具有自己的规则的元逻辑或超逻辑。\footnote{}

对大多数专家来说，让他们生动地叙述最初是什么激励他们进入他们的领域，这往往是很困难的。他们甚至可能根本记不起来了。与此相反，某些站在外面的人却可能理解一个领域所特有的传奇故事，而且能清楚地把它表达出来。我想这就是上面引用的乌兰姆那段话吸引我的原因，因为它富有想象力地传达了人工智能事业的奇妙之处，而且显示了对它所抱有的信心。现在我们必须依靠我们的信心，因为前面的路还长着呢！

十个问题与推测

下面我要提出十个关于人工智能的``问题与推测''，以此作为本章的结束。我还没有大胆到称它们为``答案''的地步------它们只是我的个人见解。一旦我学到更多东西，或人工智能有了进一步发展，这些意见都可能以某种方式改变（在下面，``人工智能程序''一词意指那种远比现在运行的种种程序先进得多的程序，是一种``确实有智能的''程序。而``程序''、``计算机''这些词恐怕是带有更多的机械论内涵了，不过让我们还是依旧用这些词吧）。

\dialogname{问题：}是否计算机程序终将谱写出优美的乐曲？

\dialogname{推测：}会的，但不是近期的事情。音乐是种情感的语言，在程序能有像我们那么复杂的情感之前，它无论如何无法谱出什么优美的旋律。``赝品''------对已有音乐的句法进行肤浅的模仿------是容易达到的，但无论人们开始时是怎么想的，音乐表现远比句法规则所能捕捉到的东西要多得多。在相当长的一段时间之内，计算机作曲程序谱出的乐曲不会产生新的美感。下面我来把这一思想进一步展开。设想------我听到过这种暗示------我们很快就能有一种程序化了的``乐盒''台式的、可大批生产、二十块钱就能邮购，这``乐盒''中那没有生殖能力的电路产生的乐曲会让人觉得是出自肖邦或巴赫之手（如果他们能活到现在的话）。真要是这么想的话，那就是对人类精神的深度的一种荒唐可耻的低估。能有如此功能的``程序''必须得能自己走进这个世界，在纷繁的生活中抗争，并每时每刻体验到来自生活的感受。它必须懂得暗夜里的凉风所带来的喜悦与孤独，懂得对于带来温暖爱抚的手掌的渴望，懂得遥远异地的不可企及，还要能体验到一个人死去后引起的心碎与升华。它必须明了放弃与厌世、悲伤与失望、决心与胜利、虔诚与敬畏。它里面得能把诸如希望与恐惧、苦恼与欢乐、宁静与不安等等相对立的情绪混合在一起。它的核心部分必须能体验优美感、幽默感、韵律感、惊讶感------当然，也包括能精妙地觉察到清新的作品中那魔幻般的魅力。音乐的意义与源泉正是来自这些东西，也仅仅是来自这些东西。

\dialogname{问题：}情感是否能显明地在一台机器中程序化？

\dialogname{推测：}不能。那将是荒唐的。任何对于情感的直接模拟------比如``帕里''------都不可能达到人类情感的复杂性。后者间接地源于我们人类的心智组织，而程序或计算机可能以类似的方式获得它们自己的情感：作为它们的结构，即它们组织起来的方式的副产品------而非直接程序化的产物。因此，打个比方说，没人能写出一个``坠入爱河''的子程序，更不能写出什么``犯错误''子程序。``坠入爱河''是我们用于对一个复杂系统的一个复杂过程的描述。而且，这个描述并不单单来自系统中的任何一个单独的模块！

\dialogname{问题：}智能计算机是否能更快地做加法？

\dialogname{推测：}可能不会。组成我们自身的硬件的确在做着各种各样的运算，但这并不等于说``我们''所在的符号层知道如何完成同样的运算。让我们这样说：你无法把数字装进你的神经元中，以此来计算你的购货账单。幸亏你的符号层（也就是``你''）不能干预正在进行着你的思维活动的神经元------否则你非得搞昏了头不可。让我们重新解释笛卡尔的话：

``我思，故我无法进入我算的层次。''

为什么这对一个智能程序就不能同样成立呢？一定不能允许它干预正在进行着它的思维活动的电路------否则它非得搞昏了CPU不可。说正经的，一台能通过图灵测验的机器作起加法来可能和你一样慢，其原因也是类似的。它将不是把数字2仅仅表示成两个二进制位``10''，而是像我们那样把它表示成一个五脏倶全的``概念''，充满了各种联想。例如它的同义词有``两''、``一对''、``一双''，还有一大堆内心想象，如麻将牌上的点，数字``2''的形状，关于交替、偶数、奇数的概念，如此等等\ldots\ldots 背着这些``额外负担''前进，一个智能程序作起加法来会像树懒一样慢。当然，我们可以给它一个``袖珍计算器''，如果能这样说的话（或在它内部构造一个）。那样它就能很快做出回答，但它的表现和一个用袖珍计算器的人就完全一样了。在这种机器中将有两个彼此隔离的部分：一个可靠但呆板的部分和一个聪明但易错的部分。而且无法保证这样的合成系统是可靠的，正像无法保证人机合成系统必然是可靠的一样。因此如果你只需要一个正确的答案，最好单独使用那个袖珍计算器------不要往里边加智能！

\dialogname{问题：}是否会出现能击败任何人的下棋程序？

\dialogname{推测：}不会。可能会出现这样一些程序，它们在下棋时能击败任何人，但它们不是只会下棋的程序。它们将具有``通用的''智能，而且和人一样会发脾气。``你想下盘棋吗？''\,``不，我下棋都下腻了，我们来谈谈诗吧。''在你和那个能击败任何人的程序之间可能会出现这样的对话。这是因为真正的智能必然依赖于一种纵览全局的能力------即可以说是在程序中编入了一种``跳出系统外''的能力------至少大致能达到我们所具有的水平。一旦实现了这一步，你就无法再遏止这个程序了，它将越过这个特定的临界点，那时你将不得不面对面地与你造出来的东西打交道。

\dialogname{问题：}在存储器中是否有一个特定位置被用来存控制程序行为的参数，这样，如果找到并修改了它们，你就能使得该程序再灵一点或再笨一点或更有创造力或更喜欢足球？简而言之，你是否能在较低的层次摆弄它一下就可以对程序进行调整？

\dialogname{推测：}不行。它将对存储器中任意特定单元的改变毫不在意，正像我们身体中虽然每天要死掉几千个神经元（！），但我们却几乎没事一样。但如果你折腾得太厉害，就有可能伤害它，正像你不负责任地给一个人做了神经外科手术一样。在存储器中不存在``有魔力''的位置用以存储诸如程序的``智商''之类的东西。同样，那将是一个作为低层行为的结果而产生的性质，并非明确地存储在某个地方。对于像``短时记忆中所能保存的项目的数量''、``它喜欢物理学的程度''之类的东西也是一样。

\dialogname{问题：}你是否可能把一个人工智能程序的行为调整得像我或像你------或恰好介于我们两个中间？

\dialogname{推测：}不能。人工智能程序和人一样不能像条变色龙。它将依赖于它的记忆的稳定性，不能在不同的人格之间变来变去。那种想通过改变内部参数``调出一个新人格来''的想法暴露了对人格复杂性的一种可笑的低估。

\dialogname{问题：}一个人工智能程序是否有个``心脏''？还是说它只不过包含了一堆``无意义的循环和平凡操作的序列''（用马尔文·明斯基的话说\footnote{}）？

\dialogname{推测：}如果我们对它能像对一个浅潭那样一眼看到底，我们无疑只能看到``无意义的循环和平凡操作的序列''------而且肯定看不见什么``心脏''。现在关于人工智能有两种极端主义的观点：一种认为由于本质而又神秘的原因，人的心灵是无法程序化的。而另一种认为你只需去汇集适当的``启发式装置------多重优化器、模式识别技巧、规划代数、递归管理过程、诸如此类''\footnote{}，你就能得到智能。我觉得我自己在二者之间的某个位置上，相信我们将发现人工智能程序这个``水潭''很深而且黑洞洞的，无法一眼看到底。如果我们从上面看下去，那些循环是看不到的，正像今天对大多数程序员来说，载流电子是看不到的一样。一旦我们造出了一个能通过图灵测验的程序，我们将会看到其中有个``心脏''，尽管我们知道它并非实际存在。

\dialogname{问题：}人工智能程序能否具有``超智能''？

\dialogname{推测：}我不知道。我不清楚我们是否能理解``超智能''或与它发生关系，甚至不清楚这个概念是否有意义。例如，我们自己的智能是联系于我们的思维速度的。如果我们的反应速度被提高或降低了一个数量级，那我们可能已经发展出一套完全不同的概念，以此来描述这个世界。一种对世界有完全不同的看法的生物和我们之间可能不会有很多接触。我常常在想，如果存在着某种乐曲，比如说它们和巴赫的关系就像巴赫和民歌的关系一样------可以称它们为``巴赫的平方''------那我有可能理解它们吗？也许在我周围已经有这样的音乐，只不过我不能识别它们，就像狗不能理解语言一样。超智能是个很奇怪的想法。不管怎么说，我都不认为它是人工智能研究的目标，尽管如果我们终于达到了人类智能的水平，超智能无疑将成为下一个目标------不仅对我们是这样，而且对作为我们的同事的人工智能程序们也是如此。它们对人工智能和超智能也会具有同样的好奇心。很可能人工智能程序会对人工智能的一般问题怀有极大的兴趣------这是可以理解的。

\dialogname{问题：}那你的意思似乎是说人工智能程序实际上和人完全一样了。难道就一点差别也没有吗？

\dialogname{推测：}人工智能程序与人之间的差别或许要大于大多数人之间的差别。几乎无法想象人工智能程序所寓的``躯体''会对它不产生深入的影响。因此，除非它有一个人类躯体的忠实得惊人的复制品------为什么它要这样？------否则它对于什么是重要的、什么是有趣的等问题或许会有极其不同的看法。维特根斯坦曾作过一个有趣的评论，``如果狮子会说话，那我们也无法理解它。''这使我想起了卢梭的一幅画：在月光皎洁的沙漠上，一头温顺的狮子守在睡着了的吉普赛人身旁。但维特根斯坦是怎样知道的呢？我猜任何人工智能即使是我们所能理解的，那也一定是看上去和我们很不同。出于这种原因，我们到时候可能会难于断定和我们打交道的是一个真正的人工智能程序，还是说它只不过是一个``古怪''程序而已。

\dialogname{问题：}一旦我们造出了一个智能程序，我们是否就能理解什么是智能、意识、自由意志和``我''了？

\dialogname{推测：}在某种程度上是这样------这完全取决于你说的``理解''是什么意思。首先，在某个关键性的层次上看，我们中的每个人大概都对这些东西有尽可能好的理解。这就像是听音乐。是不是因为你已经把巴赫拆散了，你就算真的理解他了呢？或者说当你感到全身的每一根神经都很兴奋，那时你就算理解他了？我们是否理解光速在每个惯性参照系中保持恒定到底是怎么回事？我们可以算出它来，但世界上没有一个人具有真正的相对论性的直觉。而且大概永远没人能以一种直观的方式理解智能与意识之谜。我们中的每一个都能理解``人''，而这或许也就是你所能达到的。

\begin{enumerate}
\item
  \phantomsection\label{Chapter19.xhtml_B_1}
\item
  \phantomsection\label{Chapter19.xhtml_B_2}
\item
  \phantomsection\label{Chapter19.xhtml_B_3}
\item
  \phantomsection\label{Chapter19.xhtml_B_4}
\item
  \phantomsection\label{Chapter19.xhtml_B_5}
\item
  \phantomsection\label{Chapter19.xhtml_B_6}
\item
  \phantomsection\label{Chapter19.xhtml_B_7}
\end{enumerate}

\phantomsection\label{Dialog20.xhtml}{}

\chapter{树懒卡农}

这一回，我们看到阿基里斯和乌龟去拜访树懒。在他们谈话的时候，树懒倒吊在树上，阿基里斯坐在地上，而乌龟在下面的水塘里泅水。

\dialogname{阿基里斯：}小懒，我跟你说说我跟龟兄进行的那场滑稽至极的赛跑怎么样？

\dialogname{树懒：}请吧。

\dialogname{阿基里斯：}这件事在这一带已经尽人皆知了。我想它甚至已经被人写出来了，是芝诺写的。

\dialogname{树懒：}听上去叫人高兴。

\dialogname{阿基里斯：}是这样。你知道，龟兄在我前面的地方起跑。他比我靠前好长一截路，然而------

\dialogname{树懒：}你追上他了，不是吗？

\dialogname{阿基里斯：}没错儿------我长着一双飞毛腿嘛。我先是以一定的速度缩短我们两个之间的距离，然后很快就超过了他。

\dialogname{树懒：}差距越来越小，所以你赢了。

\dialogname{阿基里斯：}正是。哦，瞧------龟兄把他的小提琴带来了。我可以试拉两下吗，龟兄？

\dialogname{乌龟：}请别。听上去叫人心烦。

\dialogname{阿基里斯：}噢，很抱歉。可我现在特别想听音乐。我真不知道这是怎么回事。

\dialogname{树懒：}我记得你会弹钢琴，阿基。

\dialogname{阿基里斯：}是的。我马上就来两下。我刚才只是想补充说我跟龟兄在那之后还有过一场``比赛''呢。不幸的是在那场比赛里------

\dialogname{乌龟：}你追不上我了，是吗？差距越来越大，所以你臝不了。

\dialogname{阿基里斯：}完全正确。我想这场比赛甚至也已经被人写出来了，是刘易斯·卡罗尔写的。现在，小懒，我要接受你的提醒，弹弹钢琴。可我钢琴弹得太糟了，我不敢说我能弹好。

\dialogname{树懒：}你应该尝试。

（阿基里斯坐下来，开始弹一支简单的曲子。）

\dialogname{阿基里斯：}唉------听起来很古怪。不该是这种调儿啊！有什么东西不对头。

\dialogname{乌龟：}我记得你不会弹钢琴，阿基。你不应该尝试。

\dialogname{阿基里斯：}简直像镜子里的什么钢琴。高音键在左边，低音键在右边。所有曲子都成了反的，好像倒悬着似的。谁会鼓捣出这么一个荒唐的玩艺儿来？

\dialogname{乌龟：}这很符合树懒的本性。他们是------

\dialogname{阿基里斯：}哦，我知道------挂在树枝上------当然是倒悬着。这架树懒钢琴要是用来演奏某些赋格和卡农中的倒奏曲子倒挺合适。但是要想学会挂在树上弹钢琴可太难了。你必须竭尽全力。

\dialogname{树懒：}这很不符合树懒的本性。

\dialogname{阿基里斯：}对，我认为树懒是愿意过得舒适些的。他们无论干什么，都比正常速度慢得多。此外，他们还倒悬着。这是种多么奇特的生活方式啊！说到具有倒个儿和减速这类特征的事物，倒让我想起了《音乐的奉献》中《反向进行的增值卡农》。在我那本谱子里，三行谱线前各有``高''、``中''、``低''三个字。我不明白这是什么意思。不过，不管怎么说，我认为巴赫把这一技巧运用得很娴熟。你以为呢，龟兄？

\pandocbounded{\includegraphics[keepaspectratio]{./media/OEBPS/Images/img_133.png}}

图133．``树懒卡农''

选自巴赫的《音乐的奉献》，乐谱是用唐纳德·伯德的程序``斯马特''印制的，美术字由陈春光设计

\dialogname{乌龟：}这在他是很轻松自如的。至于``高''、``中''、``低''三个字，你一定猜得出它们指什么。

\dialogname{阿基里斯：}我想是代表``女高音''、``女中音''、``女低音''，因为三个声部的作品常常用这种方式构成。你说呢，小懒？

\dialogname{树懒：}它们不是------

\dialogname{阿基里斯：}等一下，我们三个现在呆的地方很有意思：你，小懒，位置最高，挂在树上；我在中间，坐在地上；龟兄最低，泡在下面的小塘里。哎------你怎么出来了，龟兄，还穿上了你的外套？你看去很疲倦。怎么啦？

\dialogname{乌龟：}精疲力尽了。我得先走一步。\kaitismall{（步履疲惫地走远了。）}

\dialogname{阿基里斯：}他真可怜------看来他确实顶不住了。他一上午都在颠来颠去的。为了同我再较量一次，他正接受训练呢。

\dialogname{树懒：}这在他可是竭尽了全力。

\dialogname{阿基里斯：}是的，但是没用。他也许能赢一只树懒\ldots\ldots 可要想赢我------决不可能！哎，你不是要告诉我那三个字的意思吗？

\dialogname{树懒：}至于``高''、``中''、``低''三个字，你根本猜不出它们指什么。

\dialogname{阿基里斯：}要是我刚才猜得不对，这倒更让我好奇了。算啦，我们先不谈这个了。我真有点想吃东西了，我上次在老蟹那儿吃过一次油炸土豆片，听说你也会做？据说那东西有一种神奇的功效：当你觉得懒洋洋时，吃上它几片，马上就能------

\dialogname{树懒：}精力充沛了。

\dialogname{阿基里斯：}对，他们是这么说的。可是，像你这样常吃它的人，说话、办事的速度居然只是以慢吞吞而著称于世的龟兄的一半，这不免让我对这种传说产生怀疑。不过，不管怎么说，它还是挺好吃的。炸的时候有什么要领吗？

\dialogname{树懒：}你得后退一步。

\dialogname{阿基里斯：}对，这样就不会被溅出的热油烫着了，是吧？我现在就去试试，我想我会做成功的，只可惜龟兄也许尝不到了。

\phantomsection\label{Chapter20.xhtml}{}

第二十章

\chapter{怪圈，或缠结的层次结构}

机器能具有创造性吗？

在第十八章，我描述过阿瑟·塞缪尔的精彩的跳棋程序，它可以击败它的设计者。据此，听一听塞缪尔自己关于``计算机与创造性''问题的感想还是很有趣的。下面引文摘自塞缪尔写于1960年的一篇反驳维纳的文章。

我坚信机器不可能具有维纳的文章所暗示的那种创造性。按照他的意思，``机器一定能够超越其设计者的某些局限性。在这种情况下，它们可能既是有效的又是危险的。''

\ldots\ldots{}

机器不是精灵，它不是靠魔法工作，它不具有意愿，而且与维纳的意见相反，不会生成事先没有放进去的东西。当然，这样说不包括偶然出故障的情况。

\ldots\ldots{}

机器所表现出来的``意向''是事先已明确化的人类程序员的意向，或是依照程序员所指定的规则从这些意向中导出的子意向。我们甚至可以像维纳那样期望更高的抽象层次，在那些层次上程序不仅修改子意向，而且也修改导出规则，或修改其修改规则的方式，如此等等。在那种层次上甚至一台机器会设计和构造另一台具有更强能力的机器。但重要的是，除非机器已经得到如何去做这些事的指令，否则它\heititext{不会也不可能}（黑体是他用的）去做这些事。这里有一条逻辑上永远存在的鸿沟，一边是实现人的愿望的过程中任何一种极力的扩展和精心的构造，另一边是在机器中开发属于它自己的意愿。如果不相信这一点，则要么就相信魔法，要么就相信人的意愿只是一种错觉，而人的行为是和机器一样机械的。也许维纳的文章和我的反驳两者都是被机械地决定了的，但我拒绝这么想。\footnote{}

这让我想起了卡罗尔的对话（《二部创意曲》）。我来解释一下为什么这样说。塞缪尔反对机器自我意识（或愿望）的论点是建立在这样的见解之上的：对意愿的任何机械化都需要一种无穷回归。类似地，卡罗尔的乌龟争辩说，任何一步推理，无论多么简单，如果不援引更高的层次上的规则来证实其合理性，都不能进行。但这种证实本身不是一个推理步骤，这样你不得不求助于更高层的规则，如此等等。结论：推理包含着一个无穷回归。

显然，乌龟的论证中有些不大对头的东西，而且我相信塞缪尔的论证也有类似的问题。为了显示两者的相似性，我现在暂时充当一下魔鬼的辩护士来``帮助魔鬼''，故意唱唱反调。（众所周知，上帝帮助自助者，因此想来魔鬼就只帮助那些不自助的人。魔鬼帮助他自己吗？）以下是我站在魔鬼的立场上从卡罗尔的对话中推出的结论：

``推理是不可能的''这一结论对人类不适用，因为无论有多少更高的层次，我们的确是在进行许多推理，这对任何人都是习以为常的。这表明我们人类的活动不需要规则：我们是``非形式的系统''。但从另一方面来说，上述结论可以作为一种论证，说明对推理进行机械化解释是不可能的，因为任何机械系统势必明确地依赖于规则。因此，除非有元规则告诉它何时使用规则，有元元规则告诉它何时使用元规则，如此等等，否则它不可能开始工作。我们可以断言，推理能力永远不可能机械化。它是为人类所独有的能力。

魔鬼辩护士的这种观点有什么错误呢？显然是如下的假设：如果没有规则告诉机器如何去做某件事，它就做不了。事实上，机器能和人一样容易地避开乌龟那笨拙的责难，而且理由是同样的：机器和人两者都是由硬件构成的，而硬件可以按照物理学定律完全独立地运行。这里根本不需要``允许你使用规则的规则''，因为最底层的规则------关于它们不再有``元规则''------是嵌入于硬件之中的，它们的运用无需经过许可。由此可见，卡罗尔的对话说到底并没讲出人和机器的任何区别（而的确，推理是可以机械化的）。

关于卡罗尔对话就谈这些。下面回到塞缪尔的论证。其要点------如果允许我夸张一下------是这样的：

没有任何计算机真正``想要''去做任何事，因为它是被别的什么人编好了程序的。只有当它能从零开始为自己编程序------这是荒唐的------它才可能有自己的关于愿望的意识。

在他的论证中，塞缪尔重蹈了乌龟的覆辙，只不过是把``推理''换成了``意愿''。他的意思是说：在任何愿望的机械化的背后，必定存在一个无穷回归------或者更糟些，存在一个封闭循环。如果这是计算机不具有自己意愿的原因，那么人是怎么回事呢？从同样的评价标准中可以推出：

除非一个人设计他自己并选择他自己的需要（同时对自己需要的选择进行选择，等等），否则他不能算是具有自己的意愿。

这迫使你暂停片刻，想想你自己关于``具有意愿''这一意识是来自何方。除非你是一个唯灵论者，否则你大概会说它来自于你的大脑一块你既没有设计也没有选择的硬件。然而，你需要某些东西，而非另外一些东西，这种意识并不因此而减弱。你不是一个``自我编程的物体''（不管那是种什么东西），但你仍然具有关于愿望的意识，它发源于你的心智的物理基质之中。同样，尽管事实上不会有从``乌有之乡''自发地出现于存储器中的那种魔术般的程序（一个``自我编程的程序''），机器在某一天仍可能具有意愿。它们拥有意愿的原因和你相同------都是因为许多层次上硬件和软件的组织与结构。由此可见，塞缪尔的论证说到底并没有讲出机器和人的任何区别（而的确，意愿是可以机械化的）。

每个缠结的层次结构下面都有一个不受干扰的层次

紧接在《二部创意曲》之后，我说过本书的一个核心论题将是``言语和思维是否遵从形式规则？''。本书的主要努力之一就是要指出心／脑的多层性，并且我已试图表明对此问题的最终回答是``是的------条件是你要下到最底层------硬件------去发现这些规则。''

现在，塞缪尔的陈述引出了一个我将要继续探讨的概念。这是说：当我们人在思维时，我们的确是在改变我们自己的思想规则，同时我们改变那些修改规则的规则，如此下去------但这些都可以说是``软件规则''。最底层的规则并未改变。神经元始终以同样一种简单的方式活动。你不能``想服''你的神经元去改变活动方式，虽然你可以使你的头脑改变思想的形式和主题。正像《前奏曲，蚂蚁赋格》中的阿基里斯一样，你可以达到你的思想，但不能达到你的神经元。各个层次上的软件规则都可以改变，而硬件规则不变------事实上，软件的灵活性来自于硬件的稳固性！这里没有任何悖论，不过是一个关于智能机制的基本而又简单的事实。

在能进行自我修改的软件和不受侵犯的硬件之间的这种区别，正是我要在这最后一章中加以展开的主题。我要从中发展出一组变奏，其中有些变奏可能看起来很牵强，但我希望等到我完成这个循环，回到大脑、心智和自我意识等问题时，你会发现在所有这些变奏中都有一个不变的内核。

我在本章中的主要目的是讨论一些想象，它们帮助我形象化地设想意识是如何从错综复杂的神经元中产生。我将讨论一些难以捉摸的直觉，我希望这些直觉是有价值的，也许对别人能有所帮助，使他们能更清晰地表述他们自己关于心智活动的想象。我仅仅希望我自己思想中的关于思想和想象的模糊想象能够促成其他人思想中的关于思想和想象的更清晰的想象。

一种自我修改的棋

第一个变奏关系到这样一种棋赛：当轮到你走的时候，你可以修改规则。以国际象棋为例，显然规则始终是同样的，而每一步只是改变了局势。但让我们发明一种变形的国际象棋，在下这种棋时，一旦轮到你走，你可以走一步棋，或是对规则进行修改。可是怎么改呢？随便吗？你能把它改成跳棋吗？显然这样乱来将是毫无意义的，必须有某些限制。例如这样一种修改：允许你重新定义马的走法，让它不再走``横1竖2''（当然还有``竖1横2''），而改走``横m竖n''（或``竖m横n''），这里m和n是任意的自然数。当轮到你走棋时，你可以对m或n加1或减1。因此它可以从``1-2''到``1-3''到``0-3''到``0-4''到``0-5''到``1-5''到``2-5''\ldots\ldots 然后将会有一些规则来重新定义象的走法以及其它棋子的走法。还可以有些规则被用于在棋盘上添增新线和删除旧线\ldots\ldots{}

现在我们有两层规则：一层是说明如何移动棋子的，另一层是说明如何修改规则的。这样就有了规则和元规则。下一步是很显然的：引进用于修改元规则的元元规则。但怎样才能做到这一点，这可并不显然。我们之所以能很容易地建立移动棋子的规则，是因为棋子的移动是在一个形式化空间中进行，即棋盘。如果你能为规则和元规则设计一种简单的形式化表示法，那处理它们就会象形式化地处理符号串------甚至可以像处理棋子一样。若把问题推向极端，你甚至可以用辅助棋盘上的一些位置来表示规则和元规则。那时，棋盘上的任意一种局势既可以被看成一盘棋，也可以被看成一组规则，还可以被看成是一组元规则，这取决于你用哪种方式解释它。当然，两名棋手必须在这种表示法的解释上达成协议。

这时，我们就可以有任意多个相邻的棋盘：一个用于下棋、一个表示规则、一个表示元规则、一个表示元元规则------如此上去，你愿意有多少都行。当轮到你走时，你可以在除最顶上那个棋盘之外的任何一个棋盘上走一步，这一步所遵循的规则是位于该层次结构中上一层的棋盘所提供的。无疑，两个棋手会由于几乎所有事情------虽然并非每件事情------都可以改变而被搞得晕头转向。根据定义，最顶上的棋盘是不能改变的，因为你没有规则来说明如何修改它。它是``不受干扰的''。还有一些不受干扰的东西：关于不同棋盘的解释方式的协议、轮流走棋的约定，每人每次可以改变一个棋盘的约定------如果你仔细检查这个想法，还能发现更多。

现在就有可能进一步取消使这种定位得以建立的支柱。一步一步地来\ldots\ldots，我们首先把这一整套棋盘压缩成一个棋盘。这是什么意思呢？这就是说对于这个棋盘可以有两种解释方式：（1）作为要移动的一些棋子；（2）作为一些移动棋子的规则。轮到你走时，你移动了棋子------你必然也就修改了规则！这样，规则就会不断修改它们自己。这使我们想到了``印符遗传学''------或因此想到了真正的遗传学。在棋局、规则、元规则、元元规则之间的差别已经消失。一度整齐、清楚的层次结构变成了一个怪圈，或缠结的层次结构。棋的着法改变了规则，而规则又决定了棋的着法，鸡生蛋，蛋生鸡\ldots\ldots 虽然仍有不同的层次，但``低层''和``高层''之间的区别已经被抹掉了。

现在，有一部分过去不受干扰的东西已经成为可修改的。但仍有很多不受干扰的东西。正像以前一样，在你和你的对手之间有一些约定，据此可以把棋盘解释成一组规则。还有一个约定是关于轮流走棋的------或许还会有其它的隐含约定。因此必须注意，关于不同层次的概念以一种出人意料的方式存活下来了。有一个不受干扰的层次------在其上是关于解释的约定，还有一个缠结的层次------在其上是那个缠结的层次结构。因此这两层仍可构成一个层次结构：不受干扰层控制着缠结层上发生的事件，但缠结层没有也不可能影响不受干扰层。尽管缠结层本身是个缠结的层次结构，它仍然被存在于它之外的一组约定所控制。而这就是关键所在。

正如你无疑已经想到的那样，没有什么可以阻止我们去做那种``不可能做到''的事------那就是说，根据棋盘上的位置把关于解释的约定本身作为修改对象，以此把不受干扰层和缠结层缠结在一起。但为了进行这样一种``超缠结''，你必须和你的对手达成某些与这两个层次有关的进一步的协议------而这样做就创建了一个新层次------一个新的不受干扰的层次，它位于那个``超缠结''层之上（你愿意说``之下''也行）。这一过程可以继续进行下去。事实上，这里所做的``跳跃''很像在《生日大合唱哇乌阿乌阿\ldots\ldots》中------或在对TNT的各种改进形式反复进行哥德尔化的过程中------所描绘的那样。每当你以为你已经达到了终点，总会有``跳出系统''这一主题的某种新的变奏出现，而想发现它则需要一种创造力。

再谈作者三角形

但我并不想进一步讨论在自我修改的国际象棋中会出现的那种极其复杂的缠结现象。我举这个例子只是想形象地说明下述观点：在任意一个系统中总存在某个``受保护的''层次，它不会被其它层上的规则所攻击，尽管这些规则之间的相互作用可能是复杂地缠结在一起的。在第四章中有个有趣的谜，它在一个略有不同的环境中表示了同样的想法。也许它会打你个冷不防：

有三个作者------Z、W和G。现在Z仅仅存在于W所写的一部小说中。相似地，W仅仅存在于G所写的小说中。更怪的是，G也仅仅存在于小说中------当然是Z写的。那么，这样一个``作者三角形''真是可能的吗？（见图134）

\pandocbounded{\includegraphics[keepaspectratio]{./media/OEBPS/Images/img_134.png}}

图134．``作者三角形''

这当然是可能的。但这里耍了一个花招\ldots\ldots 这三个作者Z（芝某）、W（乌某）和G（吉某）全都是侯某所写的另一部小说中的角色！你可以把芝-乌-吉组成的三角形看成一个怪圈，或缠结的层次机构，但作者侯某是处于发生这个缠结的空间之外的------作者侯某某位于一个不受干扰的空间之中。虽然芝某、乌某和吉某都能直接地或间接地彼此相互作用，而且能在他们各自的小说中互相诽谤，但他们都无法接触到侯某的生活！他们甚至无法想象到他------正像如果你是一本小说中的角色，那你也无法想象该书的作者一样。如果我想描写作者侯某，我只能在书外的某个地方来表示它。这当然会带来问题，因为要描写一个东西就必须把它放进书中\ldots\ldots 不管怎么说，侯某的确是在芝某、乌某和吉某所处的世界之外的，而且应当被表示成这样。

艾舍尔的《画手》

\pandocbounded{\includegraphics[keepaspectratio]{./media/OEBPS/Images/img_135.jpg}}

图135．画手

艾舍尔作（石版画，1948）

在我们的主题上的另一个经典性的变奏就是艾舍尔的作品《画手》（图135）。在这张画中，一只左手在画一只右手，而与此同时，右手又在画左手。又一次地，平时被看作组成了层次结构的那些层次------画的和被画的------彼此重叠，构成了一个缠结的层次结构。但这恰恰说明了本章的一个主题，因为在全部这些东西之后隐藏着一直未画出但正在画的手，它属于艾舍尔，左手和右手二者的创作者。艾舍尔处于这两只手所在的空间之外，而在我为这幅画所作的示意图中（图136），你可以清楚地看到这一点。在这张为艾舍尔的画所配的示意图中，你在上部可以看到那个怪圈，或缠结的层次结构；同时，你又能在它下面看到那个不受干扰的层次，下一层使上一层成为可能。我们可以使艾舍尔这幅画进一步``艾舍尔化''，方法是拍一张一只手正在画它的照片，如此下去。

\pandocbounded{\includegraphics[keepaspectratio]{./media/OEBPS/Images/img_136.png}}

图136．艾舍尔的《画手》的抽象示意图

上面部分似乎是个悖论，下面部分是它的解

大脑和心智：一个神经元缠结支持一个符号缠结

现在我们可以把这个问题和大脑以及人工智能程序联系起来了。在我们的思维过程中，符号激活其它符号，所有符号相互作用而形成异层结构的形式。更进一步来说，符号可以造成相互间的内部变化，就像程序作用于其它程序一样。由于符号组成了缠结的层次结构，这就造成了一种假象，即``不存在不受干扰的层次''。之所以有人认为不存在这样的层次，是因为它处于我们的视野之外。

如果可能画出整个这幅图像的示意图，那么其中将有一处是个庞大的符号森林，符号之间由缠结的线彼此相联，就像热带丛林中的藤蔓一样------这将是其顶层，是思维实体在其中流来流去的那个缠结的层次结构。这就是难以捉摸的``心智''层：类似于``左手''和``右手''。在这幅示意图的下端，类似于那个不可见的``原动力''艾舍尔，将会是无数神经元的表示。那个使得上面的缠结得以产生的``不受干扰的基质''。有趣的是，下面这一层本身也是个不折不扣的``缠结''------几十亿个细胞被上千亿条轴突联结在一起。

在这种有趣的情况下，一个由符号构成的软件缠结被一个由神经元构成的硬件缠结支持着。但只有符号缠结是一个缠结的层次结构。神经元缠结只是一个``简单''缠结而已。这种差别就像我在第十六章中提到的怪圈和反馈之间的差别一样。一旦一个你原以为具有规整的层次结构的系统令人惊讶地以一种扰乱了层次的方式绕了回来，一个缠结的层次结构就出现了。``惊讶''这个因素还是很重要的，正因为如此我才说怪圈是``怪''的。一个简单缠结，像反馈，并没有涉及到扰乱预先设定的层次划分。例如你在浴室里可以先用右手洗左手，再用左手洗右手，在这种情况下没有什么可奇怪的。艾舍尔并非无缘无故地选择去画画手的手！

像两只手相互洗这样的事件随时随地都在发生，我们从不特别注意它们。我对你说了几句话，然后你又回答了我几句话。这是悖论吗？不，我们的相互感知并不需要以一种层次结构作为前提，因此这里就没有什么奇怪的东西。

在另一方面，在语言谈论它自己的时候，不论这种谈论是直接的还是间接的，它的确构成了一个怪圈。在这里，某些系统之中的东西跳了出去并且作用于系统之上，仿佛它是在系统之外似的。使我们感到麻烦的大概是某种界说不清的拓扑学上的不对头：内外差别被搞乱了，就像在那个被称为``克林瓶''的著名图形中一样。尽管这个系统是一种抽象，我们的心智的确使用着带有某种精神拓扑的空间想象力。

让我们回到符号缠结上来。如果我们只看着它，忘掉那个神经元缠结，那我们似乎看到了一个自我编程序的物体------正像如果我们看着``画手''，由于忘记了艾舍尔的存在，就会产生一种错觉，好像我们看到了一幅自己画自己的图画。对这幅画而言，这种情况不太可能发生------但对于人以及他们观察自己心智的方式来说，这正是通常所发生的情况。我们``觉得''我们是在给自己编程序。的确，我们无法产生别样的感觉，因为我们被屏蔽于底层------也就是神经元缠结------之外了。我们的思维好像是在自己的空间中运行，创造新思想，而且我们从未注意到任何神经元会给我们帮助！但这是意料之中的事。我们没办法。

一种类似的双关现象可能出现在某些Lisp程序之中，这些程序被设计得可以进入并修改它们自身的结构。如果你在Lisp层上观察它们，你会说它们是在自我修改，但如果你换一个层次，把Lisp程序看成Lisp解释程序（见第五章）的数据，那么实际上只有一个解释程序在运行，而且所进行的修改只不过是对某些数据的修改。Lisp解释程序本身是不被修改的。

你应如何描述这样一个缠结的情境，这取决于你在开始描述前先要往回走多远。如果你走得足够远，那你往往可以发现能把你引向无缠结事物的线索。

政府中的怪圈

一个极有趣的有层次缠结的领域是美国政府------尤其是法院。在一般情况下，你可以想象两个辩护人就案情在法庭上展开辩论，然后由法院对案件进行裁决。法院和辩护人是处在不同层次上的。但如果法院本身卷入了一场法律纠纷，那怪事就可能开始出现了。通常存在一个处于纠纷之外的更高级的法院。即使两个低级法院卷入了某种奇怪的争斗之中，每家都声称对另一家有裁判权，这时仍有某个更高级的法院居于事外，这在某种意义下类似于我们在那种变了形的国际象棋中所讨论的不受干扰的解释协定。

但如果连最高法院自己也卷入了法律纠纷，不再有更高级的法院，这会造成什么结果呢？这和``水门事件''中发生的麻烦差不多。总统威胁说他只服从最高法院的``决定性裁决''------然后又声称他有权确定什么是``决定性的''。现在看来这一威胁并没起作用。但假如当时真的管用了，它就会在政府的两个层次之间触发一场尖锐的对抗。二者之中的每一方都可以通过某种方式合法地宣称它是``高于''另一方的------这时该找谁来确定谁是谁非呢？找国会并不解决问题，因为国会可能会命令总统服从最高法院，而总统仍有可能拒绝，并宣称他根据法律有权在特殊情况下不服从最高法院（和国会！）。这将构成一起新的法律案件，并把整个系统搞乱，因为这实在太令人惊讶、太纠缠不清------太怪了！

具有讽刺意味的是，一旦你这样碰到了房顶------也就是说无法跳出系统寻求更高层的权威，那时唯一办法是求助于那些看上去没有用规则定义清楚的力量，而它们才是更高层规则的唯一来源------这就是低层规则，在上述例子中就是指社会的普遍反应。应当记住，在像美国这样的社会里，法律制度在某种意义上是被千百万人所共同承认的一种文雅的姿态------它可以很容易地被无视，就像河水会漫过堤坝一样。那时就会出现某种程度的无政府状态。但无政府状态也有它自身的规则，而且一点也不比文明社会少：只不过它们的作用方式是自下而上的，不是自上而下的。一个研究无政府状态的人可以设法去发现某些规则，它们控制了无政府状态随时间的发展过程，而且很可能确实存在这样一些规则。

不妨用一个物理学中的现象做类比。正像本书前面提到过的那样，平衡态下的气体遵从某些简单的规律，这些规律把它们的温度、压力和体积联系起来了。但是，气体有可能违反这些规律（这就像总统有可能违反法律一样）------条件是它并非处于平衡态。在非平衡状态下，一个物理学家如果要描述所发生的情况，他只能求助于统计力学------即转到某个非宏观的描述层次。因为对气体行为的最终解释永远是在分子水平上完成的，正像对社会的政治行为的最终解释永远是在``基层水平''上完成的一样。非平衡态热力学这个领域中的研究目标就是发现一些宏观规律，以此来描述偏离平衡态的气体（以及其它系统）的行为。它和政治学中研究无政府社会所遵循的规律的那个分支相类似。

美国政府中还有另外一些稀奇古怪的缠结，诸如联邦调查局调查它本身的错误、一个县长负担着县长的责任的同时又被送入监狱、议会程序规则的自我应用，如此等等（在其它国家的政坛上也有类似的情形）。我所听说过的一个最稀奇古怪的法律案件牵连到一个自称有特异功能的人。实际上，他自称能用他的特异功能测验人的特征，这样就可以协助律师挑选陪审员。那么如果有一天这种``特异功能''被用于对他自己的审判，会出现什么情况呢？这对一个坚信``超感官知觉''存在的陪审员会产生什么影响？他会在多大程度上觉得受到了特异功能的影响（不论这种特异功能是真是假）？开拓这个领域的时机已经成熟------这里充满了自我实现的预言。

与科学和鬼话有关的缠结

谈到特异功能和超感官知觉，生活中充满了怪圈的领域就是伪科学。伪科学的所作所为总是使正统科学的许多标准信念和过程成为疑问，因此对科学的客观性构成了挑战。一些与已有的证据解释方式相竞争的新方式被提出来了，但你该怎样评价一种解释证据的方式呢？这难道不正是又回到了客观性问题，只不过是在一个更高的层次上吗？当然如此。刘易斯·卡罗尔那个无穷回归的悖论以一种新的形式出现了。乌龟可能会争辩说，如果你想说明A是个事实，你需要证据B，但你怎样确定B是A的证据呢？为了说明这一点，你需要元证据C。为了说明这个元证据的有效性，你又需要元元证据------如此令人厌烦地束缚下去。尽管有这种论点，人们还是对证据具有一种直觉。这是因为------我又要重弹老调了------人们在脑中具有预先构成的硬件，其中包括某些对证据进行解释的基本方式。我们可以在此基础上构造并积累新的证据解释方式，我们甚至要学习应该以什么方式、在什么情况下摆脱我们最基本的证据解释机制，例如我们如果想看出魔术的诀窍，就必须这样做。

关于证据的二难推理的具体例子往往和伪科学的许多现象一起出现。例如，超感官知觉常常在实验室之外得到证明，可是一旦进入实验室，它就神秘地消失了。对这种情况的标准科学解释是：超感官知觉不是一种真实现象，经不住严格的审查。但是，部分相信超感官知觉的人（绝非全部）有一种奇特的反击方式。他们说：``不，超感官知觉是真实的，只不过在人们企图对它进行科学观察时就会跑掉------它与科学世界观的本性相互排斥。''这是一种无耻的手段，我们可能该管它叫``矛头向上''，这就是说，不是怀疑手头的东西，而是去怀疑那些可靠性更高的理论。那些相信超感官知觉的人想说明：出了毛病的并不是他们的想法，而是科学的信念系统。这是一种极其狂妄的观点，除非它能找到无可置疑的证据，否则我们就要对它表示怀疑。但那样我们又遇到了同样的问题：我们在谈论``无可置疑的证据''，仿佛所有的人都能对它所指的东西保持一致意见似的！

证据的本质

在第十三章和第十五章提到过的萨哲杜-辛普利奇奥-萨尔维亚蒂缠结，为证据评价的复杂性提供了另一个例子。萨哲杜试图在辛普利奇奥和萨尔维亚蒂的相反意见之间找到某种客观的妥协，如果这是可能的话。但妥协并不总是可能的。你怎样才能``公平地''在正确和错误之间达成妥协呢？怎样才能在公平和不公平之间达成妥协？这些问题在关于日常事物的争论中一次又一次地以各种不同形式出现。

是否有可能定义什么是证据？是否有可能制订一些法律来说明应当如何从情境中生成意义？或许不能，因为任何刻板的规则无疑都会有例外，而不刻板的规则又不能算是规则。就算我们有个智能化的程序也无济于事，因为作为一个证据处理器，它会和人一样容易犯错误。因此，如果证据确实是这样一种难以捉摸的东西，那我为什么还要对解释证据的新方式提出警告呢？我是不是自相矛盾了？在这种情况下，我不这样认为。我的感觉是我们可以划出一些界线，在其外边可以进行一种有机综合。但这时不可避免地会出现某些判断和知觉------这些东西是因人而异的。它们也是``因人工智能程序而异''的。最后，还会有一些复杂的标准，它们可以被用来确定一个评价证据的方法是否得当。这包括根据这种推理方式所得到的想法的``可用性''。导致生活中的新事物的思维方式在某种意义下被认为是``有效的''。但``可用''这个词带有太强的主观色彩。

我的感觉是，我们确定事物的有效性和真实性的过程是一门艺术，它深深地依赖于一种对美和简单性的感受力，正像它依赖于坚固的逻辑或推理或任何其它能被客观地加以形式化的原则一样。我既不是说（1）真理是一种幻想，也不是说（2）人的智能本质上是无法程序化的。我是在说（1）真理是如此难以捉摸，以致于无法被任何人或人群所彻底获得，以及（2）对人工智能来说，一旦它达到了人类智能的水平------或甚至超过了这一水平------它仍将被艺术、美感及简单性等问题所折磨，而且会在它自己对知识和理解的探求过程中不断向这些目标前进。

什么是证据？这绝不仅仅是个哲学问题，因为它闯进了生活的各个领域。你每时每刻都面临着极其大量的关于如何解释证据的选择。不论你走进哪个书店或碰见个什么书摊，差不多总能看到一些关于特异功能、UFO、百慕大三角、气功治病、偏方验方、占卜、相面、黑洞、生物回授、血型测人格、心理学的新理论等方面的书。在科学界，则有关于灾变论、基本粒子理论、黑洞、数学中的真理和存在、自由意志、人工智能、简化论与整体论等等的激烈论战。在生活中实用意义更强的方面，争论的题目有维生素C的功效、石油的真实贮量（不论地下的还是库存的）、通货膨胀或失业的原因------如此等等。还有佛教禅宗、芝诺悖论、精神分析学，等等等等。从在店里把书摆在哪个架子上这种微不足道的小事，到在学校里该教给孩子们什么样的思想这种至关重要的大事，证据的解释方式都起着难以估量的作用。

认识自己

在和解释证据有关的所有问题之中，最严峻的问题之一，就是设法通过对来自外界的混乱信号的解释来说明一个人到底是谁。在这种情况下，层次内部和层次之间都蕴藏着大量的冲突。每个人的心智都不得不同时应付个人对自我评价的内在需要和外界不断流入的影响自我印象的证据。其结果是信息在人格的不同层次之间乱转，而在它转来转去的时候，其中一部分放大了，一部分缩小了，一部分否定了，一部分变形了，然后这些部分又依次被卷入到一种旋转之中，如此反复进行下去------所有这一切，都是为了要在一个人是谁和他希望自己是谁之间进行调和。（见图81）

最终结局是，关于``我是谁''的完整画面是在整个精神结构中通过某种极其复杂的方式被拼出来的。而对我们每个人来说，这幅画面中都包含大量尚未解决、可能是无法解决的矛盾。这无疑提供了大量的动态张力，而这种张力对人来说起着很大作用。从这种张力之中，在关于我是谁的内部观念和外部观念之间，产生了指向各种不同目标的心理驱力，这就使我们每个人都成为独一无二的。这样一来，具有讽刺意义的是，某些为我们大家所共同具有的东西------作为具有自我反思意识的生物这一事实------反而导致了我们以各式各样的方式对关于各种事物的证据进行内在化，而这最终又成为创造不同的个性的主要力量之一。

哥德尔定理和其它学科

在人和充分复杂的形式系统间进行类比，这是很自然的，而后者和人一样，也有某种意义下的``自我印象''。\heititext{哥德尔定理}说明，对具有自我印象的一致的形式系统来说，存在着某些根本性的限制。但这一定理能否具有更普遍的意义？比如说，是否有``心理学中的\heititext{哥德尔定理}''？

如果我们把\heititext{哥德尔定理}当做一种比喻，当做灵感的源泉，而不是企图把它逐字逐句地翻译成心理学或任何其它学科的语言，那么或许它可能会为在心理学或其它领域中发现新的真理提供启发。但把它直接翻译成其它学科中的判断，并认为后者有同样的有效性，这是不合理的。认为数理逻辑中那些极其精致的结果能不加修改地适用于一个完全不同的领域，那是个极大的错误。

内省与精神失常：一个哥德尔式的问题

我认为把\heititext{哥德尔定理}翻译到其它领域之中还是有启发价值的，条件是要预先说明这种翻译是比喻性的，不能死抠字眼。在这个前提下，我发现了两种在哥德尔定理和人类思维之间建立类比的主要方式。其中之一涉及到考虑自己是否精神正常的问题。你怎样才能弄清你是否精神正常呢？这的确是个怪圈。一旦你开始探究自己精神的正常性，你就可能陷入一个极其讨厌的``信之则有''的漩涡之中，尽管这种情况绝非不可避免。每个人都知道，精神失常的人会用他们自己古怪的内部一致的逻辑去解释世界，但如果你只能用你自己的逻辑去检查它本身，那你怎样才能弄清你的逻辑是否``古怪''呢？我看这得不到任何答案。我只是由此想到了\heititext{哥德尔第二定理}，它说明那种断定了自身一致性的形式数论系统是不一致的\ldots\ldots{}

我们能否理解我们自己的心智或大脑

我所发现的另一个关于\heititext{哥德尔定理}的比喻性类比是一种有争论的见解，即认为我们最终也无法理解我们自己的心智和大脑。这种易于挑起争端的观点内容复杂，并涉及到许多层次，因此在提出它来时必须十分谨慎。``理解我们自己的心智和大脑''到底是什么意思？这可能是指要对其工作方式有个一般性的认识，就像力学对汽车的工作方式有一个一般性的认识一样；这可能是指要完全地解释为什么人们会做他们所做的一切事情；这可能是指要完全地理解人脑物理结构的所有层次；这可能是指要在一本书（或一座图书馆、一台计算机）中建立大脑的一幅完整的线路图；这可能是指要在每个时刻都精确地了解在自己大脑的神经层次所发生的事件------每次发射，每次突触传递，等等；这可能是指要写出一个能通过图灵测验的程序；这可能是指要如此完备地了解自己，以致于连潜意识和知觉这种概念都不再有意义，因为所有东西都是明摆着的；还可能是指任何别的什么。

上述自我反映中的哪一种------如果有的话------和\heititext{哥德尔定理}的自我反映最为相似呢？我有点说不准。其中的某些想法很糊涂。想要做到对自己大脑状态的一切细节进行控制，这简直是白日做梦。以此为前提是十分荒唐和无聊的。如\heititext{果哥德尔定理}说明上述想法是无法实现的，这也很难算是一种新发现。在另一方面，以某种深刻的方式认识自我，这个历史悠久的目标------让我们称其为``理解你自己的心智结构''------似乎还是具有一定的可能性的。但是否存在某种模糊的哥德尔式的循环，它限制着任何个人认识自己心灵时所能达到的深度？正像我们不能用自己的眼睛直接看到自己的面容一样，设想我们无法把我们完整的心智结构反映在实现它们的符号之中，这不是合情合理的吗？

所有元数学和计算理论中的限制性\heititext{定理都}说明，一旦你表示自身结构的能力达到某个临界点，那就该乐极生悲了：这就保证了你永远不能完整地表示自己。\heititext{哥德尔不完全性定理}、\heititext{丘奇不可判定性定理}、\heititext{图灵停机定理}、\heititext{塔斯基真理性定理}------都具有某些古老的神话故事的味道，它们警告你说：``寻找关于自我的知识就像踏上了这样一条旅途\ldots\ldots 它并非无处不通，它不能被绘制在任何图纸上，它永无尽头，它不可描述。''

但这些限制性定理对人有意义吗？这里是对此提出异议的一种方式：我要么是一致的，要么就是不一致的（后一种情况可能性更大，但为了完备性的缘故，我两种可能性都考虑）。如果我是一致的，那又有两种情况。（1）``低保真''情况：我的自我理解低于某个特定的临界点。这已经隐含设了我是不完全的。（2）``高保真''情况：我的自我理解已达到临界点，这样就能用比喻的方式用上述限制性\heititext{定理}来做类比，因此我的自我理解以一种哥德尔式的方法挖了它自己的墙角，因此我也就成了不完全的。情况（1）和（2）的前提是我具有百分之百的一致性------一种出现可能性极小的事态。更有可能的情况是我不具有一致性------但那就更坏了，因为那样在我内部就会有矛盾，而我怎么可能理解它呢？

不管一致还是不一致，任何人都无法摆脱自我的神秘性。或许我们都是不一致的。这个世界实在太复杂了，每个人都无法做到使他们的所有信念彼此协调。对一个需要快速做出许多决定的世界，张力和混乱是必不可少的。米格尔·德·乌纳穆诺曾说过，``如果一个人从未有过自相矛盾之处，那除非是他什么也不说。''我认为我们和那个禅宗大师处于相同的境地，他在一句话中出现多次自相矛盾之后，对摸不着头脑的道悟说：``吾亦不自了''。

哥德尔定理和个人的不存在

或许在我们的生活中最大的、也是最难把握的矛盾，就是下述知识：``曾经有一个时期我还不曾存在，而且还会有一个时期我将不复存在。''在某个层次上，一旦你``灵魂出窍''，把自己看成是``只不过是另一个人''，那么上述说法是完全有意义的。但在另一个层次------或许是个更深的层次------之上，个人的不存在是毫无意义的。我们所知道的一切都是嵌入在我们的心智之中的，而所有这些从宇宙中消失则是不可思议的。这在生活中是个无法回避的基本问题，也许正是\heititext{哥德尔定理}的一个最好的隐喻类比。每当你企图设想你自己已不存在时，你不得不设法跳到你自己之外，即把你自己映射成另外的什么人。你想自欺欺人地相信你能够在你自己身上引进一个旁观者对你的看法，就像TNT``相信''它能把它的元理论反映在它自身之内一样。但TNT只能在一定程度上包含它自己的元理论------并非全部。而对你来说，尽管你可以幻想已经跳出你自己之外，实际上你永远无法做到------就像艾舍尔的龙无法从它所在的二维平面中跳入三维空间一样。在任何情况下，这一矛盾都是巨大的，所以我们在生活中的大多数情况下往往对它采取视而不见的态度，因为设法解决它将完全是徒劳无功的。

在另一方面，禅宗的思想揭示了这种不可调和性。他们一次又一次地面对两种信念的冲突，一种是东方式的：``我即宇宙，因此我不存在这种想法是自相矛盾的''；另一种是西方式的：``我只是宇宙的一部分，我将死去，但宇宙没有我仍会存在。''

科学与二元论

科学常常被批评为过分``西方化''或``二元化''------即渗透着主体和客体、观察者和观察对象之间的两分法。虽然科学直到二十世纪仍是仅关心那些可以完全与观察者相分离的事物------如氧和碳、光和热、恒星和行星、加速度和轨道，等等------但这个时期的科学对更现代的科学时代来说只是一个必要的序幕。在新的时代中，生命本身成了研究的对象。``西方''科学已经逐步确定不移地走向对人类精神的研究------也就是说对观察者的研究。人工智能的探索是迄今为止沿这条路线所迈出的最远一步。在人工智能出现之前，科学研究中把主体和客体相混合而产生出奇妙结果这种事有过两次主要的预演。其一就是量子力学的革命，它的认识论问题涉及到观察者与观察对象的相互干扰。另一次就是元数学中主、客体的混合，从\heititext{哥德尔定理}开始，贯穿了我们讨论过的所有其它限制性定理。或许继人工智能之后的下一个步骤将是科学的自我应用：科学把自身当作研究对象。这是主、客观混合的一种不同方式------也许比人研究自己的精神更加纠缠不清。

顺便提一下，一个有趣的现象是，所有那些本质上依赖于主体和客体的聚变的结果都是限制性的结果。除这些限制性定理外，还有海森堡的测不准原理，其内容是说对一个量的测量会使得对一个相关量的同时测量成为不可能。我不知道为什么所有这些结果都是限制性的。随你怎么想吧。

现代音乐与绘画中的符号与对象之别

与主体-客体两分法密切相关的是符号-对象两分法，这已经由路德维希·维特根斯坦在二十世纪初深入地探索过了。后来``使用''和``谈论''这两个词被用来完成同样的区分。蒯恩等人也曾对记号和它们所代表的东西之间的联系进行了详尽的讨论。但对这个深奥抽象的问题贡献了许多思想的还不仅仅是哲学家。在二十世纪音乐和绘画所经历的危机中也反映出了与这个问题的深刻联系。例如，传统上音乐与绘画是把思想和情感通过一组``符号''来表现（即视觉形象、和声、节奏或别的东西），现在的一种趋势是去探索使音乐和绘画具有不表示任何东西的能力------它们只能是存在而已。这就是说它们就是纯粹的色块或声音，而在任何情况下都不再有任何作为符号的价值。

在音乐界，由于约翰·卡奇把一种禅宗式的音响方案带了进来，他产生的影响尤其大。他的许多作品都表达一种对``使用''声音------即用声音表达情感状态------的蔑视，但他喜欢``谈论''声音------即构造一些任意的声音序列，不考虑任何先前编好的编码，以使听众通过译码把它们变成消息。一个典型例子就是在第六章中描述过的那个用多台收音机构成的《大地幻景第4号》。我不准备对卡奇评头论足，但在我看来他的许多作品似乎是用来把无意义引入音乐，而且在某种程度上，是使无意义具有意义。偶然音乐是这一方向上的一种典型的探索。有许多别的当代作曲家也走上了卡奇的道路，但似乎没人具有同样多的创造性。安娜·洛克伍德有个作品叫《钢琴在燃烧》，是这样干的------把弦绷到最紧，把它们砸得尽可能地响。在莱蒙特·扬的一部作品中，为发出所需要的杂音，要推着钢琴满台跑，还要越过障碍物，就像推着一门土炮一样。

绘画在二十世纪中也经历了许多次这种巨变。开始是对于表达的屏弃，这的确是带有革命性的：这是抽象绘画的开端。在皮厄特·蒙德里安的作品中，揭示了从纯粹的表达到最抽象的模式的逐渐发展。在这个世界习惯了非表达性的艺术之后，超现实主义出现了。这是一种古怪的大转弯，有些像音乐中的新古典主义，即把最具表达性的绘画``败坏''了，并根据全新的理由来使用它：用它来产生振动、混乱和惊愕。这个流派的奠基人是安德烈·布洛东，开始于巴黎，其中较有影响的成员有达利、马格里特、德·蔡里科、唐居伊。

马格里特的语义错觉画

在所有这些画家中，马格里特最明确地意识到了这种符号-对象关系的奥秘（我认为这是使用-谈论区分的一种深入扩展）。他以此在观众中唤起了强烈的反应，尽管观众并没有把这种区分明确地用语言这样表示出来。例如，考虑一下他在静物画这个主题上进行的奇怪变奏（图137），这幅画的标题是《常识》。在画中有一个装满水果的托盘，这正是通常表现在静物画中的那种东西，它被放在一块空白画布上。符号与事实之间的冲突是尖锐的。但这还不是嘲弄的全部内容，因为整个这个东西本身当然只不过是一幅画------事实上，是一幅画着非标准对象的静物画。

\pandocbounded{\includegraphics[keepaspectratio]{./media/OEBPS/Images/img_137.jpg}}

图137．常识

马格里特作（1945-46）

马格里特的烟斗画系列令人既着迷又困惑。请看《两个谜》（图138）。当把注意力集中在里面那张画上时，你得到的消息是``符号和烟斗是不同的''。然后你的目光转向那个浮在空中的``真''烟斗------你感到它是真的。而另一个只是个符号。但这显然全部都错了：二者都是在你眼前的同一个平面上的。如果以为一个烟斗是在一幅两重嵌套的画中，因此在某种意义上比另一个烟斗``更不真实''，那将是彻头彻尾的谬误。一旦你想``进入这间屋子''，你就已经上当了：你错把图像当成了现实。为了使你的这种轻信保持前后一致，你应当愉快地进入下一层，把图像中的图像和事实也混在一起。不受这种骗的唯一办法是把这两个烟斗都看成只不过是离你鼻子几寸远的一个平面上的一些色斑。那时，也只有那时，你才能懂得那一行字的全部含义：Ceci n\textquotesingle est pas une pipe{[}法语：``这不是烟斗''{]}------但具有讽刺意味的是，在所有的东西变成色斑那一瞬间，这句话也变成了色斑，因此就失去了它的意义！换句话说，在那个时刻，这幅画的语言消息以一种非常哥德尔化的方式自我销毁了。

\pandocbounded{\includegraphics[keepaspectratio]{./media/OEBPS/Images/img_138.jpg}}

图138．两个谜

马格里特作（1966）

《咏叹调与歌曲》（图82）是取自于马格里特烟斗画系列中的一幅，它完成了《两个谜》所做的一切，不过是在一个层次上，而非两个。我所画的《以烟为号》和《如烟似梦》（图139，140）构成了``马格里特主题变奏曲''。请盯着《以烟为号》看一阵，你不久就会发现一条隐藏在其中的消息：``Ceci n\textquotesingle est pas un message''{[}法文：``这不是消息''{]}。这样，如果你发现了这条消息，它就否定了自己------但如果你没发现它，你就根本没抓住要害。由于我画的这两幅烟斗都具有间接的``自我熄灭性''，它们可以不严格地对应于哥德尔串G。

\pandocbounded{\includegraphics[keepaspectratio]{./media/OEBPS/Images/img_139.png}}

图139．以烟为号

作者绘

绘画中把``使用''与``谈论''相混淆的一个经典范例就是在画中出现一块调色板。虽然这块调色板是画家的表示技巧所创造出来的错觉，但这块画出来的调色板上的颜料又的确是从画家的调色板上来的颜料。颜料代表它自己------它不是代表任何其他东西的符号。在《唐·乔万尼》中，莫扎特使用了一个与此相关的技巧：他在乐谱中明确地写出了管弦乐队调弦的声音。与此相类似，如果我想让``我''这个字代表它自己（而不是作为我的符号），我就把``我''直接写进我的文章中，然后我给``我''加上引号。这就得到一个``\,`我'\,''（而非``我''或`\,``\,`我'\,''\,'）。明白了吗？

\pandocbounded{\includegraphics[keepaspectratio]{./media/OEBPS/Images/img_140.png}}

图140．如烟似梦

作者绘

现代绘画的``编码''

有许多颇有影响的人物参与了对绘画中的符号一对象二元论的进一步探索，没人能无一遗漏地把这些人的特点都列举出来。毫无疑问，对禅宗抱有浓厚兴趣的约翰·卡奇对绘画也像对音乐一样产生了重大影响。他的朋友贾斯泊·约翰斯和罗伯特·罗斯琴伯格为了探索符号和对象之间的区别，直接用对象作自己的符号------或反过来说，用符号自身作为对象。所有这些常识都可能导致下述观念的瓦解：绘画与现实之间有一段距离------或者说绘画通过``编码''来说话，而观众必须像译员那样来接收。这些现代画家的想法是取消翻译这个步骤，让对象本身呈现出来，仅此而已。如果这就是他们的目的，那他们可碰了个大钉子，而且这可能是无法避免的。

一旦一个对象在画廊中展出或被称为一件``作品''，它就会带有一种意味深长的气氛------不论你怎样告诫观众不要寻求它的意义。实际上事与愿违，你越告诉观众不要把这些对象神秘化，他们越觉得神秘。不管怎样说，如果一个放在博物馆地板上的破木箱只不过是一个放在博物馆地板上的破木箱的话，那看门的为什么不把它拖出去扔进垃圾堆呢？为什么还要标上一个艺术家的姓名？为什么这个艺术家想让艺术失去神秘性？为什么不把外面的那块烂泥也标上一个艺术家的名字？这是在捉弄人吗？是我疯了，还是这些艺术家疯了？越来越多的问题会涌入观众的脑海，他无法不这样想。这就是艺术品本身自动造成的``框架效应''。谁也没办法抑制好奇的心智产生疑虑。

当然，如果其目的是灌输一种禅宗式的观念，即认为世界既无规律也无意义，那这种艺术品大概仅仅能作为------就像用理智把握禅宗思想的效果那样------一种催化剂，其作用是激发观众灵感，使他们跳出作品之外，了解这种主张抛弃``内在意义''、把世界作为一个整体来把握的哲学。在这种情况下，从短时期看，这种艺术品是挖了自己的墙脚，因为观众的确是在考虑它的意义，但它在一些人身上达到了长期目标，即把他们引向了它的来源。但不论在哪一种情况下，说不存在把思想传送给观众的编码，这都不是真的。实际上，这是一种更复杂的编码，其中表述了诸如``没有编码''等等的东西------也就是说它部分是编码、部分是元编码，如此等等。在这些非常禅宗式的艺术对象中传送了一个消息的缠结的层次结构，或许这就是许多人以为现代艺术晦涩难懂的原因。

再谈主义

卡奇领导了一个旨在打破艺术与自然之间的界线的运动。在音乐上，主题就是认为所有音响都是平等的------一种声音民主主义。这样看来``无声''恰如``有声''同样重要。列奥纳德·毕·迈尔在他的著作《音乐、美术与思想》中，把这场运动称为音乐中的``超验主义''，他写到：

如果区分艺术和自然是种错误的举动，那么审美评价就会是毫无意义的。这样，一首钢琴奏鸣曲和一块石头、一场雷雨、一个海星相比，并没有更高的价值。卢西诺·伯里奥{[}一位现代作曲家{]}曾写到：``使用范畴的判断，例如对与错、美与丑这些音乐美学中的典型理性思维方式，已经不能再被用来理解一个当代作曲家在创作可听形式和音乐活动时的动机和过程。''

接下来，迈尔进一步描述了超验主义的哲学立场：

\ldots\ldots 所有事物在所有时空条件下都不可避免地相互联系在一起。在宇宙中发现的任何区别、分类或组织都是随意的。世界是一个复杂、连续、单一的事件。\footnote{}［这里能看见芝诺的影子！］

我觉得用``超验主义''作为这场运动的名字未免过于冗长。我用``主义''来代替它。由于``主义''一般是某个词的后一部分，用它作名字正好暗示了一种没有内容的意识形态------对这种意识形态不论你做怎样的揭示，都是一样的。主义是体现在艺术中的禅宗精神。正如禅宗的中心问题是要揭出自我的本来面目一样，二十世纪中艺术的中心问题似乎是要指出艺术到底是什么。上面提到的那些东冲西撞的行为正是其同一性危机的部分表现。我们已经看到，使用-谈论两分法一旦推向深入就会转化为符号-对象二元论这一哲学问题，而这又与思维之谜相联系。马格里特为他的画《人类的处境Ⅰ》（图141）写过下面的话：

我在一间屋子的窗前放了一幅画，从屋内向外看去，画的内容恰好表现了这幅画所遮挡住的那一部分风景。因此，画上那棵树表示着被画挡住的那棵位于屋外的树。而在观众的脑海里，它的存在既表示屋里画中的那棵树，同时又表示外面真实景色中的那棵树。我们正是这样看世界的：我们把它看作存在于我们之外，尽管它只不过是我们在内心中所体验到的一个对于世界的心智表示而已。\footnote{}

\pandocbounded{\includegraphics[keepaspectratio]{./media/OEBPS/Images/img_141.jpg}}

图141．人类的处境Ⅰ

马格里特作（1933）

理解心智

首先通过画中含蓄的图像，然后又通过直接的语言，马格里特表述了下列两个问题间的联系：一是``符号是怎样工作的？''，二是``我们的心智是怎样工作的？''因而他又把我们带回到那个前面提出的问题：``我们是否有希望理解我们的心与脑？''

我们是否会遇到某种阻止我们理解我们的心智的魔怪似的哥德尔式命题？如果你并未接受对``理解''的一种完全非理性化的定义，那我看在最终理解我们心智的旅途上就不会有哥德尔式的障碍。例如，期望着能理解大脑的一般工作原理，就像我们理解汽车发动机的一般工作原理那样，我认为是完全合理的。这完全不同于试图理解任何特定大脑的每个具体细节------更不必说试图对自己的大脑做到这一点了！我看不出\heititext{哥德尔定理}怎样才能对这种设想的可行性产生影响，不管对\heititext{哥德尔定理}怎样解释都是这样。我看不出\heititext{哥德尔定理}能用何种理由来限制我们刻划与核实那种发生在神经细胞介质上的思维过程的一般机制的能力。我看不出\heititext{哥德尔定理}怎么会设置障碍，阻止我们在计算机（或它们的后代）之上实现各种能获得和大脑大致相同结果的符号处理功能。那种试图把某个特定的人的心智复制在一个程序中的想法完全是另外一个问题------但构成一个智能程序却是一个有限得多的目标。\heititext{哥德尔定理}并没有禁止我们通过程序再现我们自己的智力水平，正像它并没有禁止我们通过DNA中遗传信息的传递，再继之以教育，以再现我们自己的智力水平一样。实际上我们在第十六章中已经看到，一个非常哥德尔化的机制------蛋白质和DNA构成的怪圈------恰好是使智能得传递的条件。

这样说来，\heititext{哥德尔定理}是否与我们关于自己心智的思考毫无关系呢？我想还是有关系的。尽管作用不像某些人所想象的那样神秘和具有限制性。我认为对哥德尔的证明进行理解的过程，理解那些涉及了任意编码的构造、复杂的同构、高低不同的解释层次、以及自我反映的能力，会在我们关于符号和符号处理的表象中注入某种丰富的底蕴和风味，这将深化我们对于不同层次上的心智结构之间关系的直觉。

智能是碰巧不可说明的吗？

在提出哥德尔证明的一种在哲学意义上迷惑人的``应用''之前，我想先介绍一下智能的``碰巧不可说明''性。这种观点含义如下：我们的大脑可能与汽车发动机不同，是一个我们以任何方式也无法分解的难以对付的系统。在目前这种时候，我们不知道我们的大脑是否能通过反复尝试而被劈成整齐的层次，其中每一层都能用更低的层次加以解释------或者我们的大脑是否可能挫败我们的所有进行分解的尝试。

但即使我们在自我理解时遭到失败，在其背后也并非一定有某种哥德尔式的``缠绕''。这可能仅仅是命运的一种偶然结果，是由于我们的大脑碰巧未能强到能理解它自身。不妨以鸡为例，它们的大脑虽然是远远低于自我理解所要求的水平，但仍和我们自己的脑很相像。事实上，螃蟹、树獭、食蚁兽的大脑------甚至乌龟或某种远比我们聪明的未知生物的大脑------可能都基本上在用同一组原理进行操作。螃蟹所具有的智能可能远远低于理解那些原理是怎样相互结合以产生出心智的特性所需要达到的某个阈值，而人可能离那个阈值更近一些------或许正好在它下面，甚至可能在它上面。关键之处是，并没有根本性的（即哥德尔式的）理由说明那些特性是不可理解的，它们对更高级的智能生物来说可能是一目了然的。

不可判定性与高层观点不可分离

除去这种关于脑的碰巧不可说明性的悲观看法之外，哥德尔的证明能为我们解释我们的心与脑提供哪些见解呢？哥德尔的证明提供了这样的观点：一个系统的高层观点可能会包含某种在低层上完全不具有的解释能力。我这样说是这个意思：假设某人给你哥德尔的不可判定串G，形式上是TNT中的一个串。再假设你对哥德尔配数法一无所知。你所要解答的问题是：``为什么这个串不是TNT的一个定理？''。现在你对这种问题已经很熟悉了，例如，如果对于S0=0问你同一个问题，你可能有一种现成的解释：``它的否定\textasciitilde S0=0是一个定理''。这样，再加上你知道TNT是一致的，就解释了为什么所给出的串是个非定理。这就是我所谓的``TNT层次上的''解释。请注意，这完全不同于对WU为什么不是WJU系统的一个定理的解释：前者是在J态，后者是在W态。

现在G怎么样呢？能处理S0=0的TNT层上的解释对G无能为力，因为\textasciitilde G不是一个定理。那些对于TNT没有整体观点的人搞不懂他为什么不能依照规则生成G，因为作为一个算术命题，它看上去没有任何毛病。事实上，一旦G变成了一个被全称量词约束的串，对G中的变量用数字去替换所得到的每个例都可以被推导出来。解释G的非定理性的唯一途径是发现哥德尔配数概念，在一个完全不同的层次上来看TNT。想在TNT层次上写出解释决不单是困难和复杂的问题，而是根本不可能做到。这种解释根本就不存在。在TNT层上从根本上看缺乏一种为更高层次所具有的解释能力。可以说，G的非定理性本质上是一个高层事实。我猜测所有不可判定命题都是这种情况，这也就是说：每个不可判定命题实际上都是一个哥德尔语句，它通过某种编码在某一系统中陈述了它自身的非定理性。

意识是一种高层所具有的现象

这样看来，哥德尔的证明就提示了------虽然绝非证明了！------可能存在某种观察心与脑的高层方式，涉及到在低层不出现的概念，而且在这个层次上可能会有在低层次上不存在------甚至从原则上讲也如此------的解释能力。这将意味着某些事实在高层可以很容易地得到解释，但在低层则根本不行。无论一个低层陈述被搞得多长、多复杂，它也无法对问题中涉及到的现象加以解释。这可以类比于下面这个事实：如果你在TNT中反复推导，不论你搞得多长、多复杂，你也绝不可能推出G来------尽管事实上在一个更高的层次上你能看到G是真的。

这种高层概念是什么东西呢？许多具有整体论或``唯灵论''倾向的科学家和人本主义者早就提出，``意识''就是一种无法用大脑成分来加以解释的现象，所以至少它可以算一个候选者。另一个更令人迷惑不解的概念是``自由意志''。因此这些性质大概是``浮现''出来的，也就是说，对它们的解释不能仅借助于生理学来完成。但关键是要明白，如果我们在提出这种大胆的假说时想要得到哥德尔证明的引导，我们就必须把这一类比贯彻到底。尤其至关紧要的是，要记住G的非定理性是有一个解释的------它不是一个彻头彻尾的谜！这个解释依赖于同时在不止一个层次上进行理解，同时一个层次上的理解方式要反映它的元层次，还依赖于这种反映的结果。如果我们的类比是成立的，那么``浮现''的现象就会变得可以根据心智系统的不同层次的关系来进行说明了。

意识的核心是怪圈

我确信，对我们大脑中那些``浮现''出来的现象------例如想法、希望、表象、类比、以至于意识和自由意志------的解释都基于一种怪圈，一种层次相互作用，其中顶层下到底层并对之产生影响，而与此同时它自身又被底层所确定。换句话说，是不同层次之间的一种自我强化的``共鸣''------正像汉肯句子那样，只是通过断定自己的可证明性，使它真的变成可证明的了。自我成为一种存在的时刻也恰是它具有反映自身能力的时刻。

上面的观点不应当被看作是一种反简化论立场。它只是意味着，一种对心智的简化论解释，如果要成为可理解的，就必须引进像层次、映射、意义等等``软''概念。原则上，我不怀疑存在着关于大脑的彻底简化论的、但却是极其费解的解释，问题在于怎样才能将其翻译成为我们自己所了解的语言。当然我们不要一个根据粒子的位置和动量所作出的描述，我们想要一个把神经活动与``信号''（中层现象）联系起来的描述------然后再把信号联系于``符号''和``子系统''，包括那个预先假设其存在的``代表自我的符号''。这种从低层物理硬件到高层心理软件的翻译类似于从数论判断到元数学判断的翻译。你还会记得，正是在这个翻译中出现的层次构成了哥德尔的不完全性和汉肯句子自我证明的特征。我猜测正是一种类似的层次交叉造成了我们几乎不可分析的自我意识。

为了应付内容丰富的脑-心系统，我们必须能在层次之间随心所欲地滑移。更进一步来说，我们还必须承认各种类型的``因果关系''，即一个描述层次上的事件能以各种方式``导致''其它层次上的事件的发生。有时候说事件A``导致''了事件B，只不过是因为其中一个是另一个在别的描述层次上的翻译而已。有时候``因果关系''具有它通常所具有的意义：物理因果关系。这两种因果关系------可能还有其它种------都必须在对心智的任何解释中被承认，因为我们在精神的缠结层次结构中必须同时承认因果链的向上和向下传播，就像在中心法则映射中那样。

在我们自我理解过程的核心，将会是对我们心智中的层次结构的理解。我的立场很接近于神经科学家罗杰·斯珀里在他的《心智、大脑与人道主义的价值观》这篇精彩的文章中所表达的观点。下面是引自那篇文章中的一段话：

在我自己所设想的大脑模型中，意识被表示成一种非常真实的动因。我认为它在脑中事件的因果序列和控制链中发挥着重要作用。它在这些过程中表现为一种活跃的、起作用的力量\ldots\ldots 简而言之，它导向这样一个问题：在充斥着脑海的各种动力之间到底是谁在推着谁转。换言之，这就是要列出脑中控制成分之间的等级关系。在人脑中存在着一个由形形色色的动力所组成的世界，而且动力里面又包含着动力。在我们所认识的宇宙里，没有比这几十厘米见方的空间中所包含的东西更为复杂的现象了。\ldots\ldots 长话短说，如果我们不断沿脑中命令链向上攀登，我们会在其顶端发现那些全面组织起来的动力和大规模的脑兴奋模式，它们和心理状态或精神活动相互关联\ldots\ldots 在大脑里的这个命令系统中离顶点不远的地方\ldots\ldots 我们发现了思想。人之所以优越于黑猩猩，就是因为人有思想和观念。在这里提出的大脑模型中，一个想法或观念的潜在动力，变得像一个分子、一个细胞、或一个神经脉冲所具有的动力那样真实。思想引出了思想，并且为新思想的形成提供了帮助。它们彼此相互作用，和同一个大脑或邻近的大脑中的心理力量相互作用，而且借助于全球通讯系统，还能和千里之外的其它国家中的大脑里的心理力量相互作用。它们还和外部环境相互作用，在进化过程中产生了突破性进展，这一进展远远超过了迄今为止在进化过程中发生的其它进展，包括有生命的细胞的突然出现在内。\footnote{}

在主观语言和客观语言这两种论述语言之间有一个著名的分裂。例如，``主观的''红色感受和``客观的''红光波长。对很多人来说，这二者似乎永远是不可调和的。我不这样认为。正像关于艾舍尔的《画手》的两种观念并非不可调和一样------一种是``在系统内看''，此时两只手在互相画；另一种是从外面看，此时它们都是艾舍尔画的。对红色的主观感受来自大脑的自我感知中心，而客观的波长则属于你退出系统之外时的观察事物方式。尽管我们之中没人能退得足够远，以至于可以把一切都看成一副``大画''，但我们不应忘记这幅大画是存在的。我们应当记住，物理定律是所发生的一切的原因------它们藏在神经网络的犄角旮旯的深处，是我们高层次的内省式探究所无法企及的。

自我符号与自由意志

在第十二章中我曾提出，我们所谓的自由意志实际上是代表自我的符号（或子系统）与脑中的其它符号相互作用的结果。如果我们把符号看成上面附有意义的高层实体，那我们就可以试着解释一下符号、自我符号和自由意志的相互关系了。

为了要认识自由意志问题，一种办法是先用一个我认为是等价的问题来替换它，而这个问题相对来说涉及的概念不那么复杂，也不太引起异议。我们不再问``系统X是否有自由意志？''，而是问``系统X是否在进行选择？''。通过仔细探讨当我们把一个系统------机械的或生物的------描述成``能够进行选择''时，到底是什么意思，我想我们能对自由意志有进一步的理解。下面我们来看几个不同的系统，在各种各样的环境中，我们会觉得应当把它们都描绘成是``在进行选择''。考察这样的系统会对我们有所帮助，而且从这些例子中，我们将对``在进行选择''到底是什么意思有所了解。

让我们以下述系统作为范例：一个小球滚下一个高低不平的山坡；用一个袖珍计算器依次求2的平方根的小数展开式中的各位数字；一个能下一手好棋的复杂程序；一个机器人在走T型迷津（这种迷津中只有一处岔路，在岔路的一端放有奖品）；还有一个正在对付复杂矛盾的人。

首先，小球滚下山坡时是一种什么情况呢？它是在选择吗？我想我们大家都会同意说它没有，尽管我们当中没人能预测它所走的路径，甚至连预测其中一小段都做不到。我们觉得它在实际所通过的路径之外不可能走别的路，而它只不过是被无情的自然规律所迫才这样做的。当然，根据我们组块化了的``心智物理学''，我们能够为这个小球想象出许多各不相同的``可能的''路径，而我们看到它在现实世界中只通过其中的一条。因此，在我们心智的某个层次上，我们可能不禁会觉得这个小球是在那一大堆想象的路径里``选择''了唯一的一条。但是，在我们心智中的另一个层次上，我们会本能地理解到这种心智物理学只是在我们为世界建立内心模型时起辅助作用，而使真实的物理事件序列得以发生的机制并不要求自然界在某个假想的宇宙（``上帝之脑''）中先经过一个类似的过程来制造出各种可能性，然后再在它们之中进行选择。因此我们不会把这一过程称之为``选择''------尽管我们承认在类似情况下使用这个词常常是具有实用意义的，因为它具有一些启发性。

对那个用来求2的平方根的可编程序计算器和那个弈棋程序，我们又能作何评论呢？在这里我们可以说碰到了一个``想象的''小球，滚下一个``想象的''山坡。事实上，在这种情况下反对``作出选择说''的证据要比小球那种情况更强。原因是如果你要想重复那个小球实验，你无疑会看到一条完全不同的通向山坡下面的路径；而如果你重新运行那个求2的平方根的程序，你每次都会得到同样的结果。那个小球似乎是每次都``选择''了一条不同的路径，不管你怎样精确地试图重现它初次往下滚时的条件，而那个程序每次都通过完全相同的路径。

对那个假想的弈棋程序的情况来说，存在着各种不同的可能性。如果你和某些特定的程序对弈，然后第二盘采用和第一盘完全相同的走法，这些程序也将准确地重复以前的走法，丝毫看不出学到了什么东西或具有什么变化的愿望。另外，有一些程序带有随机装置，能够产生一些变化，但并非出自于任何深切的期望。这种程序中的内部随机数发生器可以被重置成初始状态，然后就能导致同样的对局。还有一些程序的确能从它们的失误中吸取教训，而且能根据对局的结果改变它们的策略。这种程序就不会连续重复同一盘对局。当然，你也可以使时光倒流，抹去存储器中表示学习的所有变化，就像你能重置随机数发生器一样，但这样做似乎有些不太友好。此外，我们有理由怀疑，如果所有细节------当然包括你的脑------都被重置成它们初次发生时的状态，那么你就能改变你自己过去所做的决定吗？

不过还是让我们转回到原来的问题，即在这里``选择''这个概念是否适用。如果程序只不过是``滚下想象的山坡的想象的小球''，那它们是否在作选择呢？当然答案必然是主观的，不过我要说在这里所考虑的情况和小球的情况基本相同。但是，我必须补充说明，使用``选择''这个词的要求变得非常强烈，尽管你只是把这个词作为一种方便的或带启发性的简称而已。一个弈棋程序能够超前搜索许多可能的岔路，这完全不同于一个滚动的小球。这个事实使得它更像一个生物，而不像一个计算2的平方根的程序。但是，在这里仍没有深刻的自我认识------也不存在自由意志。

现在让我们接着来设想一个拥有一堆符号的机器人。这个机器人被放在一个T型迷津中。但是，给它编的程序不是让它去寻找奖品，而是让它在2的平方根中的下一位数字是偶数时就向左转弯，是奇数时就向右转弯。现在这个机器人能够用它的符号来为外部环境建立模型，这样它就能观察到它自己所作的选择。每当走到T型路口时，如果你问这个机器人说：``你知道你这次会往那边转弯吗？''它会回答说``不知道''。然后为了继续前进，它将激活它的``决策器''子程序，该子程序能算出2的平方根中的下一位数字，并据此作出决策。但机器人对决策器的内部机制一无所知------它用机器人的符号来表示时只不过是个黑箱，以某种神秘的、而且好像是随机的方式输出``左''和``右''。除非这个机器人的符号可以摸出这个决策器的``脉搏''，否则它也会对它自己所作的``选择''大惑不解的。那么这个机器人是否在作选择呢？你可以设身处地的去想想。如果你被封装在一个滚下山坡的小球里，完全无法影响它的滚动路径，但能够用你全部的人类智能去观察它，你是否会觉得这个小球通过的路径与某些选择有关呢？当然不会。除非你的思想在影响着结果，否则即使有符号存在也无济于事。

因此，现在我们要对我们的机器人做一个修改：我们允许它的符号------包括自我符号------被用来影响所作的决定。那样，在这个完全依照物理定律运行的程序的例子中，就似乎比前面那些例子更多地深入了选择的本质。当这个机器人自己的组块化``自我''概念出现后，我们就开始与这个机器相认同了，因为它的所作所为与我们相似。它不像计算2的平方根时那样，似乎不用任何符号来控制作出决定的过程。诚然，如果我们在一个相当低的层次上来看这个机器人的程序，那和平方根程序也差不多。它也是一步一步地执行，最后输出``左''或``右''。但在一个较高的层次上我们可以看到，事实上符号正在被用来摹写环境和影响决定。这就极大地影响了我们对这个程序的看法。在这个阶段，``意义''已经进入了视野------正像我们用我们自己的心智所处理的那种意义一样。

一个层次交错的哥德尔漩涡

这时，如果外边的某个动因建议这个机器人下一次选择``左''，这个建议就会被送到那一大堆彼此纠缠不休的符号中去。在那里，它将被无情地吞没在与自我符号的相互作用之中，就像一叶小舟被吸进漩涡一样。这就是系统中的漩涡，那里，所有层次都相互交错。这时，那个``左''碰到了一个符号的缠结层次结构，并在层次间被上下传送。自我符号并不能控制它的所有内部过程，因此当一个实际决定出现时------``左''、``右''或系统之外的什么东西------系统将无法说清它来自何处。这不像一个标准的弈棋程序，在那种情况下程序不控制它自己，因此也就不知道它的棋步是哪里来的。在这里，程序是在进行自我控制，也具有关于它的想法的想法------但它不能控制它自己行动的每个细节，因此会对它的工作过程有一种``直觉''，但没有完全的理解。自由意志的感觉就产生于这种``有自知之明''和``无自知之明''的平衡中。

例如，让我们考虑一个作家，他想把他心理表象中的一些想法表达出来。他拿不准这些表象在他心目中是如何相互联系的，因此就不断地尝试，在表达事物时开始用这种方式，后来又用那种方式，最后才选定某种说法。但他知道这一切都是哪里来的吗？只有一些模糊的感觉。其来源像一座冰山一样，大部分都深藏在水下，是看不见的------而他知道这一点。我们还可以考虑一个作曲程序------就像我们前面所讨论过的那种------并想想什么时候我们才能觉得应当称它为一个作曲家，而不仅是一个人类作曲家的工具。或许促使我们这样做的条件是在程序内部存在着用符号表示的关于自我的知识，而且该程序在有自知之明和无自知之明之间平衡得恰到好处。这与系统是否以确定的方式运行毫无关系，使得我们将它说成是``在做选择''的是：我们是否能认同于程序运行时所产生的过程的一个高层描述。在一个低级层次上（机器语言层），这个程序看起来很像任何其它程序，而在一个高级层次上（组块化层），会浮现出``意愿''、``直觉''、``创造性''和``意识''等性质。

关键的思想是这个漩涡本身就决定了心理过程的缠结性和``哥德尔式的''特性。有人曾对我说：``关于自指等等这些东西很有趣，但你真的认为这里面有严肃的问题吗？''我当然是这样认为的。我认为最终将会发现这些问题处于人工智能的核心之中，而且是所有理解人类心智工作方式的尝试的焦点。因此，我才把哥德尔这块稀世瑰璧深深地镶嵌在我这本书之中。

一个层次交错的艾舍尔漩涡

\pandocbounded{\includegraphics[keepaspectratio]{./media/OEBPS/Images/img_142.jpg}}

图142．画廊

艾舍尔作（石版画，1956）

在艾舍尔的《画廊》（图142）中，为我们提供了对一个缠结的层次结构中的``台风眼''的一个美妙得令人吃惊的、同时又荒诞得令人烦乱的描绘。我们在其中所看到的是一个画廊，里面站着一个青年，正在看一幅画。画里有一艘船停泊在一个城镇的港湾中，这个城镇可能是在马耳他。之所以这么说，是因为这里的建筑物中有一些小塔楼，偶尔可以见到小圆屋顶，还有平的石屋顶。其中的一个屋顶上面坐着个男孩，正悠闲地在晒太阳。比他低两层楼的地方有个妇女------也许是他母亲------正从她房间的窗户朝外看，而她的房间下面正是一个画廊，里面站着一个青年，正在看一幅画。画里有一艘船停泊在一个城镇的港湾中，这个城镇可能是在马耳他------怎么！？我们又回到了我们出发时所在的那个层次，尽管从逻辑上说我们决不可能这样做。让我们把所看到的东西画成一张图（图143）。

\pandocbounded{\includegraphics[keepaspectratio]{./media/OEBPS/Images/img_143.png}}

图143．《画廊》的抽象图示

这张图所表示的是三种``之中''。画廊是物理地处于城镇之中（``包含''），城镇是艺术地处于图画之中（``描绘''），而图画又是心理地处于人脑之中（``表现''）！现在虽然这张图似乎已经令人满意了，事实上它还是带任意性的，因为所示层次的数目有非常大的任意性。下面是表示该图上半部分的另一种方式（图144）。

\pandocbounded{\includegraphics[keepaspectratio]{./media/OEBPS/Images/img_144.png}}

图144．对图143的一种压缩形式

我们已经消去了``城镇''那个层次，从概念上说它是有用的，但缺了它也无关大局。图144看上去很像那张对应于《画手》的示意图：都是一个由两个步骤构成的怪圈。步骤的划分是任意的，尽管它们在我们的头脑看来可能很自然。

为了进一步强调这一点，我们可以为《画廊》画出一张再次``压缩''了的示意图，如图145所示：

\pandocbounded{\includegraphics[keepaspectratio]{./media/OEBPS/Images/img_145.png}}

图145．对图143的进一步压缩形式

这就以赤裸裸的形式把图画中的悖论展现出来了。那么------如果说这幅画是``在它自身之中的''，则那个青年不也是在他自身之中的吗？这个问题的答案在图146里。

这样，我们就看到这个青年``在他自身之中''，这是在``之中''的三种不同意思联合起来所形成的一种有趣意义下而言的。

这张图使我们想起了通过一步就构成自我相关的说谎者悖论，而那张包括两个步骤的图则类似于互相谈论的一对句子。我们不能使这个循环变得更紧缩了，但我们可以把它张得更大，即选择插入任意数目的中间层次，例如``画框画廊''和``建筑物''。如果我们这样做的话，我们将得到多步骤的怪圈，与《瀑布》（图5）及《上升与下降》（图6）同构。层次的数目要依我们看来怎样才算``自然''而定，这可能会取决于环境、目的或心智的框架。一个人在哪些地方能感知到层次，这是一个与直觉和审美偏好有关的问题。

\pandocbounded{\includegraphics[keepaspectratio]{./media/OEBPS/Images/img_146.png}}

图146．对图143进行压缩的另一种方式

那么对于我们这些``画廊''的观众来说，是否也会由于看了这幅画就被吞没在自身之中了呢？并非如此。我们因为身在该系统之外，所以能逃离那个漩涡。而且当我们看这幅画的时候，我们能看到那个青年无法看到的东西，例如在中心的那个``疵点''上的艾舍尔的签名``MCE''。虽然这个疵点看起来像一个缺陷，但这缺陷是在我们的期望之中的，因为事实上艾舍尔无法完成图画的那一部分，除非他违背他作画时所一贯遵循的规则。漩涡的中心是------而且必须是------不完全的。艾舍尔可以使这一部分任意地小，但他无法彻底摆脱这个问题。这样，我们在外面可以知道``画廊''从根本上说是不完全的------这个事实是画中那个青年所无法得知的。艾舍尔就是这样地为\heititext{哥德尔不完全性定理}提供了一个形象化的比喻。因此，我才把哥德尔和艾舍尔这两块稀世瑰璧深深地镶嵌在我这本书之中。

一个层次交错的巴赫漩涡

\pandocbounded{\includegraphics[keepaspectratio]{./media/OEBPS/Images/img_147.png}}

图147．``无穷升高的卡农''的六角形示意图

在使用谢泼德音调的条件下，它形成了一个不折不扣的闭合循环

当我们在看怪圈的图示时，不禁会想起《音乐的奉献》中的无穷升高的卡农。它的示意图包含夫个步骤，就像图147所示的那样。不幸的是当它返回C调时，它和开始时的音高相比高了一个八度。令人惊讶的是，使用所谓的``谢泼德音调''，就可以使它恰好返回到起始时的音高。

\pandocbounded{\includegraphics[keepaspectratio]{./media/OEBPS/Images/img_148.png}}

图148．谢泼德音阶的两个完整周期的钢琴谱

每个音符的音量与其大小成正比，因此，当最上面的声部消失时，下面一个新的声部出现了。{[}该图为唐纳德·伯德的程序``斯马特''所印出{]}

这种想法是心理学家罗杰·谢泼德发现的，因此也就以他的名字命名了。谢泼德音阶的原理如图148所示。如果用语言来描绘是这样的：在若干不同的八度音域中演奏平行的音阶，每个音符独立地加权，而且当音符上升时，权重就发生变化，让最上面一个八度逐渐减弱，而与此同时逐渐增强最下面的一个八度。当在一般情况下你已经升高了一个八度的时候，权重的变化使得你恰好重现了原来的音高\ldots\ldots，这样你就可以``永远向上走''，其实一点也没升高！你不妨找架钢琴试一试。如果音高是在计算机控制下精确地合成出来的，效果还会更好。那时这种错觉就会强得足以使人上当了。

这个奇妙的音乐发现使得``无穷升高的卡农''可以这样演奏：当它``升高''了一个八度之后，又回到了自身之中。这个主意是我和斯科特·凯姆一起想出来的，而且已经用一个计算机音乐系统实现在一盘磁带上了。其效果极其微妙------但非常寘实。有意思的是巴赫本人在某种意义下显然是知道这种音阶的，因为在他的曲子中不时能发现一些段落在大体上利用了谢泼德音调的普遍原理------例如在他为管风琴所写的《G小调幻想曲和赋格中幻想曲》的中部。

汉斯·西奥多·大卫在他的著作《约·塞·巴赫的〈音乐的奉献〉》中写到：

在整部《音乐的奉献》中，读者、听众和演奏者始终都在寻找各种形式的国王主题。因此这个作品可以说是一部名副其实、不折不扣的``探求曲''。\footnote{}

我认为这是真实的。没人能把《音乐的奉献》看透。当一个人自以为知道了其中的一切的时候，总还有他所不知道的东西。例如，在其中那个他不情愿即席演奏的六部无插入赋格的结尾处，他狡猾地把自己的名字藏在了上面两个声部之间。在《音乐的奉献》中，事情往往在许多层次上发展。有关于音符和字母的技巧，有国王主题的精巧变奏，有各种原始形式的卡农，有复杂得异乎寻常的赋格，有优美并极其深沉的情感，甚至还有由作品发展的多层次性所带来的喜悦。《音乐的奉献》是一部赋格的赋格，很像艾舍尔和哥德尔所构造的那种缠结的层次结构，是一个智慧的结晶。它以一种我无法表达的方式使我想起了人类思维这个美妙的多声部赋格。因此，我才在本书中把\heititext{哥德尔}、\heititext{艾舍尔}、\heititext{巴赫}这三块我精心收\heititext{集}的\heititext{异}彩夺目的瑰\heititext{璧}嵌为一体，并使\heititext{之}发扬光大、辉映\heititext{成}章。而这三块有\heititext{异}曲同工之妙的奇珍，也因此凝\heititext{集}成了一个珠联\heititext{璧}合的整体。

\begin{enumerate}
\item
  \phantomsection\label{Chapter20.xhtml_B_1}
\item
  \phantomsection\label{Chapter20.xhtml_B_2}
\item
  \phantomsection\label{Chapter20.xhtml_B_3}
\item
  \phantomsection\label{Chapter20.xhtml_B_4}
\item
  \phantomsection\label{Chapter20.xhtml_B_5}
\end{enumerate}

\phantomsection\label{Dialog21.xhtml}{}


\end{document}
